% Generated mechanically from frozen input p04/ko/properties.tex.
% Do not edit this generated file; edit the bound locale source instead.

% OK, start here.
%

\setcounter{chapter}{27}
\chapter{스킴의 성질}



\phantomsection
\label{properties-section-phantom}


\section{서론}
\label{properties-section-introduction}

\noindent
이 장에서는 스킴의 몇 가지 절대적 성질을 도입한다.
기본 참고문헌은 \cite{EGA}이다.




\section{구성가능 집합}
\label{properties-section-constructible}

\noindent
구성가능 집합과 국소 구성가능 집합은
Topology의 절 \ref{topology-section-constructible}에서 도입하였다.
스킴의 국소 구성가능 부분집합은 다음과 같이 특징지을 수 있다.

\begin{lemma}
\label{properties-lemma-locally-constructible}
$X$를 스킴이라 하자.
$E$를 $X$의 부분집합이라 하자. 이 부분집합이 $X$에서 국소 구성가능일
필요충분조건은 모든 아핀 열린집합에 대하여 $E \cap U$가 $U$에서
구성가능한 것이다. 여기서 $U$는 $X$의 아핀 열린집합이다.
\end{lemma}

\begin{proof}
$E$가 국소 구성가능하다고 가정하자. 그러면 열린 덮개
$X = \bigcup U_i$가 존재하여 $E \cap U_i$가 $U_i$에서 각 $i$에 대해
구성가능하다. $V \subset X$를 임의의 아핀 열린집합이라 하자.
유한 아핀 열린 덮개 $V = V_1 \cup \ldots \cup V_m$을 각 $j$에 대하여
$V_j \subset U_i$이고 어떤 $i = i(j)$가 존재하도록 택할 수 있다.
Topology의 보조정리
\ref{topology-lemma-open-immersion-constructible-inverse-image}에 의해
각 $E \cap V_j$는 $V_j$에서 구성가능하다. 포함사상
$V_j \to V$는 준콤팩트이므로(Schemes의 보조정리
\ref{schemes-lemma-quasi-compact-affine}를 보라),
Topology의 보조정리
\ref{topology-lemma-collate-constructible}에 의해
$E \cap V$가 $V$에서 구성가능하다고 결론 내린다.
역은 자명하다.
\end{proof}

\begin{lemma}
\label{properties-lemma-generic-point-in-constructible}
$X$를 스킴이라 하고 $E \subset X$를 국소 구성가능 부분집합이라 하자.
$\xi \in X$를 $X$의 한 기약 성분의 일반점이라 하자.
\begin{enumerate}
\item $\xi \in E$이면, $\xi$의 어떤 열린 근방이 $E$에 포함된다.
\item $\xi \not \in E$이면, $\xi$의 어떤 열린 근방이 $E$와 만나지 않는다.
\end{enumerate}
\end{lemma}

\begin{proof}
국소 구성가능 부분집합의 여집합도 국소 구성가능하므로 (2)만 보이면
충분하다. $X$가 아핀이고 따라서 $E$가 구성가능하다고
가정해도 된다(보조정리 \ref{properties-lemma-locally-constructible}).
이 경우 $X$는 스펙트럴 공간이다
(Algebra의 보조정리 \ref{algebra-lemma-spec-spectral}).
그러면 Topology의 보조정리
\ref{topology-lemma-constructible-stable-specialization-closed}와
$\xi \not \in E$이면 $\xi \not \in \overline{E}$라는 사실을 얻는다.
여기에는 $X$에서 $\xi$와 다른 점 가운데 $\xi$로 특수화되는 점이
없다는 사실도 사용하였다.
\end{proof}

\begin{lemma}
\label{properties-lemma-quasi-separated-quasi-compact-open-retrocompact}
$X$를 준분리 스킴이라 하자. $X$의 임의의 두 준콤팩트 열린집합의
교집합은 $X$의 준콤팩트 열린집합이다.
$X$의 모든 준콤팩트 열린집합은 $X$에서 레트로콤팩트이다.
\end{lemma}

\begin{proof}
$U$와 $V$가 준콤팩트 열린집합이면
$U \cap V = \Delta^{-1}(U \times V)$이고, 여기서
$\Delta : X \to X \times X$는 대각사상이다.
$X$가 준분리이므로 $\Delta$는 준콤팩트이다.
따라서 $U \cap V$는 준콤팩트이다. 실제로 $U \times V$가 준콤팩트이다
(세부사항은 생략한다. $U \times V$가 아핀들의 유한 합집합임을 보이려면
Schemes의 보조정리
\ref{schemes-lemma-affine-covering-fibre-product}를 사용하라).
나머지 주장들은 첫 주장과 Topology의 보조정리
\ref{topology-lemma-topology-quasi-separated-scheme}에서 따른다.
\end{proof}

\begin{lemma}
\label{properties-lemma-quasi-compact-quasi-separated-spectral}
$X$를 준콤팩트 준분리 스킴이라 하자.
그러면 $X$의 기저 위상공간은 스펙트럴 공간이다.
\end{lemma}

\begin{proof}
Topology의 정의 \ref{topology-definition-spectral-space}에 의해,
$X$가 소버이고 준콤팩트이며 준콤팩트 열린집합들로 이루어진 기저를 갖고,
임의의 두 준콤팩트 열린집합의 교집합이 준콤팩트임을 확인해야 한다.
이는 Schemes의 보조정리 \ref{schemes-lemma-scheme-sober}와
\ref{schemes-lemma-basis-affine-opens}, 그리고 위의 보조정리
\ref{properties-lemma-quasi-separated-quasi-compact-open-retrocompact}에서 따른다.
\end{proof}

\begin{lemma}
\label{properties-lemma-constructible-quasi-compact-quasi-separated}
$X$를 준콤팩트 준분리 스킴이라 하자.
$X$의 모든 국소 구성가능 부분집합은 구성가능하다.
\end{lemma}

\begin{proof}
$X$가 준콤팩트이므로 유한 아핀 열린 덮개
$X = V_1 \cup \ldots \cup V_m$을 택할 수 있다. $X$가 준분리이므로
보조정리 \ref{properties-lemma-quasi-separated-quasi-compact-open-retrocompact}에 의해
각 $V_i$는 $X$에서 레트로콤팩트이다.
따라서 Topology의 보조정리
\ref{topology-lemma-collate-constructible}에 의해,
$E \subset X$가 $X$에서 구성가능할 필요충분조건은 각
$E \cap V_j$가 $V_j$에서 구성가능한 것이다. 그러므로 보조정리
\ref{properties-lemma-locally-constructible}를 적용하면 결론을 얻는다.
\end{proof}

\begin{lemma}
\label{properties-lemma-retrocompact}
$X$를 스킴이라 하자. $E$를 $X$의 부분집합이라 하자. 이 부분집합이
$X$에서 레트로콤팩트일 필요충분조건은 $E \cap U$가 $U$에서
준콤팩트인 것이 $X$의 모든 아핀 열린집합에 대해 성립하는 것이다.
\end{lemma}

\begin{proof}
$X$의 모든 준콤팩트 열린집합은 아핀 열린집합들의 유한 합집합이라는
사실에서 바로 따른다.
\end{proof}

\begin{lemma}
\label{properties-lemma-stratification-locally-finite-constructible}
$X = \coprod_{i \in I} X_i$를 스킴 $X$의 레트로콤팩트 부분들로
이루어진 분할이라 하자. 이 분할이 국소 유한일 필요충분조건은 각 부분이
국소 구성가능한 것이다.
\end{lemma}

\begin{proof}
레트로콤팩트, 분할, 국소 유한의 정의는 Topology의 정의
\ref{topology-definition-quasi-compact},
\ref{topology-definition-paritition}, 그리고
\ref{topology-definition-locally-finite}를 보라.

\medskip\noindent
분할이 국소 유한이고 $U \subset X$가 아핀 열린집합이라고 하자. 그러면
$U = \coprod_{i \in I} U \cap X_i$는 유한 분할이다(더 정확히 말하면,
유한 개를 제외한 모든 부분이 공집합이다). 따라서 $U \cap X_i$는
준콤팩트이고, 그 여집합은 레트로콤팩트 부분들의 유한 합집합이므로
$U$에서 레트로콤팩트이다. 따라서 Topology의 보조정리
\ref{topology-lemma-locally-closed-constructible-image}에 의해
$U \cap X_i$는 구성가능하다. 그러므로 보조정리
\ref{properties-lemma-locally-constructible}에 의해 $X_i$는 국소 구성가능하다.

\medskip\noindent
각 부분이 국소 구성가능하다고 가정하자. 그러면 임의의 아핀 열린집합
$U \subset X$에 대하여 구성가능 부분집합들로 이루어진 덮개
$U = \coprod X_i \cap U$를 얻는다. 구성가능 위상이 준콤팩트이므로
(Topology의 보조정리
\ref{topology-lemma-constructible-hausdorff-quasi-compact}를 보라),
이 덮개는 유한 세분을 갖는다. 즉, 분할은 국소 유한이다.
\end{proof}




\section{정수적, 기약, 축약 스킴}
\label{properties-section-integral}

\begin{definition}
\label{properties-definition-integral}
$X$를 스킴이라 하자. $X$가 공집합이 아니고 모든 공집합이 아닌 아핀
열린집합 $\Spec(R) = U \subset X$에 대하여 환 $R$이 정역이면,
이 스킴이 {\it 정수적}이라고 한다.
\end{definition}

\begin{lemma}
\label{properties-lemma-characterize-reduced}
$X$를 스킴이라 하자.
다음은 서로 동치이다.
\begin{enumerate}
\item 스킴 $X$는 축약이다. Schemes의 정의
\ref{schemes-definition-reduced}를 보라.
\item 아핀 열린 덮개 $X = \bigcup U_i$가 존재하여
각 $\Gamma(U_i, \mathcal{O}_X)$가 축약환이다.
\item 모든 아핀 열린집합 $U \subset X$에 대하여 환
$\mathcal{O}_X(U)$가 축약환이다.
\item 모든 열린집합 $U \subset X$에 대하여 환 $\mathcal{O}_X(U)$가
축약환이다.
\end{enumerate}
\end{lemma}

\begin{proof}
Schemes의 보조정리 \ref{schemes-lemma-reduced}와
\ref{schemes-lemma-affine-reduced}를 보라.
\end{proof}

\begin{lemma}
\label{properties-lemma-characterize-irreducible}
$X$를 스킴이라 하자.
다음은 서로 동치이다.
\begin{enumerate}
\item 스킴 $X$는 기약이다.
\item 아핀 열린 덮개 $X = \bigcup_{i \in I} U_i$가 존재하여
$I$가 공집합이 아니고, $U_i$가 모든 $i \in I$에 대하여 기약이며,
$U_i \cap U_j \not = \emptyset$가 모든 $i, j \in I$에 대하여 성립한다.
\item 스킴 $X$는 공집합이 아니고 모든 공집합이 아닌 아핀 열린집합
$U \subset X$는 기약이다.
\end{enumerate}
\end{lemma}

\begin{proof}
(1)을 가정하자. Schemes의 보조정리
\ref{schemes-lemma-scheme-sober}에 의해 $X$는 유일한 일반점
$\eta$를 갖는다. 그러면 $X = \overline{\{\eta\}}$이다. 따라서
$\eta$는 모든 공집합이 아닌 아핀 열린집합 $U \subset X$의 원소이다.
이는 $\eta \in U$가 조밀하고 따라서 $U$가 기약임을 뜻한다.
또한 임의의 두 공집합이 아닌 아핀 열린집합이 만남을 뜻한다.
그러므로 (1)은 (2)와 (3)을 모두 함의한다.

\medskip\noindent
(2)를 가정하자. $X = Z_1 \cup Z_2$가 두 닫힌 부분집합의 합집합이라고 하자.
모든 $i$에 대하여 $U_i \subset Z_1$ 또는 $U_i \subset Z_2$이다.
어떤 $i \in I$를 택하고 $U_i \subset Z_1$이라고 가정하자
(필요하면 $Z_1$, $Z_2$의 번호를 바꾼다). 임의의 $j \in I$에 대하여
열린 부분집합 $U_i \cap U_j$는 $U_j$에서 조밀하고 닫힌 부분집합
$Z_1 \cap U_j$에 포함된다. 따라서 $U_j \subset Z_1$이다.
그러므로 원하는 대로 $X = Z_1$이다.

\medskip\noindent
(3)을 가정하자. 아핀 열린 덮개 $X = \bigcup_{i \in I} U_i$를 택하자.
각 $U_i$가 공집합이 아니라고 가정해도 된다.
$X$가 공집합이 아니므로 $I$도 공집합이 아니다.
가정에 의해 각 $U_i$는 기약이다.
$U_i \cap U_j = \emptyset$인 어떤 쌍 $i, j \in I$가 있다고 하자.
그러면 열린집합 $U_i \amalg  U_j = U_i \cup U_j$는 아핀이다.
Schemes의 보조정리 \ref{schemes-lemma-disjoint-union-affines}를 보라.
따라서 가정에 의해 이 열린집합은 기약인데, 이는 모순이다.
그러므로 (3)은 (2)를 함의한다. 보조정리가 증명되었다.
\end{proof}

\begin{lemma}
\label{properties-lemma-characterize-integral}
스킴 $X$가 정수적일 필요충분조건은 그것이 축약이고 기약인 것이다.
\end{lemma}

\begin{proof}
$X$가 기약이면 모든 아핀 열린집합 $\Spec(R) = U \subset X$는
기약이다. $X$가 축약이면 위의 보조정리
\ref{properties-lemma-characterize-reduced}에 의해 $R$은 축약환이다. 따라서
$R$은 축약환이고 $(0)$은 소 아이디얼이다. 즉, $R$은 정역이다.

\medskip\noindent
$X$가 정수적이면 모든 공집합이 아닌 아핀 열린집합
$\Spec(R) = U \subset X$에 대하여 환 $R$은 축약이고, 따라서 보조정리
\ref{properties-lemma-characterize-reduced}에 의해 $X$는 축약이다.
더욱이 모든 공집합이 아닌 아핀 열린집합은 기약이다.
따라서 보조정리 \ref{properties-lemma-characterize-irreducible}에 의해 $X$는
기약이다.
\end{proof}

\noindent
Examples의 절
\ref{examples-section-connected-locally-integral-not-integral}에서는
모든 국소환이 정역이지만 정수적이지 않은 연결 아핀 스킴을 구성한다.










\section{환의 성질로 정의되는 스킴의 유형}
\label{properties-section-properties-rings}

\noindent
이 절에서는 환의 어떤 성질들로 스킴의 국소 성질을 정의할 수 있는지
연구한다.

\begin{definition}
\label{properties-definition-property-local}
$P$를 환의 성질이라 하자.
다음이 성립할 때 $P$가 {\it 국소적}이라고 한다.
\begin{enumerate}
\item 임의의 환 $R$과 임의의 $f \in R$에 대하여
$P(R) \Rightarrow P(R_f)$이다.
\item 임의의 환 $R$과 $f_i \in R$에 대하여
$(f_1, \ldots, f_n) = R$이면
$\forall i, P(R_{f_i}) \Rightarrow P(R)$이다.
\end{enumerate}
\end{definition}

\begin{definition}
\label{properties-definition-locally-P}
$P$를 환의 성질이라 하자. $X$를 스킴이라 하자.
$X$가 {\it 국소적으로 $P$}라고 함은, 임의의 $x \in X$에 대하여
아핀 열린 근방 $U$가 $x$를 포함하도록 $X$ 안에 존재하고
$\mathcal{O}_X(U)$가 성질 $P$를 가지는 경우이다.
\end{definition}

\noindent
이는 해당 성질이 국소적일 때에만 좋은 개념이다.
$P$가 국소 성질이더라도, 다른 곳에서 그 정의를 명시적으로
서술하지 않는 한 이 정의를 자동으로 사용하여 스킴이
``국소적으로 $P$''라고 말하지는 않을 것이다.

\begin{lemma}
\label{properties-lemma-locally-P}
$X$를 스킴이라 하자. $P$를 환의 국소 성질이라 하자.
다음은 서로 동치이다.
\begin{enumerate}
\item 스킴 $X$는 국소적으로 $P$이다.
\item 모든 아핀 열린집합 $U \subset X$에 대하여 성질
$P(\mathcal{O}_X(U))$가 성립한다.
\item 아핀 열린 덮개 $X = \bigcup U_i$가 존재하여
각 $\mathcal{O}_X(U_i)$가 $P$를 만족한다.
\item 열린 덮개 $X = \bigcup X_j$가 존재하여
각 열린 부분스킴 $X_j$가 국소적으로 $P$이다.
\end{enumerate}
더욱이 $X$가 국소적으로 $P$이면 모든 열린 부분스킴도
국소적으로 $P$이다.
\end{lemma}

\begin{proof}
물론 (1) $\Leftrightarrow$ (3)이고 (2) $\Rightarrow$ (1)이다.
(3) $\Rightarrow$ (2)이면 보조정리의 마지막 주장도 성립하며,
(4)도 (1)과 동치라는 것이 쉽게 따른다.
따라서 (3) $\Rightarrow$ (2)를 보이면 된다.

\medskip\noindent
$X = \bigcup U_i$를 아핀 열린 덮개라 하고
$U_i = \Spec(R_i)$라고 하자. $P(R_i)$라고 가정하자.
$\Spec(R) = U \subset X$를 임의의 아핀 열린집합이라 하자.
Schemes의 보조정리 \ref{schemes-lemma-good-subcover}에 의해
$U = \Spec(R)$에는 표준 열린집합 $D(f_j)$들로 이루어진 표준 덮개가
존재하여 각 환 $R_{f_j}$는 환들 $R_i$ 가운데 하나의 주 국소화이다.
정의 \ref{properties-definition-property-local} (1)에 의해 $P(R_{f_j})$를 얻는다.
그러므로 정의 \ref{properties-definition-property-local} (2)에 의해 $P(R)$이다.
\end{proof}

\noindent
응용의 한 예는 다음과 같다.

\begin{lemma}
\label{properties-lemma-reduced-is-locally-reduced}
$X$를 스킴이라 하자. 그러면 $X$가 축약일 필요충분조건은
정의 \ref{properties-definition-locally-P}의 의미에서 $X$가
``국소적으로 축약''인 것이다.
\end{lemma}

\begin{proof}
보조정리 \ref{properties-lemma-characterize-reduced}에서 명백하다.
\end{proof}

\begin{lemma}
\label{properties-lemma-properties-local}
환 $R$의 다음 성질들은 국소적이다.
\begin{enumerate}
\item (코언--매콜리.)
환 $R$은 뇌터이고 CM이다. Algebra의 정의
\ref{algebra-definition-ring-CM}을 보라.
\item (정칙.)
환 $R$은 뇌터이고 정칙이다. Algebra의 정의
\ref{algebra-definition-regular}를 보라.
\item (절대 뇌터.)
환 $R$은 $Z$ 위에서 유한형이다.
\item 필요에 따라 여기에 더 추가한다.\footnote{다만 다른 곳에서 이미
별도로 다룬 성질들은 여기에 열거하지 않는다.}
\end{enumerate}
\end{lemma}

\begin{proof}
생략한다.
\end{proof}















\section{뇌터 스킴}
\label{properties-section-noetherian}

\noindent
환 $R$이 아이디얼들의 오름 사슬 조건을 만족하면 {\it 뇌터}라고
부른다는 것을 상기하자. 동치로, $R$의 모든 아이디얼이 유한 생성이다.

\begin{definition}
\label{properties-definition-noetherian}
$X$를 스킴이라 하자.
\begin{enumerate}
\item $X$가 {\it 국소 뇌터}라고 함은 모든 $x \in X$가 아핀 열린 근방
$\Spec(R) = U \subset X$를 가지며 환 $R$이 뇌터인 경우이다.
\item $X$가 국소 뇌터이고 준콤팩트이면 $X$가 {\it 뇌터}라고 한다.
\end{enumerate}
\end{definition}

\noindent
국소 뇌터 스킴을 특징짓는 표준 결과는 다음과 같다.

\begin{lemma}
\label{properties-lemma-locally-Noetherian}
$X$를 스킴이라 하자. 다음은 서로 동치이다.
\begin{enumerate}
\item 스킴 $X$는 국소 뇌터이다.
\item 모든 아핀 열린집합 $U \subset X$에 대하여 환
$\mathcal{O}_X(U)$는 뇌터이다.
\item 아핀 열린 덮개 $X = \bigcup U_i$가 존재하여
각 $\mathcal{O}_X(U_i)$가 뇌터이다.
\item 열린 덮개 $X = \bigcup X_j$가 존재하여
각 열린 부분스킴 $X_j$가 국소 뇌터이다.
\end{enumerate}
더욱이 $X$가 국소 뇌터이면 모든 열린 부분스킴도 국소 뇌터이다.
\end{lemma}

\begin{proof}
뇌터라는 것이 환의 국소 성질임을 보이면 충분하다.
보조정리 \ref{properties-lemma-locally-P}를 보라.
뇌터환의 모든 국소화는 뇌터이다. Algebra의 보조정리
\ref{algebra-lemma-Noetherian-permanence}를 보라.
Algebra의 보조정리 \ref{algebra-lemma-cover}에 의해
정의 \ref{properties-definition-property-local}의 두 번째 성질을 얻는다.
\end{proof}

\begin{lemma}
\label{properties-lemma-immersion-into-noetherian}
임의의 몰입 $Z \to X$는 $X$가 국소 뇌터이면 준콤팩트이다.
\end{lemma}

\begin{proof}
닫힌 몰입은 명백히 준콤팩트이다.
준콤팩트 사상들의 합성은 준콤팩트이다. Topology의 보조정리
\ref{topology-lemma-composition-quasi-compact}를 보라.
따라서 국소 뇌터 스킴으로 들어가는 열린 몰입이 준콤팩트임을 보이면
충분하다.
Schemes의 보조정리 \ref{schemes-lemma-quasi-compact-affine}를 사용하면
$X$가 아핀인 경우로 환원된다.
뇌터환의 스펙트럼의 모든 열린 부분집합은 준콤팩트이다
(예를 들어 Algebra의 보조정리
\ref{algebra-lemma-Noetherian-topology}와 Topology의 보조정리
\ref{topology-lemma-Noetherian} 및
\ref{topology-lemma-Noetherian-quasi-compact}를 결합하라).
\end{proof}

\begin{lemma}
\label{properties-lemma-locally-Noetherian-quasi-separated}
국소 뇌터 스킴은 준분리이다.
\end{lemma}

\begin{proof}
Schemes의 보조정리
\ref{schemes-lemma-characterize-quasi-separated}에 의해,
두 아핀 열린집합의 교집합 $U \cap V$가 준콤팩트임을 보이면 된다.
여기서 두 열린집합은 $X$의 열린집합이다.
이는 예를 들어 열린 몰입 $U \cap V \to U$에 위의 보조정리
\ref{properties-lemma-immersion-into-noetherian}를 적용하면 따른다.
(그러나 실제로는 뇌터환의 스펙트럼의 모든 열린집합이
준콤팩트이기 때문이다.)
\end{proof}

\begin{lemma}
\label{properties-lemma-Noetherian-topology}
(국소) 뇌터 스킴의 기저 위상공간은 (국소) 뇌터이다.
Topology의 정의 \ref{topology-definition-noetherian}를 보라.
\end{lemma}

\begin{proof}
뇌터 스킴이 뇌터환들의 스펙트럼들의 유한 합집합이라는 사실과
Algebra의 보조정리 \ref{algebra-lemma-Noetherian-topology} 및
Topology의 보조정리
\ref{topology-lemma-finite-union-Noetherian}에서 따른다.
\end{proof}

\begin{lemma}
\label{properties-lemma-locally-closed-in-Noetherian}
(국소) 뇌터 스킴의 모든 국소닫힌 부분스킴은 (국소) 뇌터이다.
\end{lemma}

\begin{proof}
생략한다. 힌트: 뇌터환의 모든 몫환과 모든 국소화는 뇌터이다.
뇌터인 경우에는 뇌터 공간의 모든 부분집합이 유도위상에 관하여
뇌터 공간이라는 사실을 다시 사용하라.
\end{proof}

\begin{lemma}
\label{properties-lemma-Noetherian-irreducible-components}
뇌터 스킴은 유한 개의 기약 성분을 갖는다.
\end{lemma}

\begin{proof}
뇌터 스킴의 기저 위상공간은 뇌터이고
(보조정리 \ref{properties-lemma-Noetherian-topology}),
뇌터 위상공간은 유한 개의 기약 성분만 가지므로 결론이 따른다
(Topology의 보조정리 \ref{topology-lemma-Noetherian}).
\end{proof}

\begin{lemma}
\label{properties-lemma-morphism-Noetherian-schemes-quasi-compact}
임의의 스킴 사상 $f : X \to Y$는 $X$가 뇌터이면 준콤팩트이다.
\end{lemma}

\begin{proof}
보조정리 \ref{properties-lemma-Noetherian-topology}와, 뇌터 위상공간의 모든
부분집합이 준콤팩트라는 사실을 사용하라(Topology의 보조정리
\ref{topology-lemma-Noetherian}와
\ref{topology-lemma-Noetherian-quasi-compact}를 보라).
\end{proof}

\noindent
다음은 재미있는 보조정리이다.
모든 국소 뇌터 스킴에는 닫힌점이 충분히 많다는 것
(각 닫힌 부분집합마다 적어도 하나가 있다는 것)을 말한다.

\begin{lemma}
\label{properties-lemma-locally-Noetherian-closed-point}
공집합이 아닌 모든 국소 뇌터 스킴은 닫힌점을 갖는다.
국소 뇌터 스킴의 공집합이 아닌 모든 닫힌 부분집합은 닫힌점을 갖는다.
동치로, 국소 뇌터 스킴의 모든 점은 한 닫힌점으로 특수화된다.
\end{lemma}

\begin{proof}
두 번째 주장은 첫 번째 주장으로부터 따른다
(Schemes의 보조정리 \ref{schemes-lemma-reduced-closed-subscheme}와
보조정리 \ref{properties-lemma-locally-closed-in-Noetherian}를 사용하라).
공집합이 아닌 임의의 아핀 열린집합 $U \subset X$를 생각하자.
$x \in U$를 닫힌점이라 하자. $x$가 $X$의 닫힌점이면 끝난다.
그렇지 않으면 $X_0 \subset X$를 $\overline{\{x\}}$ 위의
유도된 축약 닫힌 부분스킴 구조라 하자.
그러면 $U_0 = U \cap X_0$는 Schemes의 보조정리
\ref{schemes-lemma-closed-subspace-scheme}에 의해 $X_0$의 아핀
열린집합이고, $U_0 = \{x\}$이다. $y \in X_0$와 $y \not = x$를
$x$의 한 특수화라 하자.
국소환 $R = \mathcal{O}_{X_0, y}$를 생각하자.
보조정리 \ref{properties-lemma-locally-closed-in-Noetherian}에 의해 $X_0$가
뇌터이므로 이는 뇌터 국소환이다. $V \subset \Spec(R)$를 표준 사상
$U_0$의 $\Spec(R)$ 안의 역상이라 하자. 여기서 사용한 표준 사상은
$\Spec(R) \to X_0$이다
(Schemes의 절 \ref{schemes-section-points}를 보라).
구성에 의해 $V$는 $x$에 대응하는 유일한 점을 갖는 한원소집합이다
(Schemes의 보조정리 \ref{schemes-lemma-specialize-points}를 사용하라).
Algebra의 보조정리
\ref{algebra-lemma-Noetherian-local-domain-dim-2-infinite-opens}에 의해
$\dim(R) = 1$임을 알 수 있다.
달리 말해 $y$는 $x$의 바로 다음 특수화이다
(Topology의 정의 \ref{topology-definition-dimension-function}를 보라).
다시 말해 $y \not = x$이고 $x \leadsto y$인 모든 점은
$x$의 바로 다음 특수화이다. 이 점들은 각각 명백히 닫힌점이므로
원하는 결론을 얻는다.
\end{proof}


\begin{lemma}
\label{properties-lemma-locally-Noetherian-specialization-dvr}
$X$를 국소 뇌터 스킴이라 하자.
$x' \leadsto x$를 $X$의 점들의 특수화라 하자. 그러면 다음이 성립한다.
\begin{enumerate}
\item 이산 값매김환 $R$과 사상 $f : \Spec(R) \to X$가 존재하여,
$\eta$는 $\Spec(R)$의 일반점이고 $x'$으로 가며 특수점은 $x$로 간다.
\item $x \not = x'$이라고 하자. 유한 생성 체 확대 $K/\kappa(x')$가
주어지면, 확대 $\kappa(\eta)/\kappa(x')$가 $f$에 의해 유도되고
주어진 확대와 동형이 되도록 할 수 있다.
\end{enumerate}
\end{lemma}

\begin{proof}
먼저 $x' \leadsto x$가 $X$에서의 특수화이고 $x \not = x'$이라고
가정하며, $K/\kappa(x')$를 유한 생성 체 확대라 하자.
Schemes의 보조정리 \ref{schemes-lemma-specialize-points}와
Schemes의 보조정리 \ref{schemes-lemma-characterize-points} 뒤의
논의에 의해 환 사상
$\mathcal{O}_{X, x} \to \kappa(x') \to K$를 얻는다.
$x \not = x'$이므로 $\mathcal{O}_{X, x}$의 $K$ 안의 상은
체가 아니다(세부사항은 생략한다).
$R \subset K$를 분수체가 $K$이고 $\mathcal{O}_{X, x} \to K$의
상을 지배하는 임의의 이산 값매김환이라 하자. Algebra의 보조정리
\ref{algebra-lemma-exists-dvr}를 보라.
환 사상 $\mathcal{O}_{X, x} \to R$은 사상
$f : \Spec(R) \to X$를 유도한다. Schemes의 보조정리
\ref{schemes-lemma-morphism-from-spec-local-ring}을 보라.
구성에 의해 이 사상은 원하는 모든 성질을 갖는다.
$x = x'$이면 $R = \kappa(x)[t]_{(t)}$로 둘 수 있다.
\end{proof}

\begin{lemma}
\label{properties-lemma-thin-infinite-sequence}
$S$를 뇌터 스킴이라 하자. $T \subset S$를 무한 부분집합이라 하자.
그러면 무한 부분집합 $T' \subset T$가 존재하여 $T'$의 점들 사이에는
자명하지 않은 특수화가 없다.
\end{lemma}

\begin{proof}
$T_0 \subset T$를, $t \in T$ 중 $T$의 다른 점으로 특수화되지 않는
모든 점의 집합이라 하자. $T_0$가 무한이면 $T' = T_0$로 두면 된다.
따라서 $T_0$가 유한이라고 가정해도 되고 그렇게 하자.
귀납적으로 $i > 0$에 대하여 집합 $T_i \subset T$를 다음을 만족하는
$t \in T$들로 정의하자.
\begin{enumerate}
\item $t \not \in T_{i - 1} \cup T_{i - 2} \cup \ldots \cup T_0$이다.
\item 자명하지 않은 특수화 $t \leadsto t'$가 존재하여
$t' \in T_{i - 1}$이다.
\item 임의의 자명하지 않은 특수화 $t \leadsto t'$에 대하여
$t' \in T$이면 $t' \in T_{i - 1} \cup T_{i - 2} \cup \ldots \cup T_0$이다.
\end{enumerate}
다시, $T_i$가 무한이면 $T' = T_i$로 두면 된다.
$d$를 국소환 $\mathcal{O}_{S, t}$의 차원들의 최댓값이라 하자.
여기서 $t \in T_0$이다. $d$가 정수인 것은 $T_0$가 유한이고 Algebra의 명제
\ref{algebra-proposition-dimension}에 의해 국소환들의 차원이
유한하기 때문이다.
그러면 $T_i = \emptyset$이다($i > d$).
실제로 $t \in T_i$이면 자명하지 않은 특수화들의 열
$t = t_i \leadsto t_{i - 1} \leadsto \ldots \leadsto t_0$를 찾을 수 있고,
$t_0 \in T_0$이다. 점들
$t = t_i, t_{i - 1}, \ldots, t_0$는
$\Spec(\mathcal{O}_{S, t_0})$에 속하므로
(Schemes의 보조정리 \ref{schemes-lemma-specialize-points}),
$i \leq d$임을 알 수 있다.
따라서 $\bigcup T_i = T_d \cup \ldots \cup T_0$는 $T$의 유한
부분집합이다.

\medskip\noindent
$t \in T$가 $\bigcup T_i$에 속하지 않는다고 하자. 그러면 특수화
$t \leadsto t'$가 존재하여 $t' \in T$이고 $t' \not \in \bigcup T_i$이어야
한다. (실제로 $t$의 모든 특수화가 유한 집합
$T_d \cup \ldots \cup T_0$에 속하면, 최대 $i$가 존재하여 특수화
$t \leadsto t'$와 $t' \in T_i$가 있고, 구성에 의해
$t \in T_{i + 1}$이다.) 따라서
$$
t \leadsto t' \leadsto t'' \leadsto \ldots
$$
와 같이 $T \setminus \bigcup T_i$의 점들 사이에 자명하지 않은
특수화들의 무한 열을 얻는다. 이는 $S$의 기저 위상공간이 보조정리
\ref{properties-lemma-locally-Noetherian-quasi-separated}에 의해 뇌터이므로
불가능하다.
\end{proof}

\begin{lemma}
\label{properties-lemma-maximal-points}
$S$를 뇌터 스킴이라 하자. $T \subset S$를 부분집합이라 하자.
$T_0 \subset T$를, $t \in T$ 중 자명하지 않은 특수화
$t' \leadsto t$와 $t' \in T'$가 존재하지 않는 모든 점의 집합이라 하자.
그러면 (a) $T_0$의 점들 사이에는 특수화가 없고, (b) $T$의 모든 점은
$T_0$의 한 점의 특수화이며, (c) $T$와 $T_0$의 폐포는 같다.
\end{lemma}

\begin{proof}
$\dim(\mathcal{O}_{S, s}) < \infty$임이 모든 $s \in S$에 대하여
성립함을 상기하자. Algebra의 명제
\ref{algebra-proposition-dimension}을 보라. $t \in T$라 하자.
$t' \leadsto t$이면 차원론에 의해
$\dim(\mathcal{O}_{S, t'}) \leq \dim(\mathcal{O}_{S, t})$이고,
등호가 성립할 필요충분조건은 $t' = t$인 것이다. 따라서
$t' \leadsto t$를 $\dim(\mathcal{O}_{T, t'})$가 최소가 되도록
택하면 $t' \in T_0$이다. 달리 말해, 모든 $t \in T$는
$T_0$의 한 원소의 특수화이다.
\end{proof}

\begin{lemma}
\label{properties-lemma-countable-dense-subset}
$S$를 뇌터 스킴이라 하자. $T \subset S$를 무한 조밀 부분집합이라 하자.
그러면 가산 부분집합 $E \subset T$가 존재하며 $S$에서 조밀하다.
\end{lemma}

\begin{proof}
$T'$을, 모든 점 $s \in S$ 가운데 $\overline{\{s\}} \cap T$가
폐포가 $\overline{\{s\}}$인 가산 부분집합을 포함하는 점들의 집합이라
하자. 유한 집합은 가산이므로 $T \subset T'$이다.
$s \in T'$에 대하여 그러한 가산 부분집합
$E_s \subset \overline{\{s\}} \cap T$를 택하자.
$E' = \{s_1, s_2, s_3, \ldots\} \subset T'$을 가산 부분집합이라 하자.
그러면 $E'$의 $S$에서의 폐포는 가산 부분집합
$\bigcup_n E_{s_n}$의 폐포이며, 이 부분집합은 $T$에 포함된다.
따라서 $Z$가 $E'$의 폐포의 한
기약 성분이면 $Z$의 일반점은 $T'$에 속한다.

\medskip\noindent
$T'_0 \subset T'$을, 보조정리 \ref{properties-lemma-maximal-points}에서처럼
$t \in T'$ 중 자명하지 않은 특수화 $t' \leadsto t$와
$t' \in T'$가 존재하지 않는 모든 점의 부분집합이라 하자.
그 결과들을 더 언급하지 않고
사용할 것이다. $T'_0$가 무한이면 가산 부분집합
$E' \subset T'_0$을 택한다. 첫 문단의 논증에 의해 $E'$의 폐포의
기약 성분들의 일반점들은 $T'$에 속한다. 그러나 이 점들 가운데 하나는
$E' \subset T'_0$의 서로 다른 원소 무한 개로 특수화되므로 모순이다.
따라서 $T'_0$는 유한이다. 이를
$T'_0 = \{s_1, \ldots, s_m\}$이라 하자. 그러면 $S$는 $T$의 폐포이고
$\{s_1, \ldots, s_m\}$의 폐포에 포함되고, 이는 다시 가산 부분집합
$E_{s_1} \cup \ldots \cup E_{s_m} \subset T$의 폐포에 포함된다.
이로써 원하는 결론을 얻는다.
\end{proof}








\section{Jacobson 스킴}
\label{properties-section-jacobson}

\noindent
모든 닫힌 부분집합에서 닫힌점들이 조밀한 공간을 {\it Jacobson}이라고
한다는 것을 상기하자. Topology의 절
\ref{topology-section-space-jacobson}을 보라.

\begin{definition}
\label{properties-definition-jacobson}
스킴 $S$의 기저 위상공간이 Jacobson이면 이를 {\it Jacobson}이라고 한다.
\end{definition}

\noindent
환 $R$의 모든 근기 아이디얼이 극대 아이디얼들의 교집합이면
$R$이 Jacobson이라고 한다는 것을 상기하자. Algebra의 정의
\ref{algebra-definition-ring-jacobson}을 보라.

\begin{lemma}
\label{properties-lemma-affine-jacobson}
아핀 스킴 $\Spec(R)$가 Jacobson일 필요충분조건은 환 $R$이
Jacobson인 것이다.
\end{lemma}

\begin{proof}
이는 Algebra의 보조정리 \ref{algebra-lemma-jacobson}이다.
\end{proof}

\noindent
Jacobson 스킴을 특징짓는 표준 결과는 다음과 같다.
직관적으로 Jacobson $\Leftrightarrow$ 국소적으로 Jacobson임을 주장한다.

\begin{lemma}
\label{properties-lemma-locally-jacobson}
$X$를 스킴이라 하자. 다음은 서로 동치이다.
\begin{enumerate}
\item 스킴 $X$는 Jacobson이다.
\item 스킴 $X$는 정의 \ref{properties-definition-locally-P}의 의미에서
``국소적으로 Jacobson''이다.
\item 모든 아핀 열린집합 $U \subset X$에 대하여 환
$\mathcal{O}_X(U)$는 Jacobson이다.
\item 아핀 열린 덮개 $X = \bigcup U_i$가 존재하여
각 $\mathcal{O}_X(U_i)$가 Jacobson이다.
\item 열린 덮개 $X = \bigcup X_j$가 존재하여
각 열린 부분스킴 $X_j$가 Jacobson이다.
\end{enumerate}
더욱이 $X$가 Jacobson이면 모든 열린 부분스킴도 Jacobson이다.
\end{lemma}

\begin{proof}
보조정리의 마지막 주장은 Topology의 보조정리
\ref{topology-lemma-jacobson-inherited}에 의해 성립한다.
(5)와 (1)의 동치는 Topology의 보조정리
\ref{topology-lemma-jacobson-local}이다.
따라서 보조정리 \ref{properties-lemma-affine-jacobson}를 사용하면
(1) $\Leftrightarrow$ (2)임을 알 수 있다.
보조정리의 증명을 끝내려면 ``Jacobson''이 환의 국소 성질임을 보이면
충분하다. 보조정리 \ref{properties-lemma-locally-P}를 보라.
Jacobson 환을 한 원소에서 국소화한 환은 Jacobson이다. Algebra의 보조정리
\ref{algebra-lemma-Jacobson-invert-element}를 보라.
$R$을 환이라 하고 $f_1, \ldots, f_n \in R$이 단위 아이디얼을 생성하며
각 $R_{f_i}$가 Jacobson이라고 가정하자. 그러면
$\Spec(R) = \bigcup D(f_i)$는 모두 Jacobson인 열린 부분집합들의
합집합이므로, Topology의 보조정리
\ref{topology-lemma-jacobson-local}에 의해 $\Spec(R)$도 Jacobson이다.
이는 정의 \ref{properties-definition-property-local}의 두 번째 성질을 증명한다.
\end{proof}

\noindent
대수기하에서 흔히 쓰이는 많은 스킴은 Jacobson이다. Morphisms의 보조정리
\ref{morphisms-lemma-ubiquity-Jacobson-schemes}를 보라.
여기서는 다음 흥미로운 경우를 언급한다.

\begin{lemma}
\label{properties-lemma-complement-closed-point-Jacobson}
뇌터 Jacobson 스킴의 예들은 다음과 같다.
\begin{enumerate}
\item $(R, \mathfrak m)$이 뇌터 국소환이면 구멍 뚫린 스펙트럼
$\Spec(R) \setminus \{\mathfrak m\}$은 Jacobson 스킴이다.
\item $R$이 Jacobson 근기 $\text{rad}(R)$을 갖는 뇌터환이면
$\Spec(R) \setminus V(\text{rad}(R))$은 Jacobson 스킴이다.
\item $(R, I)$가 Zariski 쌍이고(More on Algebra의 정의
\ref{more-algebra-definition-zariski-pair}), $R$이 뇌터이면
$\Spec(R) \setminus V(I)$는 Jacobson 스킴이다.
\end{enumerate}
\end{lemma}

\begin{proof}
(3)의 증명. $\Spec(R) - V(I)$는 아핀 열린집합
$\Spec(R_f)$들로 이루어진 덮개를 가짐을 관찰하자. 여기서 $f \in I$이다.
환들 $R_f$는
More on Algebra의 보조정리
\ref{more-algebra-lemma-noetherian-zariski-jacobson-complement}에 의해
Jacobson이다. 따라서 보조정리 \ref{properties-lemma-locally-jacobson}에 의해
$\Spec(R) \setminus V(I)$는 Jacobson이다.
(1)과 (2)는 (3)의 특수한 경우이다.

\medskip\noindent
(1)의 직접 증명.
$\Spec(R)$가 뇌터 스킴이므로
$S$는 뇌터 스킴이다(보조정리 \ref{properties-lemma-locally-closed-in-Noetherian}).
따라서 $S$는 소버 뇌터 위상공간이다
(Schemes의 보조정리 \ref{schemes-lemma-scheme-sober}를 사용하라).
$S$가 Jacobson이 아니라고 가정하여 모순을 얻자.
Topology의 보조정리
\ref{topology-lemma-non-jacobson-Noetherian-characterize}에 의해
어떤 닫히지 않은 점 $\xi \in S$가 존재하여 $\{\xi\}$는
국소닫혀 있다. 이는 다음을 만족하는 소 아이디얼
$\mathfrak p \subset R$에 대응한다. (1) 두 포함이 모두 진포함인
소 아이디얼 $\mathfrak q$와
$\mathfrak p \subset \mathfrak q \subset \mathfrak m$이 존재하고,
(2) $\{\mathfrak p\}$는 $\Spec(R/\mathfrak p)$에서 열려 있다.
이는 Algebra의 보조정리
\ref{algebra-lemma-Noetherian-local-domain-dim-2-infinite-opens}에 의해
불가능하다.
\end{proof}

\section{정규 스킴}
\label{properties-section-normal}

\noindent
환 $R$의 모든 국소환이 정규 정역이면 그 환을 정규라고 함을 상기하자.
Algebra의 정의 \ref{algebra-definition-ring-normal}를 보라.
정규 정역은 자신의 분수체 안에서 정수적으로 닫힌 정역이다.
Algebra의 정의 \ref{algebra-definition-domain-normal}를 보라.
따라서 정규 스킴을 다음과 같이 정의하는 것이 자연스럽다.

\begin{definition}
\label{properties-definition-normal}
스킴 $X$가 {\it 정규}라는 것은 모든 $x \in X$에 대하여 국소환
$\mathcal{O}_{X, x}$가 정규 정역이라는 것과 같다.
\end{definition}

\noindent
이는 EGA에서 사용되는 정의인 듯하다. \cite[0, 4.1.4]{EGA}를 보라.
$X = \Spec(A)$이고 $A$가 축약이라고 하자. 이때 $X$가 정규라는 것과
$A$가 자신의 전분수환 안에서 정수적으로 닫혀 있다는 것은 동치가 아니다.
그러나 $A$가 뇌터이면 둘은 동치이다
(Algebra의 보조정리 \ref{algebra-lemma-characterize-reduced-ring-normal}를 보라).

\begin{lemma}
\label{properties-lemma-locally-normal}
$X$를 스킴이라 하자. 다음은 서로 동치이다.
\begin{enumerate}
\item 스킴 $X$는 정규이다.
\item 모든 아핀 열린집합 $U \subset X$에 대하여 환 $\mathcal{O}_X(U)$는
정규이다.
\item 아핀 열린 덮개 $X = \bigcup U_i$가 존재하여 각
$\mathcal{O}_X(U_i)$가 정규이다.
\item 열린 덮개 $X = \bigcup X_j$가 존재하여 각 열린 부분스킴
$X_j$가 정규이다.
\end{enumerate}
더욱이 $X$가 정규이면 모든 열린 부분스킴도 정규이다.
\end{lemma}

\begin{proof}
정의에서 바로 따른다.
\end{proof}

\begin{lemma}
\label{properties-lemma-normal-reduced}
정규 스킴은 축약이다.
\end{lemma}

\begin{proof}
정의에서 즉시 따른다.
\end{proof}

\begin{lemma}
\label{properties-lemma-integral-normal}
$X$를 정수적 스킴이라 하자. 그러면 $X$가 정규라는 것은 모든 공집합이
아닌 아핀 열린집합 $U \subset X$에 대하여 환 $\mathcal{O}_X(U)$가
정규 정역이라는 것과 같다.
\end{lemma}

\begin{proof}
Algebra의 보조정리 \ref{algebra-lemma-normality-is-local}에서 따른다.
\end{proof}

\begin{lemma}
\label{properties-lemma-normal-locally-finite-nr-irreducibles}
모든 준콤팩트 열린집합이 유한 개의 기약 성분을 갖는 스킴 $X$를 생각하자.
다음은 서로 동치이다.
\begin{enumerate}
\item $X$는 정규이고,
\item $X$는 정규 정수적 스킴들의 서로소 합집합이다.
\end{enumerate}
\end{lemma}

\begin{proof}
(2)가 (1)을 함의함은 정의에서 즉시 따른다.
모든 준콤팩트 열린집합이 유한 개의 기약 성분을 갖는 정규 스킴 $X$를
생각하자. $X$가 아핀이면
Algebra의 보조정리 \ref{algebra-lemma-characterize-reduced-ring-normal}에
의해 $X$는 (2)를 만족한다. 일반적인 $X$에 대하여
$X = \bigcup X_i$를 아핀 열린 덮개라 하자. 각 $X_i$ 역시 유한 개의
기약 성분만을 가지며, 각 $X_i$에 대하여 보조정리가 성립한다.
$T \subset X$를 기약 성분이라 하자. 아핀인 경우에 의해 각 교집합
$T \cap X_i$는 $X_i$에서 열려 있고 정규 정수적 스킴이다.
따라서 $T \subset X$는 열린 정규 정수적 스킴이다.
이는 $X$가 자신의 기약 성분들의 서로소 합집합이고, 그 성분들이 정규
정수적 스킴임을 보인다.
\end{proof}

\begin{lemma}
\label{properties-lemma-normal-Noetherian}
$X$를 뇌터 스킴이라 하자. 다음은 서로 동치이다.
\begin{enumerate}
\item $X$는 정규이고,
\item $X$는 유한 개의 정규 정수적 스킴들의 서로소 합집합이다.
\end{enumerate}
\end{lemma}

\begin{proof}
뇌터 스킴의 밑 위상공간은 뇌터이므로, 이는 보조정리
\ref{properties-lemma-normal-locally-finite-nr-irreducibles}의 특수한 경우이다
(보조정리 \ref{properties-lemma-Noetherian-topology} 및
Topology의 보조정리 \ref{topology-lemma-Noetherian}를 보라).
\end{proof}

\begin{lemma}
\label{properties-lemma-normal-locally-Noetherian}
$X$를 국소 뇌터 스킴이라 하자. 다음은 서로 동치이다.
\begin{enumerate}
\item $X$는 정규이고,
\item $X$는 정수적 정규 스킴들의 서로소 합집합이다.
\end{enumerate}
\end{lemma}

\begin{proof}
생략. 힌트: 보조정리 \ref{properties-lemma-normal-Noetherian}에 의해 이는 순전히
위상적인 사실이다.
\end{proof}

\begin{remark}
\label{properties-remark-normal-connected-irreducible}
$X$를 정규 스킴이라 하자. $X$가 국소 뇌터이면 보조정리
\ref{properties-lemma-normal-locally-Noetherian}에 의해 $X$가 정수적이라는 것은
$X$가 연결이라는 것과 같다. 그러나 연결 아핀 스킴 $X$ 중에는
$\mathcal{O}_{X, x}$가 모든 $x \in X$에 대하여 정역이지만
$X$는 기약이 아닌 것이 존재한다. Examples의 절
\ref{examples-section-connected-locally-integral-not-integral}을 보라.
이 예는 심지어 정규 스킴이기도 하므로(증명은 생략한다) 주의해야 한다!
\end{remark}

\begin{lemma}
\label{properties-lemma-normal-integral-sections}
\begin{slogan}
정규 스킴 위의 함수환은 정규이다.
\end{slogan}
$X$를 정수적 정규 스킴이라 하자.
그러면 $\Gamma(X, \mathcal{O}_X)$는 정규 정역이다.
\end{lemma}

\begin{proof}
$R = \Gamma(X, \mathcal{O}_X)$로 놓자. $R$이 정역임은 명백하다.
그 분수체의 원소 $f = a/b$가 $R$ 위에서 정수적이라고 하자.
$f^d + \sum_{i = 0, \ldots, d - 1} a_i f^i = 0$이고
$a_i \in R$인 관계가 있다고 하자.
$U \subset X$를 공집합이 아닌 아핀 열린집합이라 하자.
$b \in R$은 영이 아니고 $X$는 정수적이므로
$b|_U \in \mathcal{O}_X(U)$ 역시 영이 아니다. 따라서 $a/b$는
$\mathcal{O}_X(U)$의 분수체의 원소이며 $\mathcal{O}_X(U)$ 위에서
정수적이다(실제로 같은 다항식
$f^d + \sum_{i = 0, \ldots, d - 1} a_i|_U f^i = 0$을
$U$ 위에서 사용할 수 있다).
$\mathcal{O}_X(U)$는 정규 정역이므로
(보조정리 \ref{properties-lemma-integral-normal}),
$f_U = (a|_U)/(b|_U) \in \mathcal{O}_X(U)$임을 알 수 있다.
$f_U|_V = f_V$임은 공집합이 아닌 아핀 열린집합들이
$V \subset U \subset X$를 만족할 때마다 명백하다. 따라서 국소 절단 $f_U$들은
$g \in R = \Gamma(X, \mathcal{O}_X)$인 원소로 짜깁기된다.
그러면 $bg$와 $a$는 $\mathcal{O}_X(U)$의 원소로서 위와 같은 모든
$U$에 대하여 일치하므로 $bg = a$이다. 다시 말해, $g$의 상은
$f$이며 이는 $R$의 분수체 안에서의 등식이다.
\end{proof}

\section{코언--매콜리 스킴}
\label{properties-section-Cohen-Macaulay}

\noindent
Algebra의 정의 \ref{algebra-definition-local-ring-CM}에서 보았듯이,
국소 뇌터환 $(R, \mathfrak m)$이 코언--매콜리라는 것은
$\text{depth}_{\mathfrak m}(R) = \dim(R)$이라는 뜻이다.
뇌터환 $R$이 코언--매콜리라는 것은 $R_{\mathfrak p}$ 꼴의 모든
국소환, 즉 $R$의 각 소 아이디얼에서의 국소화가 코언--매콜리라는 뜻이다.
Algebra의 정의 \ref{algebra-definition-ring-CM}를 보라.

\begin{definition}
\label{properties-definition-Cohen-Macaulay}
$X$를 스킴이라 하자. $X$가 {\it 코언--매콜리}라는 것은 모든
$x \in X$에 대하여 아핀 열린집합 $U \subset X$가 존재하여 이것이
$x$의 근방이고 환 $\mathcal{O}_X(U)$가 뇌터이며 코언--매콜리라는 뜻이다.
\end{definition}

\begin{lemma}
\label{properties-lemma-characterize-Cohen-Macaulay}
$X$를 스킴이라 하자. 다음은 서로 동치이다.
\begin{enumerate}
\item $X$는 코언--매콜리이고,
\item $X$는 국소 뇌터이며 그 모든 국소환은 코언--매콜리이고,
\item $X$는 국소 뇌터이며 임의의 폐점 $x \in X$에 대하여 국소환
$\mathcal{O}_{X, x}$는 코언--매콜리이다.
\end{enumerate}
\end{lemma}

\begin{proof}
Algebra의 보조정리 \ref{algebra-lemma-localize-CM}에 의하면 코언--매콜리
국소환의 국소화는 코언--매콜리이다. 이를 보조정리
\ref{properties-lemma-locally-Noetherian}, 국소 뇌터 스킴에서 폐점의 존재
(보조정리 \ref{properties-lemma-locally-Noetherian-closed-point}), 그리고 정의들과
결합하면 보조정리를 얻는다.
\end{proof}

\begin{lemma}
\label{properties-lemma-locally-Cohen-Macaulay}
$X$를 스킴이라 하자. 다음은 서로 동치이다.
\begin{enumerate}
\item 스킴 $X$는 코언--매콜리이다.
\item 모든 아핀 열린집합 $U \subset X$에 대하여 환
$\mathcal{O}_X(U)$는 뇌터이고 코언--매콜리이다.
\item 아핀 열린 덮개 $X = \bigcup U_i$가 존재하여 각
$\mathcal{O}_X(U_i)$가 뇌터이고 코언--매콜리이다.
\item 열린 덮개 $X = \bigcup X_j$가 존재하여 각 열린 부분스킴
$X_j$가 코언--매콜리이다.
\end{enumerate}
더욱이 $X$가 코언--매콜리이면 모든 열린 부분스킴도 코언--매콜리이다.
\end{lemma}

\begin{proof}
보조정리 \ref{properties-lemma-locally-Noetherian}과
\ref{properties-lemma-characterize-Cohen-Macaulay}를 결합하면 된다.
\end{proof}

\noindent
코언--매콜리 스킴과 깊이에 관한 더 자세한 내용은
Cohomology of Schemes의 절 \ref{coherent-section-depth}에서 찾을 수 있다.

\section{정칙 스킴}
\label{properties-section-regular}

\noindent
Algebra의 정의 \ref{algebra-definition-regular-local}에서 보았듯이,
국소 뇌터환 $(R, \mathfrak m)$에서 $\mathfrak m$이
$\dim(R)$개의 원소로 생성되면 이 환을 {\it 정칙}이라고 한다.
뇌터환 $R$이 {\it 정칙}이라는 것은 $R_{\mathfrak p}$ 꼴의 모든
국소환, 즉 $R$의 각 소 아이디얼에서의 국소화가 정칙이라는 뜻이다.
Algebra의 정의 \ref{algebra-definition-regular}를 보라.

\begin{definition}
\label{properties-definition-regular}
$X$를 스킴이라 하자. $X$가 {\it 정칙}, 또는 {\it 비특이}라는 것은
모든 $x \in X$에 대하여 아핀 열린집합 $U \subset X$가 존재하여
이것이 $x$의 근방이고 환 $\mathcal{O}_X(U)$가 뇌터이며 정칙이라는 뜻이다.
\end{definition}

\begin{lemma}
\label{properties-lemma-characterize-regular}
$X$를 스킴이라 하자. 다음은 서로 동치이다.
\begin{enumerate}
\item $X$는 정칙이고,
\item $X$는 국소 뇌터이며 그 모든 국소환은 정칙이고,
\item $X$는 국소 뇌터이며 임의의 폐점 $x \in X$에 대하여 국소환
$\mathcal{O}_{X, x}$는 정칙이다.
\end{enumerate}
\end{lemma}

\begin{proof}
Algebra의 정의 \ref{algebra-definition-regular} 앞에 있는 논의에 의하면
정칙 국소환의 국소화는 정칙이다. 이를 보조정리
\ref{properties-lemma-locally-Noetherian}, 국소 뇌터 스킴에서 폐점의 존재
(보조정리 \ref{properties-lemma-locally-Noetherian-closed-point}), 그리고 정의들과
결합하면 보조정리를 얻는다.
\end{proof}

\begin{lemma}
\label{properties-lemma-locally-regular}
$X$를 스킴이라 하자. 다음은 서로 동치이다.
\begin{enumerate}
\item 스킴 $X$는 정칙이다.
\item 모든 아핀 열린집합 $U \subset X$에 대하여 환
$\mathcal{O}_X(U)$는 뇌터이고 정칙이다.
\item 아핀 열린 덮개 $X = \bigcup U_i$가 존재하여 각
$\mathcal{O}_X(U_i)$가 뇌터이고 정칙이다.
\item 열린 덮개 $X = \bigcup X_j$가 존재하여 각 열린 부분스킴
$X_j$가 정칙이다.
\end{enumerate}
더욱이 $X$가 정칙이면 모든 열린 부분스킴도 정칙이다.
\end{lemma}

\begin{proof}
보조정리 \ref{properties-lemma-locally-Noetherian}과
\ref{properties-lemma-characterize-regular}를 결합하면 된다.
\end{proof}

\begin{lemma}
\label{properties-lemma-regular-normal}
정칙 스킴은 정규이다.
\end{lemma}

\begin{proof}
Algebra의 보조정리 \ref{algebra-lemma-regular-normal}를 보라.
\end{proof}

\section{차원}
\label{properties-section-dimension}

\noindent
스킴의 차원은 바로 그 밑 위상공간의 차원이다.

\begin{definition}
\label{properties-definition-dimension}
$X$를 스킴이라 하자.
\begin{enumerate}
\item $X$의 {\it 차원}은 위상공간으로서의 $X$의 차원이다.
Topology의 정의 \ref{topology-definition-Krull}를 보라.
\item $x \in X$에 대하여 $\dim_x(X)$는 $X$의 밑 위상공간이
$x$에서 갖는 차원을 나타낸다. Topology의 정의
\ref{topology-definition-Krull}를 보라.
$\dim_x(X)$를 {\it $X$의 $x$에서의 차원}이라고 한다.
\end{enumerate}
\end{definition}

\noindent
스킴의 밑 위상공간은 소버이므로
(Schemes의 보조정리 \ref{schemes-lemma-scheme-sober}),
$X$의 차원은 기약 닫힌 부분집합들의 사슬
$$
T_0 \subset T_1 \subset \ldots \subset T_n
$$
의 길이 $n$들의 상한으로 계산할 수 있다. 여기서 이 부분집합들은
$X$ 안에 있다. 또는 특수화 사슬
$$
\xi_n \leadsto \xi_{n - 1} \leadsto \ldots \leadsto \xi_0
$$
의 길이 $n$들의 상한으로 계산할 수도 있다. 여기서 이 점들은
$X$의 점들이다.

\begin{lemma}
\label{properties-lemma-dimension}
$X$를 스킴이라 하자. 다음 세 수는 같다.
\begin{enumerate}
\item $X$의 차원.
\item $X$의 국소환들의 차원들의 상한.
\item $\dim_x(X)$의 상한. 여기서 $x \in X$는 모든 점을 돈다.
\end{enumerate}
\end{lemma}

\begin{proof}
특수화 사슬
$$
\xi_n \leadsto \xi_{n - 1} \leadsto \ldots \leadsto \xi_0
$$
이 주어졌다고 하자. 이 사슬의 점들은 $X$의 점들이며, 모든
$\xi_i$는 Schemes의 보조정리
\ref{schemes-lemma-specialize-points}에 의해 $X$의
$\xi_0$에서의 국소환의 소 아이디얼에 대응한다. 따라서 $X$의 차원은 그
국소환들의 차원들의 상한이다. 특히
$\dim_x(X) \geq \dim(\mathcal{O}_{X, x})$이다. 실제로 $\dim_x(X)$는 $x$의 열린
근방들의 차원들의 최솟값이다. 따라서
$\sup_{x \in X} \dim_x(X) \geq \dim(X)$이다. 한편
$\sup_{x \in X} \dim_x(X) \leq \dim(X)$임은 명백하다. 실제로
임의의 열린 부분집합에 대하여 $\dim(U) \leq \dim(X)$이고 이 열린집합은
$X$ 안에 있다.
\end{proof}

\begin{lemma}
\label{properties-lemma-codimension-local-ring}
$X$를 스킴이라 하자. $Y \subset X$를 기약 닫힌 부분집합이라 하고
$\xi \in Y$를 그 일반점이라 하자. 그러면
$$
\text{codim}(Y, X) = \dim(\mathcal{O}_{X, \xi})
$$
이다. 여기서 여차원은 Topology의 정의
\ref{topology-definition-codimension}와 같다.
\end{lemma}

\begin{proof}
Topology의 보조정리
\ref{topology-lemma-codimension-at-generic-point}에 의해 $X$를
$\xi$의 아핀 열린 근방으로 바꾸어도 된다. 이 경우 결과는
Algebra의 보조정리
\ref{algebra-lemma-irreducible-components-containing-x}에서 쉽게 따른다.
\end{proof}

\begin{lemma}
\label{properties-lemma-generic-point}
$X$를 스킴이라 하고 $x \in X$라 하자. $x$가 $X$의 어떤 기약 성분의
일반점일 필요충분조건은 $\dim(\mathcal{O}_{X, x}) = 0$인 것이다.
\end{lemma}

\begin{proof}
예를 들어 보조정리 \ref{properties-lemma-codimension-local-ring}에서 따른다.
\end{proof}

\begin{lemma}
\label{properties-lemma-locally-Noetherian-dimension-0}
차원이 $0$인 국소 뇌터 스킴은 아르틴 국소환들의 스펙트럼들의
서로소 합집합이다.
\end{lemma}

\begin{proof}
차원이 $0$인 뇌터환은 아르틴 국소환들의 유한 곱이다.
Algebra의 명제 \ref{algebra-proposition-dimension-zero-ring}을 보라.
따라서 국소 뇌터 스킴 $X$의 차원이 $0$이면 그 아핀 열린집합은 이산인
밑 위상공간을 갖는다. 이는 $X$의 위상이 이산임을 뜻하며,
이 관찰들에서 보조정리가 쉽게 따른다.
\end{proof}

\begin{lemma}
\label{properties-lemma-dimension-zero}
\begin{reference}
2016년 6월 4일 자 Ofer Gabber의 전자우편
\end{reference}
$X$를 차원 영인 스킴이라 하자. 다음은 서로 동치이다.
\begin{enumerate}
\item $X$는 준분리이다.
\item $X$는 분리이다.
\item $X$는 하우스도르프이다.
\item 모든 아핀 열린집합은 닫혀 있다.
\end{enumerate}
이 경우 $X$의 연결 성분들은 점이고, $X$의 모든 준콤팩트 열린집합은
아핀이다. 특히 $X$가 준콤팩트이면 $X$는 아핀이다.
\end{lemma}

\begin{proof}
$X$의 차원이 영이므로, 임의의 아핀 열린집합 $U \subset X$에 대하여
공간 $U$는 프로유한이며 아래에서 자유롭게 사용할 여러 다른 성질을
만족한다. Algebra의 보조정리
\ref{algebra-lemma-ring-with-only-minimal-primes}를 보라.
아핀 열린 덮개 $X = \bigcup U_i$를 선택한다.

\medskip\noindent
(4)가 성립한다고 하자. 그러면 $U_i \cap U_j$는 $U_i$의 닫힌
부분집합이므로 준콤팩트이다. 따라서 Schemes의 보조정리
\ref{schemes-lemma-characterize-quasi-separated}에 의해 $X$는
준분리이고, (1)이 성립한다.

\medskip\noindent
(1)이 성립한다고 하자. 그러면 $U_i \cap U_j$는 $U_i$의 준콤팩트
열린집합이므로 $U_i$에서 닫혀 있다. 따라서
$U_i \cap U_j \to U_i$는 상이 닫힌 열린 몰입이고, 따라서 닫힌
몰입이다. 특히 $U_i \cap U_j$는 아핀이며
$\mathcal{O}(U_i) \to \mathcal{O}_X(U_i \cap U_j)$는 전사이다.
그러므로 Schemes의 보조정리
\ref{schemes-lemma-characterize-separated}에 의해 $X$는 분리이고,
(2)가 성립한다.

\medskip\noindent
(2)를 가정하고 $x, y \in X$라 하자. $x \in U_i$라고 하자.
$y \in U_i$이기도 하면 $x$와 $y$의 서로소 열린 근방을 찾을 수 있다.
실제로 $U_i$는 하우스도르프이다. 이제 $y \not \in U_i$이고
$y \in U_j$라고 하자. 그러면 $y \not \in U_i \cap U_j$이다.
이 교집합은 $U_j$의 아핀 열린집합이므로 $U_j$에서 닫혀 있다.
따라서 $y$의 열린 근방 가운데 $U_i$와 만나지 않는 것을 찾을 수 있다.
결국 $X$는 하우스도르프이고, (3)이 성립한다.

\medskip\noindent
(3)을 가정하자. $U \subset X$를 아핀 열린집합이라 하자.
그러면 Topology의 보조정리
\ref{topology-lemma-quasi-compact-in-Hausdorff}에 의해 $U$는 $X$에서
닫혀 있다. 이는 (4)를 증명한다.

\medskip\noindent
$X$가 서로 동치인 조건 (1) -- (4)를 만족한다고 하자.
보조정리의 마지막 주장들을 증명하자. $x, y \in X$이고
$x \not = y$라고 하자. $y$가 $x$로 특수화되지 않으므로
$U \subset X$인 아핀 열린집합을 $x \in U$이고 $y \not \in U$가
되도록 선택할 수 있다. 그러면
$X = U \amalg (X \setminus U)$는 열린닫힌 부분집합들로의 분해이고,
따라서 $x$와 $y$는 $X$의 같은 연결 성분에 속하지 않는다.
다음으로 $U \subset X$를 준콤팩트 열린집합이라 하자.
$U = U_1 \cup \ldots \cup U_n$을 아핀 열린집합들의 합집합으로 쓰자.
$n$에 대한 귀납법으로 $U$가 아핀임을 보일 것이다. 이는 곧바로
$n = 2$인 경우로 귀착된다. 이 경우
$U = (U_1 \setminus U_2) \amalg (U_1 \cap U_2) \amalg (U_2 \setminus U_1)$
이고, 위의 논증에 의해 각 조각은 아핀이다.
\end{proof}

\begin{lemma}
\label{properties-lemma-isolated-dim-0-locally-noetherian}
$x$를 국소 뇌터 스킴 $X$의 점이라 하자. 그러면
$\dim_x(X) = 0$일 필요충분조건은 $x$가 $X$의 고립점인 것이다.
\end{lemma}

\begin{proof}
$x$가 고립점이면 $\{x\}$는 $X$의 열린 부분집합이고
(Topology의 정의 \ref{topology-definition-isolated-point}),
따라서 정의에 의해 $\dim_x(X) = \dim(\{x\}) = 0$이다.
역으로 $\dim_x(X) = 0$이면 열린집합 $U \subset X$가 존재하여
$x$의 근방이고 $\dim(U) = 0$이다
(Topology의 정의 \ref{topology-definition-Krull}).
보조정리 \ref{properties-lemma-locally-Noetherian-dimension-0}에 의해 $U$의
위상은 이산이므로 $x$는 고립점이다.
\end{proof}

\section{카테너리 스킴}
\label{properties-section-catenary}

\noindent
위상공간 $X$가 {\it 카테너리}라는 것은 기약 닫힌 부분집합의 임의의
쌍 $T \subset T'$에 대하여 기약 닫힌 부분집합들의 극대 사슬
$$
T = T_0 \subset T_1 \subset \ldots \subset T_e = T'
$$
이 존재하고 그러한 모든 사슬의 길이가 같다는 뜻임을 상기하자.
Topology의 정의 \ref{topology-definition-catenary}를 보라.

\begin{definition}
\label{properties-definition-catenary}
$S$를 스킴이라 하자. $S$의 밑 위상공간이 카테너리이면
$S$를 {\it 카테너리}라고 한다.
\end{definition}

\noindent
환 $A$가 {\it 카테너리}라는 것은 소 아이디얼의 임의의 쌍
$\mathfrak p \subset \mathfrak q$에 대하여 소 아이디얼들의 극대 사슬
$$
\mathfrak p =
\mathfrak p_0 \subset \ldots \subset \mathfrak p_e
= \mathfrak q
$$
이 존재하고 이들 모두의 길이가 같다는 뜻임을 상기하자.
Algebra의 정의 \ref{algebra-definition-catenary}를 보라.

\begin{lemma}
\label{properties-lemma-catenary-local}
$S$를 스킴이라 하자. 다음은 서로 동치이다.
\begin{enumerate}
\item $S$는 카테너리이고,
\item $S$의 열린 덮개 가운데 모든 원소가 카테너리 스킴인 것이 존재하고,
\item 모든 아핀 열린집합 $\Spec(R) = U \subset S$에 대하여
환 $R$은 카테너리이고,
\item 아핀 열린 덮개 $S = \bigcup U_i$가 존재하여 각 $U_i$가
카테너리 환의 스펙트럼이다.
\end{enumerate}
더욱이 이 경우 $S$의 모든 국소닫힌 부분스킴 역시 카테너리이다.
\end{lemma}

\begin{proof}
Topology의 보조정리 \ref{topology-lemma-catenary}와
Algebra의 보조정리 \ref{algebra-lemma-catenary}를 결합하면 된다.
\end{proof}

\begin{lemma}
\label{properties-lemma-catenary-dimension-function}
$S$를 국소 뇌터 스킴이라 하자. 다음은 서로 동치이다.
\begin{enumerate}
\item $S$는 카테너리이고,
\item Zariski 위상에서 국소적으로 $S$ 위의 차원 함수가 존재한다
(Topology의 정의 \ref{topology-definition-dimension-function}를 보라).
\end{enumerate}
\end{lemma}

\begin{proof}
이는 Topology의 보조정리들
\ref{topology-lemma-catenary},
\ref{topology-lemma-dimension-function-catenary},
\ref{topology-lemma-locally-dimension-function},
Schemes의 보조정리 \ref{schemes-lemma-scheme-sober}, 그리고 마지막으로
보조정리 \ref{properties-lemma-Noetherian-topology}에서 따른다.
\end{proof}

\noindent
스킴이 카테너리일 필요충분조건은 그 국소환들이 카테너리인 것임이 밝혀진다.

\begin{lemma}
\label{properties-lemma-catenary-local-rings-catenary}
$X$를 스킴이라 하자. 다음은 서로 동치이다.
\begin{enumerate}
\item $X$는 카테너리이고,
\item 임의의 $x \in X$에 대하여 국소환 $\mathcal{O}_{X, x}$는
카테너리이다.
\end{enumerate}
\end{lemma}

\begin{proof}
$X$가 카테너리라고 가정하자. $x \in X$라 하자. 보조정리
\ref{properties-lemma-catenary-local}에 의해 $X$를 $x$의 아핀 열린 근방으로
바꾸어도 되며, 그러면 $\Gamma(X, \mathcal{O}_X)$는 카테너리 환이다.
Algebra의 보조정리 \ref{algebra-lemma-localization-catenary}에 의해
카테너리 환의 모든 국소화는 카테너리이다. 따라서
$\mathcal{O}_{X, x}$는 카테너리이다.

\medskip\noindent
역으로 $X$의 모든 국소환이 카테너리라고 가정하자.
$Y \subset Y'$를 $X$의 기약 닫힌 부분집합들의 포함이라 하자.
$\xi \in Y$를 일반점이라 하자.
$\mathfrak p \subset \mathcal{O}_{X, \xi}$를 $Y'$의 일반점에
대응하는 소 아이디얼이라 하자.
Schemes의 보조정리 \ref{schemes-lemma-specialize-points}를 보라.
같은 보조정리에 의해 $X$ 안에서 $Y$와 $Y'$ 사이에 있는 기약 닫힌
부분집합들은 $\mathfrak q \subset \mathcal{O}_{X, \xi}$이면서
$\mathfrak p \subset \mathfrak q \subset \mathfrak m_{\xi}$인
소 아이디얼들에 대응한다. $\mathcal{O}_{X, \xi}$는 카테너리 환이므로
이러한 소 아이디얼들의 모든 극대 사슬은 유한하고 길이가 같다.
\end{proof}

\section{Serre 조건}
\label{properties-section-Rk}

\noindent
자주 유용한 두 가지 기술적 개념을 소개한다.
Cohomology of Schemes의 절 \ref{coherent-section-depth}도 보라.

\begin{definition}
\label{properties-definition-Rk}
$X$를 국소 뇌터 스킴이라 하고 $k \geq 0$이라 하자.
\begin{enumerate}
\item $X$가 {\it 여차원 $k$에서 정칙}이라고 하거나, $X$가
성질 {\it $(R_k)$}를 갖는다고 함은 모든 $x \in X$에 대하여
다음이 성립한다는 뜻이다.
$$
\dim(\mathcal{O}_{X, x}) \leq k
\Rightarrow
\mathcal{O}_{X, x}\text{ is regular}
$$
\item $X$가 성질 {\it $(S_k)$}를 갖는다고 함은 모든 $x \in X$에
대하여
$\text{depth}(\mathcal{O}_{X, x}) \geq \min(k, \dim(\mathcal{O}_{X, x}))$
가 성립한다는 뜻이다.
\end{enumerate}
\end{definition}

\noindent
``여차원 $k$에서 정칙''이라는 표현은 타당하다. 실제로 절
\ref{properties-section-catenary}에서 $Y \subset X$가 일반점 $x$를 갖는 기약
닫힌 부분집합이면
$\dim(\mathcal{O}_{X, x}) = \text{codim}(Y, X)$임을 보았다.
예를 들어 조건 $(R_0)$은 모든 일반점 $\eta \in X$가 $X$의 어떤
기약 성분의 일반점일 때 국소환 $\mathcal{O}_{X, \eta}$가 체라는 뜻이다.
그러나 일반적인 뇌터 스킴에서는 $X$의 정칙 궤적이 잘 거동하지 않을 수
있으므로 주의해야 한다.

\begin{lemma}
\label{properties-lemma-scheme-regular-iff-all-Rk}
$X$를 국소 뇌터 스킴이라 하자. 그러면 $X$가 정칙일 필요충분조건은
$X$가 $(R_k)$를 모든 $k \geq 0$에 대하여 갖는 것이다.
\end{lemma}

\begin{proof}
보조정리 \ref{properties-lemma-characterize-regular}과 정의들에서 따른다.
\end{proof}

\begin{lemma}
\label{properties-lemma-scheme-CM-iff-all-Sk}
$X$를 국소 뇌터 스킴이라 하자. 그러면 $X$가 코언--매콜리일
필요충분조건은 $X$가 $(S_k)$를 모든 $k \geq 0$에 대하여 갖는 것이다.
\end{lemma}

\begin{proof}
보조정리 \ref{properties-lemma-characterize-Cohen-Macaulay}에 의해 국소환을
살피는 문제로 귀착된다. 뇌터 국소환이 코언--매콜리일 필요충분조건은
그 깊이가 차원과 같은 것이므로 보조정리가 성립한다.
\end{proof}

\begin{lemma}
\label{properties-lemma-criterion-reduced}
$X$를 국소 뇌터 스킴이라 하자. 그러면 $X$가 축약일 필요충분조건은
$X$가 성질 $(S_1)$과 $(R_0)$을 갖는 것이다.
\end{lemma}

\begin{proof}
이는 Algebra의 보조정리 \ref{algebra-lemma-criterion-reduced}이다.
\end{proof}

\begin{lemma}
\label{properties-lemma-criterion-normal}
$X$를 국소 뇌터 스킴이라 하자. 그러면 $X$가 정규일 필요충분조건은
$X$가 성질 $(S_2)$와 $(R_1)$을 갖는 것이다.
\end{lemma}

\begin{proof}
이는 Algebra의 보조정리 \ref{algebra-lemma-criterion-normal}이다.
\end{proof}

\begin{lemma}
\label{properties-lemma-normal-dimension-1-regular}
$X$를 정규인 국소 뇌터 스킴이라 하고 그 차원이 $\leq 1$이라고 하자.
그러면 $X$는 정칙이다.
\end{lemma}

\begin{proof}
보조정리 \ref{properties-lemma-criterion-normal}과 정의들에서 따른다.
\end{proof}

\begin{lemma}
\label{properties-lemma-normal-dimension-2-Cohen-Macaulay}
$X$를 정규인 국소 뇌터 스킴이라 하고 그 차원이 $\leq 2$라고 하자.
그러면 $X$는 코언--매콜리이다.
\end{lemma}

\begin{proof}
보조정리 \ref{properties-lemma-criterion-normal}과 정의들에서 따른다.
\end{proof}

\section{Japanese 스킴과 Nagata 스킴}
\label{properties-section-nagata}

\noindent
이 절에서 다루는 개념들은 EGA에서 두드러지게 정의되어 있지 않다.
``보편 Japanese 스킴''은 \cite[IV Corollary 5.11.4]{EGA}에서 언급되고
정의되어 있다. ``Japanese 스킴''은
\cite[IV Remark 10.4.14 (ii)]{EGA}에서 언급되지만 정의는 주어지지 않는다.
아래와 같은 Nagata 스킴은 문헌의 몇 곳에서 등장한다
(예를 들어 \cite[Definition 8.2.30]{Liu}와 \cite[Page 142]{Greco}를 보라).

\medskip\noindent
정역 $R$이 {\it Japanese}라는 것은 그 분수체의 임의의 유한 확대에서
$R$의 정수적 폐포가 $R$ 위에서 유한이라는 뜻임을 간단히 상기하자.
환 $R$이 {\it 보편 Japanese}라는 것은 임의의 유한형 환 사상
$R \to S$에 대하여 $S$가 정역이면 $S$가 Japanese라는 뜻이다.
환 $R$이 {\it Nagata}라는 것은 이 환이 뇌터이고
$R/\mathfrak p$가 모든 소 아이디얼 $\mathfrak p$에 대하여
Japanese라는 뜻이다. 여기서 이 소 아이디얼들은 $R$의 소 아이디얼이다.

\begin{definition}
\label{properties-definition-nagata}
$X$를 스킴이라 하자.
\begin{enumerate}
\item $X$가 정수적이라고 가정하자. $X$가 {\it Japanese}라는 것은
모든 $x \in X$에 대하여 아핀 열린 근방 $x \in U \subset X$가
존재하여 환 $\mathcal{O}_X(U)$가 Japanese라는 뜻이다
(Algebra의 정의 \ref{algebra-definition-N}을 보라).
\item $X$가 {\it 보편 Japanese}라는 것은 모든 $x \in X$에 대하여
아핀 열린 근방 $x \in U \subset X$가 존재하여 환
$\mathcal{O}_X(U)$가 보편 Japanese라는 뜻이다
(Algebra의 정의 \ref{algebra-definition-nagata}를 보라).
\item $X$가 {\it Nagata}라는 것은 모든 $x \in X$에 대하여 아핀 열린
근방 $x \in U \subset X$가 존재하여 환 $\mathcal{O}_X(U)$가
Nagata라는 뜻이다
(Algebra의 정의 \ref{algebra-definition-nagata}를 보라).
\end{enumerate}
\end{definition}

\noindent
Nagata라는 것은 국소 뇌터이고 보편 Japanese라는 것과 같다.
보조정리 \ref{properties-lemma-nagata-universally-Japanese}를 보라.

\begin{remark}
\label{properties-remark-non-integral-Japanese}
\cite{Hoobler-finite}에서는 (국소 뇌터) 스킴 $X$가 Japanese라는 것을
모든 $x \in X$와 모든 동반 소 아이디얼 $\mathfrak p$에 대하여,
이것이 $\mathcal{O}_{X, x}$의 동반 소 아이디얼일 때 환
$\mathcal{O}_{X, x}/\mathfrak p$가
Japanese라는 것으로 정의한다. 우리는 이 정의를 사용하지 않는다.
실제로 차원이 일이고 모든 국소환이 훌륭하여(특히 Japanese여서)
정규화가 유한이 아닌 뇌터 정역이 존재하기 때문이다.
\cite[Example 1]{Hochster-loci}, \cite{Heinzer-Levy}, 또는
\cite[Expos\'e XIX]{Traveaux}를 보라.
한편 스킴 $X$를 Japanese라고 부르는 방식으로 이 문제를 피할 수도 있다.
즉 모든 아핀 열린집합 $\Spec(A) \subset X$에 대하여
$A/\mathfrak p$가 모든 동반 소 아이디얼 $\mathfrak p$, 곧 $A$의
동반 소 아이디얼에 대해 Japanese라고 요구할 수 있다.
\end{remark}

\begin{lemma}
\label{properties-lemma-nagata-locally-Noetherian}
Nagata 스킴은 국소 뇌터이다.
\end{lemma}

\begin{proof}
Nagata 환은 정의상 뇌터이므로 성립한다.
\end{proof}

\begin{lemma}
\label{properties-lemma-locally-Japanese}
$X$를 정수적 스킴이라 하자. 다음은 서로 동치이다.
\begin{enumerate}
\item 스킴 $X$는 Japanese이다.
\item 모든 아핀 열린집합 $U \subset X$에 대하여 정역
$\mathcal{O}_X(U)$는 Japanese이다.
\item 아핀 열린 덮개 $X = \bigcup U_i$가 존재하여 각
$\mathcal{O}_X(U_i)$가 Japanese이다.
\item 열린 덮개 $X = \bigcup X_j$가 존재하여 각 열린 부분스킴
$X_j$가 Japanese이다.
\end{enumerate}
더욱이 $X$가 Japanese이면 모든 열린 부분스킴도 Japanese이다.
\end{lemma}

\begin{proof}
보조정리 \ref{properties-lemma-locally-P}와 Algebra의 보조정리들
\ref{algebra-lemma-localize-N},
\ref{algebra-lemma-Japanese-local}에서 따른다.
\end{proof}

\begin{lemma}
\label{properties-lemma-locally-universally-Japanese}
$X$를 스킴이라 하자. 다음은 서로 동치이다.
\begin{enumerate}
\item 스킴 $X$는 보편 Japanese이다.
\item 모든 아핀 열린집합 $U \subset X$에 대하여 환
$\mathcal{O}_X(U)$는 보편 Japanese이다.
\item 아핀 열린 덮개 $X = \bigcup U_i$가 존재하여 각
$\mathcal{O}_X(U_i)$가 보편 Japanese이다.
\item 열린 덮개 $X = \bigcup X_j$가 존재하여 각 열린 부분스킴
$X_j$가 보편 Japanese이다.
\end{enumerate}
더욱이 $X$가 보편 Japanese이면 모든 열린 부분스킴도 보편 Japanese이다.
\end{lemma}

\begin{proof}
보조정리 \ref{properties-lemma-locally-P}와 Algebra의 보조정리들
\ref{algebra-lemma-universally-japanese},
\ref{algebra-lemma-nagata-local}에서 따른다.
\end{proof}

\begin{lemma}
\label{properties-lemma-locally-nagata}
$X$를 스킴이라 하자. 다음은 서로 동치이다.
\begin{enumerate}
\item 스킴 $X$는 Nagata이다.
\item 모든 아핀 열린집합 $U \subset X$에 대하여 환
$\mathcal{O}_X(U)$는 Nagata이다.
\item 아핀 열린 덮개 $X = \bigcup U_i$가 존재하여 각
$\mathcal{O}_X(U_i)$가 Nagata이다.
\item 열린 덮개 $X = \bigcup X_j$가 존재하여 각 열린 부분스킴
$X_j$가 Nagata이다.
\end{enumerate}
더욱이 $X$가 Nagata이면 모든 열린 부분스킴도 Nagata이다.
\end{lemma}

\begin{proof}
보조정리 \ref{properties-lemma-locally-P}와 Algebra의 보조정리들
\ref{algebra-lemma-nagata-localize},
\ref{algebra-lemma-nagata-local}에서 따른다.
\end{proof}

\begin{lemma}
\label{properties-lemma-characterize-nagata}
$X$를 국소 뇌터 스킴이라 하자. 그러면 $X$가 Nagata일 필요충분조건은
모든 정수적 닫힌 부분스킴 $Z \subset X$가 Japanese인 것이다.
\end{lemma}

\begin{proof}
$X$가 Nagata라고 가정하자. $Z \subset X$를 정수적 닫힌 부분스킴이라
하고 $z \in Z$라 하자. 아핀 열린집합
$\Spec(A) = U \subset X$를 $z$를 포함하고 $A$가 Nagata가 되도록
선택하자. 그러면 $Z \cap U \cong \Spec(A/\mathfrak p)$이고,
여기서 $\mathfrak p$는 어떤 소 아이디얼이다.
Schemes의 보조정리 \ref{schemes-lemma-closed-subspace-scheme}와
정의 \ref{properties-definition-integral}을 보라. Algebra의 정의
\ref{algebra-definition-nagata}에 의해 $A/\mathfrak p$는 Japanese이다.
따라서 정의에 의해 $Z$는 Japanese이다.

\medskip\noindent
$X$의 모든 정수적 닫힌 부분스킴이 Japanese라고 가정하자.
$\Spec(A) = U \subset X$를 임의의 아핀 열린집합이라 하자.
$X$가 국소 뇌터이므로 $A$는 뇌터이다
(보조정리 \ref{properties-lemma-locally-Noetherian}). $\mathfrak p \subset A$를
소 아이디얼이라 하자. $A/\mathfrak p$가 Japanese임을 보여야 한다.
$T \subset U$를 닫힌 부분집합 $V(\mathfrak p) \subset \Spec(A)$라 하자.
$\overline{T} \subset X$를 그 폐포라 하자. 그러면 $\overline{T}$는
기약 부분집합의 폐포이므로 기약이다. 따라서 $\overline{T}$가 정의하는
축약 닫힌 부분스킴은 정수적 닫힌 부분스킴이다(이를 다시
$\overline{T}$라 부른다). Schemes의 보조정리
\ref{schemes-lemma-reduced-closed-subscheme}를 보라.
다시 말해 $\Spec(A/\mathfrak p)$는 $X$의 어떤 정수적 닫힌
부분스킴의 아핀 열린집합이다. 가정에 의해 이 부분스킴은 Japanese이고,
보조정리 \ref{properties-lemma-locally-Japanese}에 의해 $A/\mathfrak p$는
Japanese이다.
\end{proof}

\begin{lemma}
\label{properties-lemma-nagata-universally-Japanese}
$X$를 스킴이라 하자. 다음은 서로 동치이다.
\begin{enumerate}
\item $X$는 Nagata이고,
\item $X$는 국소 뇌터이며 보편 Japanese이다.
\end{enumerate}
\end{lemma}

\begin{proof}
이는 Algebra의 명제
\ref{algebra-proposition-nagata-universally-japanese}이다.
\end{proof}

\noindent
이 논의는 Morphisms의 절 \ref{morphisms-section-nagata}에서 이어진다.

\section{G-스킴}
\label{properties-section-G-scheme}

\noindent
환 $R$이 G-환이라는 것은 $R$이 뇌터이고 모든 소 아이디얼
$\mathfrak p$, 즉 $R$의 소 아이디얼에 대하여 환 사상
$R_\mathfrak p \to (R_\mathfrak p)^\wedge$가 정칙이라는 뜻임을
언급해 둔다.

\begin{definition}
\label{properties-definition-G-scheme}
$X$를 스킴이라 하자. $X$를 {\it G-스킴}\footnote{이는 표준 용어가
아닐 수 있다. $G$가
유한군 또는 군 스킴이면 $G$의 작용을 갖춘 스킴을 $G$-스킴이라고
부르기도 한다.}이라고 하는 것은 모든 $x \in X$에 대하여 아핀 열린
근방 $x \in U \subset X$가 존재하여 환 $\mathcal{O}_X(U)$가
G-환이라는 뜻이다
(More on Algebra의 정의 \ref{more-algebra-definition-G-ring}을 보라).
\end{definition}

\begin{lemma}
\label{properties-lemma-characterize-G-scheme}
$X$를 스킴이라 하자. 다음은 서로 동치이다.
\begin{enumerate}
\item $X$는 G-스킴이고,
\item $X$는 국소 뇌터이며 모든 $x \in X$에 대하여 환 사상
$\mathcal{O}_{X, x} \to \mathcal{O}_{X, x}^\wedge$가 정칙이고,
\item $X$는 국소 뇌터이며 임의의 폐점 $x \in X$에 대하여 환 사상
$\mathcal{O}_{X, x} \to \mathcal{O}_{X, x}^\wedge$가 정칙이다.
\end{enumerate}
\end{lemma}

\begin{proof}
(1)이 (2)를 함의하고 (2)가 (3)을 함의함은 명백하다.
(3)을 가정하자. $x \in X$를 임의의 점이라 하자. 보조정리
\ref{properties-lemma-locally-Noetherian-closed-point}에 의해 특수화
$x \leadsto x'$이 존재하며 $x'$은 $X$에서 닫힌점이다. 그러면
More on Algebra의 보조정리
\ref{more-algebra-lemma-check-G-ring-maximal-ideals}에 의해
$\mathcal{O}_{X, x'}$는 G-환이다.
$\mathcal{O}_{X, x}$는 $\mathcal{O}_{X, x'}$를 어떤 소 아이디얼에서
국소화한 것이므로(Schemes의 보조정리
\ref{schemes-lemma-specialize-points}),
$\mathcal{O}_{X, x} \to \mathcal{O}_{X, x}^\wedge$는 정칙 환 사상이다.
$x \in X$가 임의였으므로, 임의의 아핀 열린집합 $U \subset X$에 대하여
뇌터환(보조정리 \ref{properties-lemma-locally-Noetherian})
$\mathcal{O}_X(U)$는 G-환이라고 결론내린다. 이는 (1)을 증명한다.
\end{proof}

\begin{lemma}
\label{properties-lemma-locally-G-scheme}
$X$를 스킴이라 하자. 다음은 서로 동치이다.
\begin{enumerate}
\item 스킴 $X$는 G-스킴이다.
\item 모든 아핀 열린집합 $U \subset X$에 대하여 환
$\mathcal{O}_X(U)$는 G-환이다.
\item 아핀 열린 덮개 $X = \bigcup U_i$가 존재하여 각
$\mathcal{O}_X(U_i)$가 G-환이다.
\item 열린 덮개 $X = \bigcup X_j$가 존재하여 각 열린 부분스킴
$X_j$가 G-스킴이다.
\end{enumerate}
더욱이 $X$가 G-스킴이면 모든 열린 부분스킴도 G-스킴이다.
\end{lemma}

\begin{proof}
보조정리 \ref{properties-lemma-locally-Noetherian}과
\ref{properties-lemma-characterize-G-scheme}를 결합하면 된다.
\end{proof}

\section{특이 궤적}
\label{properties-section-singular-locus}

\noindent
정의는 다음과 같다.

\begin{definition}
\label{properties-definition-singular-locus}
$X$를 국소 뇌터 스킴이라 하자. {\it 정칙 궤적}
$\text{Reg}(X)$이라 하는 $X$의 부분집합은 점 $x \in X$ 중에서
국소환 $\mathcal{O}_{X, x}$가 정칙 국소환인 것들의 집합이다.
{\it 특이 궤적} $\text{Sing}(X)$은 여집합
$X \setminus \text{Reg}(X)$, 즉 점 $x \in X$ 중에서 국소환
$\mathcal{O}_{X, x}$가 정칙 국소환이 아닌 것들의 집합이다.
\end{definition}

\noindent
국소 뇌터 스킴의 정칙 궤적은 일반화에 대하여 안정적이다.
Algebra의 정의 \ref{algebra-definition-regular} 앞의 논의를 보라.
그러나 일반적인 국소 뇌터 스킴에서는 정칙 궤적이 열려 있을 필요가 없다.
이것이 성립하는 몇 가지 판정법은 More on Algebra의 절
\ref{more-algebra-section-singular-locus}에서 찾을 수 있다.
이 문제는 Morphisms의 절 \ref{morphisms-section-singular-locus}에서
더 논의한다.

\section{국소 기약성}
\label{properties-section-unibranch}

\noindent
More on Algebra의 절 \ref{more-algebra-section-unibranch}에서
(기하학적으로) 단일분지인 국소환의 개념을 도입했음을 상기하자.

\begin{definition}
\label{properties-definition-unibranch}
\begin{reference}
\cite[Chapter IV (6.15.1)]{EGA4}
\end{reference}
$X$를 스킴이라 하고 $x \in X$라 하자.
$X$가 {\it $x$에서 단일분지}라는 것은 국소환
$\mathcal{O}_{X, x}$가 단일분지라는 뜻이다.
$X$가 {\it $x$에서 기하학적으로 단일분지}라는 것은 국소환
$\mathcal{O}_{X, x}$가 기하학적으로 단일분지라는 뜻이다.
$X$가 모든 점에서 단일분지이면 $X$를 {\it 단일분지}라고 한다.
$X$가 모든 점에서 기하학적으로 단일분지이면
$X$를 {\it 기하학적으로 단일분지}라고 한다.
\end{definition}

\noindent
국소환 $A$가 More on Algebra의 정의
\ref{more-algebra-definition-unibranch}의 의미에서 기하학적으로
단일분지이지만, 스킴 $\Spec(A)$는 정의
\ref{properties-definition-unibranch}의 의미에서 기하학적으로 단일분지가 아닌
일이 실제로 일어날 수 있다. 예를 들어 $A$가 보통 이중점 특이점
(node)을 갖는 기약 평면곡선 위 원뿔의 꼭짓점에서의 국소환이면 그렇다.

\begin{lemma}
\label{properties-lemma-normal-geometrically-unibranch}
정규 스킴은 기하학적으로 단일분지이다.
\end{lemma}

\begin{proof}
정의에서 따른다. 실제로 스킴이 정규라는 것은 그 국소환들이 정규
정역이라는 뜻이다. More on Algebra의 정의
\ref{more-algebra-definition-unibranch}에서 정규인 국소 정역이
기하학적으로 단일분지임은 즉시 따른다.
\end{proof}

\begin{lemma}
\label{properties-lemma-geometrically-unibranch}
\begin{reference}
\cite[Proposition 2.3]{Etale-coverings}와 비교하라
\end{reference}
$X$를 뇌터 스킴이라 하자. 다음은 서로 동치이다.
\begin{enumerate}
\item $X$는 기하학적으로 단일분지이고
(정의 \ref{properties-definition-unibranch}),
\item 모든 점 $x \in X$ 가운데 $X$의 기약 성분의 일반점이 아닌 것에 대하여
강한 헨젤화 $\mathcal{O}_{X, x}^{sh}$의 구멍 뚫린 스펙트럼은 연결이다.
\end{enumerate}
\end{lemma}

\begin{proof}
More on Algebra의 보조정리
\ref{more-algebra-lemma-geometrically-unibranch}에 의해 (1)이 성립하면
(2)의 구멍 뚫린 스펙트럼들은 기약이고, 특히 연결이다.

\medskip\noindent
(2)를 가정하자. $x \in X$라 하자. $\mathcal{O}_{X, x}$가
기하학적으로 단일분지임을 보여야 한다.
$\dim(\mathcal{O}_{X, x})$에 대한 귀납법으로 $x$의 모든 자명하지 않은
일반화에 대하여 결과가 성립한다고 가정할 수 있다. $X$를
$\Spec(\mathcal{O}_{X, x})$로 바꾸어도 된다. 다시 말해
$X = \Spec(A)$이고 $A$가 국소환이며, $A_\mathfrak p$가 기하학적으로
단일분지라고 가정할 수 있다. 여기서 $\mathfrak p \subset A$는 각
비극대 소 아이디얼을 돈다.

\medskip\noindent
$A^{sh}$를 $A$의 강한 헨젤화라 하자.
$\mathfrak q \subset A^{sh}$가 $\mathfrak p \subset A$ 위에 놓인
소 아이디얼이면 $A_\mathfrak p \to A^{sh}_\mathfrak q$는
\'etale 대수들의 여과 쌍극한이다. 따라서 $A_\mathfrak p$와
$A^{sh}_\mathfrak q$의 강한 헨젤화들은 동형이다.
그러므로 More on Algebra의 보조정리
\ref{more-algebra-lemma-geometrically-unibranch}에 의해
$A^{sh}_\mathfrak q$는 모든 비극대 소 아이디얼 $\mathfrak q$, 즉
$A^{sh}$의 비극대 소 아이디얼에 대하여 유일한 극소 소 아이디얼을 갖는다.

\medskip\noindent
$\mathfrak q_1, \ldots, \mathfrak q_r$를 $A^{sh}$의 극소 소 아이디얼들이라
하자. $r = 1$임을 보여야 한다. 위의 논의에 의해
$V(\mathfrak q_1) \cap V(\mathfrak q_j) = \{\mathfrak m^{sh}\}$가
$j = 2, \ldots, r$에 대하여 성립한다.
따라서 $V(\mathfrak q_1) \setminus \{\mathfrak m^{sh}\}$는
$A^{sh}$의 구멍 뚫린 스펙트럼의 열린닫힌 부분집합이다.
이 구멍 뚫린 스펙트럼이 연결이라는 가정과 모순되지 않으려면
$r = 1$이어야 한다.
\end{proof}

\begin{definition}
\label{properties-definition-number-of-branches}
$X$를 스킴이라 하고 $x \in X$라 하자. {\it $X$의 $x$에서의 분지 수}는
국소환 $\mathcal{O}_{X, x}$의 분지 수이다. 이는 More on Algebra의
정의 \ref{more-algebra-definition-number-of-branches}에서 정의되었다.
{\it $X$의 $x$에서의 기하학적 분지 수}는 국소환
$\mathcal{O}_{X, x}$의 기하학적 분지 수이다. 이는 More on Algebra의
정의 \ref{more-algebra-definition-number-of-branches}에서 정의되었다.
\end{definition}

\noindent
흔히 이를 완비 국소환의 분지들과 비교하고 싶지만 일반적으로 그 비교는
간단하지 않다. 이에 관한 정보는 More on Algebra의 절
\ref{more-algebra-section-branches-completion}에서 찾을 수 있다.

\begin{lemma}
\label{properties-lemma-number-of-branches-irreducible-components}
$X$를 스킴이라 하고 $x \in X$라 하자. $X_i$, $i \in I$를 $X$의
기약 성분 중 $x$를 지나는 것들이라 하자. 그러면 $X$의 $x$에서의
(기하학적) 분지 수는 $i \in I$에 걸쳐 $X_i$의 $x$에서의 (기하학적) 분지 수를
합한 것과 같다.
\end{lemma}

\begin{proof}
$X_i$들을 $X$의 정수적 닫힌 부분스킴으로 본다.
Schemes의 정의 \ref{schemes-definition-reduced-induced-scheme}와
보조정리 \ref{properties-lemma-characterize-integral}을 보라.
$X_i$의 $x$에서의 (기하학적) 분지 수는 적어도 $1$임을 관찰하자.
이는 모든 $i$에 대하여 성립한다(본질적으로 정의에서 따른다).
$X_i$들은 $1$-대-$1$로 극소 소 아이디얼
$\mathfrak p_i \subset \mathcal{O}_{X, x}$들과 대응함을
상기하자. Algebra의 보조정리
\ref{algebra-lemma-irreducible-components-containing-x}를 보라.
따라서 $I$가 무한이면 $\mathcal{O}_{X, x}$는 무한히 많은 극소
소 아이디얼을 갖는다. 그러므로 $\mathcal{O}_{X, x}^h$와
$\mathcal{O}_{X, x}^{sh}$도 모두 무한히 많은 극소 소 아이디얼을 갖는다
(Algebra의 보조정리
\ref{algebra-lemma-injective-minimal-primes-in-image}와
\ref{algebra-lemma-minimal-prime-image-minimal-prime}, 그리고 사상
$\mathcal{O}_{X, x} \to  \mathcal{O}_{X, x}^h \to \mathcal{O}_{X, x}^{sh}$
의 단사성을 결합하라). 이 경우 $X$의 $x$에서의 (기하학적) 분지 수는
$\infty$로 정의되며, 이 값은 합에 대해서도 같다. 따라서 $I$가
유한이라고 가정할 수 있다. $A'$을 $\mathcal{O}_{X, x}$의 정수적
폐포라 하되, 이는 전분수환 $Q$, 즉
$(\mathcal{O}_{X, x})_{red}$의 전분수환 안에서 취한다.
$A'_i$를 $\mathcal{O}_{X, x}/\mathfrak p_i$의 정수적 폐포라 하되,
이는 전분수환 $Q_i$, 즉 $\mathcal{O}_{X, x}/\mathfrak p_i$의
전분수환 안에서 취한다. Algebra의 보조정리
\ref{algebra-lemma-total-ring-fractions-no-embedded-points}에 의해
$Q = \prod_{i \in I} Q_i$이다. 따라서 $A' = \prod A'_i$이다.
이제 More on Algebra의 보조정리
\ref{more-algebra-lemma-number-of-branches-1}는 (기하학적) 분지 수를
$A'$의 극대 아이디얼들로 나타내므로, 그 보조정리에서 원하는 등식이 따른다.
\end{proof}

\begin{lemma}
\label{properties-lemma-number-of-branches-1}
$X$를 스킴이라 하고 $x \in X$라 하자.
\begin{enumerate}
\item $X$의 $x$에서의 분지 수가 $1$일 필요충분조건은
$X$가 $x$에서 단일분지인 것이다.
\item $X$의 $x$에서의 기하학적 분지 수가 $1$일 필요충분조건은
$X$가 $x$에서 기하학적으로 단일분지인 것이다.
\end{enumerate}
\end{lemma}

\begin{proof}
정의들과 환에 대한 대응 결과에서 즉시 따른다. More on Algebra의
보조정리 \ref{more-algebra-lemma-number-of-branches-1}을 보라.
\end{proof}

\section{유한형 가군과 유한 표시 가군의 특징짓기}
\label{properties-section-characterizing-finite-type-presentation}

\noindent
$X$를 스킴이라 하자.
$\mathcal{F}$를 준연접 $\mathcal{O}_X$-가군이라 하자.
다음 보조정리는 $\mathcal{F}$가 유한형일 필요충분조건이
(Modules의 정의 \ref{modules-definition-finite-type}를 보라)
$\mathcal{F}$가 각 아핀 열린집합 $\Spec(A) = U \subset X$ 위에서
어떤 $\widetilde M$ 꼴이고, 여기서 $A$-가군 $M$이 유한인 것임을 보인다.
마찬가지로 $\mathcal{F}$가 유한 표시일 필요충분조건은
(Modules의 정의 \ref{modules-definition-finite-presentation}을 보라)
$\mathcal{F}$가 각 아핀 열린집합 $\Spec(A) = U \subset X$ 위에서
어떤 $\widetilde M$ 꼴이고, 여기서 $A$-가군 $M$이 유한 표시인 것이다.

\begin{lemma}
\label{properties-lemma-finite-type-module}
$X = \Spec(R)$을 아핀 스킴이라 하자. $\mathcal{O}_X$-가군의 준연접층
$\widetilde M$이 유한형 $\mathcal{O}_X$-가군일 필요충분조건은
$M$이 유한 $R$-가군인 것이다.
\end{lemma}

\begin{proof}
$\widetilde M$이 유한형 $\mathcal{O}_X$-가군이라고 가정하자.
이는 $X$의 어떤 열린 덮개가 존재하여 $\widetilde M$을 그 덮개의
원소들에 제한하면 유한 개의 절단으로 전역 생성된다는 뜻이다.
따라서 표준 열린 덮개
$X = \bigcup_{i = 1, \ldots, n} D(f_i)$도 존재하여
$\widetilde M|_{D(f_i)}$가 유한 개의 절단으로 생성된다.
그러므로 $M_{f_i}$는 각 $i$에 대하여 유한 생성이다.
이제 Algebra의 보조정리 \ref{algebra-lemma-cover}에서 결론이 따른다.
\end{proof}

\begin{lemma}
\label{properties-lemma-finite-presentation-module}
$X = \Spec(R)$을 아핀 스킴이라 하자. $\mathcal{O}_X$-가군의 준연접층
$\widetilde M$이 유한 표시 $\mathcal{O}_X$-가군일 필요충분조건은
$M$이 유한 표시 $R$-가군인 것이다.
\end{lemma}

\begin{proof}
$\widetilde M$이 유한 표시 $\mathcal{O}_X$-가군이라고 가정하자.
보조정리 \ref{properties-lemma-finite-type-module}에 의해 $M$은 유한
$R$-가군이다. 전사사상 $R^n \to M$을 선택하고 그 핵을 $K$라 하자.
Schemes의 보조정리 \ref{schemes-lemma-spec-sheaves}에 의해 짧은 완전열
$$
0 \to \widetilde{K} \to
\bigoplus \mathcal{O}_X^{\oplus n} \to
\widetilde{M} \to 0
$$
이 존재한다. Modules의 보조정리
\ref{modules-lemma-kernel-surjection-finite-free-onto-finite-presentation}에
의해 $\widetilde{K}$는 유한형 $\mathcal{O}_X$-가군이다.
따라서 보조정리 \ref{properties-lemma-finite-type-module}를 다시 사용하면
$K$는 유한 $R$-가군이다. 그러므로 $M$은 유한 표시 $R$-가군이다.
\end{proof}







\section{주 열린집합 위의 절단}
\label{properties-section-principal-opens}

\noindent
다음은 이러한 종류의 전형적인 결과이다. 뒤의 보조정리들에서는
더 소박하지만 더 직접적인 증명 방법을 사용할 것이다.

\begin{lemma}
\label{properties-lemma-invert-f-sections}
\begin{slogan}
준연접층의 절단은 무한대에서 유리형 특이점만 가진다.
\end{slogan}
$X$를 스킴이라 하자. $f \in \Gamma(X, \mathcal{O}_X)$라 하자.
$X_f \subset X$를 $f$가 가역인 열린집합이라 쓰자. Schemes의
보조정리 \ref{schemes-lemma-f-open}를 보라.
$X$가 준콤팩트이고 준분리이면 표준 사상
$$
\Gamma(X, \mathcal{O}_X)_f \longrightarrow \Gamma(X_f, \mathcal{O}_X)
$$
은 동형사상이다. 또한 $\mathcal{F}$가 $\mathcal{O}_X$-가군의
준연접층이면 사상
$$
\Gamma(X, \mathcal{F})_f \longrightarrow \Gamma(X_f, \mathcal{F})
$$
은 동형사상이다.
\end{lemma}

\begin{proof}
$R = \Gamma(X, \mathcal{O}_X)$라 쓰자. 스킴의 표준 사상
$$
\varphi : X \longrightarrow \Spec(R)
$$
을 생각하자. Schemes의 보조정리
\ref{schemes-lemma-morphism-into-affine}를 보라.
그러면 오른쪽의 표준 열린집합 $D(f)$의 역상은 왼쪽의
$X_f$이다. 또한 $X$가 준콤팩트이고 준분리라고 가정했으므로
사상 $\varphi$는 준콤팩트이고 준분리이다. Schemes의 보조정리
\ref{schemes-lemma-quasi-compact-affine}와
\ref{schemes-lemma-compose-after-separated}를 보라. 따라서
Schemes의 보조정리 \ref{schemes-lemma-push-forward-quasi-coherent}에
의해 $\varphi_*\mathcal{F}$는 준연접이다.
그러므로 $\varphi_*\mathcal{F} = \widetilde M$이고,
$M = \Gamma(X, \mathcal{F})$는 $R$-가군이다. 따라서
$$
\Gamma(X_f, \mathcal{F}) =
\Gamma(D(f), \varphi_*\mathcal{F}) =
\Gamma(D(f), \widetilde M) = M_f
$$
이고, 이것이 바로 보조정리의 내용이다. 보조정리의 첫 번째 표시
동형사상은 $\mathcal{F} = \mathcal{O}_X$로 놓으면 따른다.
\end{proof}

\noindent
$X$를 스킴이라 하고, $\mathcal{L}$을 $X$ 위의 가역층이라 하며,
$\mathcal{O}_X$-가군의 층 $\mathcal{F}$가 주어지면 등급환
$\Gamma_*(X, \mathcal{L}) =
\bigoplus\nolimits_{n \geq 0} \Gamma(X, \mathcal{L}^{\otimes n})$
과 등급 $\Gamma_*(X, \mathcal{L})$-가군
$\Gamma_*(X, \mathcal{L}, \mathcal{F}) =
\bigoplus\nolimits_{n \in \mathbf{Z}}
\Gamma(X, \mathcal{F} \otimes_{\mathcal{O}_X} \mathcal{L}^{\otimes n})$
을 얻는다는 것을 상기하자. Modules의 정의
\ref{modules-definition-gamma-star}를 보라.
더 나아가 절단 $s \in \Gamma(X, \mathcal{L})$가 있으면 사상
\begin{equation}
\label{properties-equation-module-invert-s}
\Gamma_*(X, \mathcal{L}, \mathcal{F})_{(s)}
\longrightarrow
\Gamma(X_s, \mathcal{F}|_{X_s})
\end{equation}
을 얻는다. 이 사상은 $t/s^n$을 보낸다. 여기서
$t \in \Gamma(X, \mathcal{F} \otimes_{\mathcal{O}_X} \mathcal{L}^{\otimes n})$
이고, 그 상은 $t|_{X_s} \otimes s|_{X_s}^{-n}$이다.
정의상 $X_s \subset X$는 $s$가 역원을 갖는 열린집합이므로 이는
잘 정의된다. Modules의 보조정리 \ref{modules-lemma-s-open}를 보라.

\begin{lemma}
\label{properties-lemma-invert-s-sections}
\begin{reference}
\cite[Section 6.8]{EGA1-second}
\end{reference}
$X$를 스킴이라 하자. $\mathcal{L}$을 $X$ 위의 가역층이라 하자.
$s \in \Gamma(X, \mathcal{L})$라 하자. $\mathcal{F}$를 준연접
$\mathcal{O}_X$-가군이라 하자.
\begin{enumerate}
\item $X$가 준콤팩트이면 (\ref{properties-equation-module-invert-s})는
단사이고,
\item $X$가 준콤팩트이고 준분리이면
(\ref{properties-equation-module-invert-s})는 동형사상이다.
\end{enumerate}
특히 표준 사상
$$
\Gamma_*(X, \mathcal{L})_{(s)}
\longrightarrow
\Gamma(X_s, \mathcal{O}_X),\quad
a/s^n \longmapsto a \otimes s^{-n}
$$
은 $X$가 준콤팩트이고 준분리이면 동형사상이다.
\end{lemma}

\begin{proof}
$X$가 준콤팩트라고 가정하자. 유한 아핀 열린 덮개
$X = U_1 \cup \ldots \cup U_m$을 선택하자. 여기서 $U_j$는
아핀이고 $\mathcal{L}|_{U_j} \cong \mathcal{O}_{U_j}$이다.
이 동형사상을 통해 상 $s|_{U_j}$는 어떤
$f_j \in \Gamma(U_j, \mathcal{O}_{U_j})$에 대응한다. 그러면
$X_s \cap U_j = D(f_j)$이다.

\medskip\noindent
(1)의 증명. (\ref{properties-equation-module-invert-s})의 핵의 원소
$t/s^n$을 잡자. 그러면 $t|_{X_s} = 0$이다.
따라서 $(t|_{U_j})|_{D(f_j)} = 0$이다. 보조정리
\ref{properties-lemma-invert-f-sections}에 의해 $f_j^{e_j} t|_{U_j} = 0$이
되는 어떤 $e_j \geq 0$이 존재한다. $e = \max(e_j)$라 하자.
그러면 $t \otimes s^e$는 $U_j$에 제한하면 모든 $j$에 대하여
영이므로 그 자체가 영이다. $t/s^n$은
$t \otimes s^e/s^{n + e}$와
$\Gamma_*(X, \mathcal{L}, \mathcal{F})_{(s)}$ 안에서 같으므로
원하는 대로 $t/s^n = 0$을 얻는다.

\medskip\noindent
(2)의 증명. $X$가 준콤팩트이고 준분리라고 가정하자.
그러면 $U_j \cap U_{j'}$는 모든 쌍 $j, j'$에 대하여 준콤팩트이다.
Schemes의 보조정리 \ref{schemes-lemma-characterize-quasi-separated}를
보라. (1)에 의해 (\ref{properties-equation-module-invert-s})가 단사임을 안다.
$t' \in \Gamma(X_s, \mathcal{F}|_{X_s})$라 하자. 모든 $j$에 대하여
정수 $e_j \geq 0$과 $t'_j \in \Gamma(U_j, \mathcal{F}|_{U_j})$가
존재하여 $t'|_{D(f_j)}$는 $t'_j/f_j^{e_j}$에 대응한다.
여기서는 보조정리 \ref{properties-lemma-invert-f-sections}의 동형사상을
사용하였다. $e = \max(e_j)$라 놓고
$$
t_j = f_j^{e - e_j} t'_j \otimes q_j^e \in
\Gamma(U_j,
(\mathcal{F} \otimes_{\mathcal{O}_X} \mathcal{L}^{\otimes e})|_{U_j})
$$
라 하자. 여기서 $q_j \in \Gamma(U_j, \mathcal{L}|_{U_j})$는
동형사상 $\mathcal{L}|_{U_j} \cong \mathcal{O}_{U_j}$에서 오는
자명화 절단이다. 특히 $s|_{U_j} = f_j q_j$이다. 이를 사용하여
계산하면 $t_j|_{U_j \cap U_{j'}}$와
$t_{j'}|_{U_j \cap U_{j'}}$는 $\mathcal{F}$의 같은 절단으로
사상되고, 그 절단은 $U_j \cap U_{j'} \cap X_s$ 위에 정의된다.
$U_j \cap U_{j'}$의 준콤팩트성과 (1)에 의해 어떤 정수
$e' \geq 0$이 존재하여
$$
t_j|_{U_j \cap U_{j'}} \otimes s^{e'}|_{U_j \cap U_{j'}} =
t_{j'}|_{U_j \cap U_{j'}} \otimes s^{e'}|_{U_j \cap U_{j'}}
$$
가 $\mathcal{F} \otimes \mathcal{L}^{\otimes e + e'}$의 절단으로서
$U_j \cap U_{j'}$ 위에서 성립한다. 모든 쌍에 대하여 같은
$e'$을 선택하여 쓸 수 있다. 여기서 쌍은 $j, j'$이다. 그러면 층
조건에 의해 절단
$t \in \Gamma(X, \mathcal{F} \otimes \mathcal{L}^{\otimes e + e'})$
가 존재하고, 그 $U_j$로의 제한은
$t_j \otimes s^{e'}|_{U_j}$이다. 간단히 계산하면
$t/s^{e + e'}$가 원하는 대로 $t'$으로 사상됨을 알 수 있다.
\end{proof}

\noindent
$X$를 스킴이라 하자. $\mathcal{L}$을 가역
$\mathcal{O}_X$-가군이라 하자. $\mathcal{F}$와 $\mathcal{G}$를
준연접 $\mathcal{O}_X$-가군이라 하자. 등급
$\Gamma_*(X, \mathcal{L})$-가군
$$
M = \bigoplus\nolimits_{n \in \mathbf{Z}} \Hom_{\mathcal{O}_X}(\mathcal{F},
\mathcal{G} \otimes_{\mathcal{O}_X} \mathcal{L}^{\otimes n})
$$
을 생각하자. 다음으로 절단 $s \in \Gamma(X, \mathcal{L})$를
잡자. 그러면 표준 사상
\begin{equation}
\label{properties-equation-hom-invert-s}
M_{(s)} \longrightarrow
\Hom_{\mathcal{O}_{X_s}}(\mathcal{F}|_{X_s}, \mathcal{G}|_{X_s})
\end{equation}
이 있다. 이 사상은 $\alpha/s^n$을
$\alpha|_{X_s} \otimes s|_{X_s}^{-n}$으로 보낸다.
다음 보조정리를 보조정리 \ref{properties-lemma-extend-finite-presentation}과
함께 사용하면, 대략적으로 $X$가 준콤팩트이고 준분리일 때 유한 표시
$\mathcal{O}_{X_s}$-가군의 범주는 유한 표시
$\mathcal{O}_X$-가군의 범주에서 사상들의 곱셈적 계
$s^n: \mathcal{F} \to
\mathcal{F} \otimes_{\mathcal{O}_X} \mathcal{L}^{\otimes n}$을
가역화하여 얻는 범주라는 것을 알 수 있다.

\begin{lemma}
\label{properties-lemma-section-maps-backwards}
$X$를 스킴이라 하자. $\mathcal{L}$을 가역
$\mathcal{O}_X$-가군이라 하자.
$s \in \Gamma(X, \mathcal{L})$를 절단이라 하자.
$\mathcal{F}$, $\mathcal{G}$를 준연접 $\mathcal{O}_X$-가군이라 하자.
\begin{enumerate}
\item $X$가 준콤팩트이고 $\mathcal{F}$가 유한형이면
(\ref{properties-equation-hom-invert-s})는 단사이고,
\item $X$가 준콤팩트이고 준분리이며 $\mathcal{F}$가 유한 표시이면
(\ref{properties-equation-hom-invert-s})는 전단사이다.
\end{enumerate}
\end{lemma}

\begin{proof}
먼저 $X = \Spec(A)$가 아핀이고
$\mathcal{L} = \mathcal{O}_X$인 경우를 증명하자. 이 경우 $s$는
원소 $f \in A$에 대응한다. $\mathcal{F} = \widetilde{M}$ 및
$\mathcal{G} = \widetilde{N}$이라고 하자. 여기서 $A$-가군은
각각 $M$과 $N$이다. 그러면 보조정리 \ref{properties-lemma-finite-type-module}와
\ref{properties-lemma-finite-presentation-module}를 통해 이 보조정리는
다음 대수 명제들로 옮겨진다.
\begin{enumerate}
\item $M$이 유한 $A$-가군이고 $\varphi : M \to N$이
$A$-가군 사상이며 유도된 사상 $M_f \to N_f$가 영이면,
$f^n\varphi = 0$이 되는 어떤 $n$이 존재한다.
\item $M$이 유한 표시 $A$-가군이면
$\Hom_A(M, N)_f = \Hom_{A_f}(M_f, N_f)$이다.
\end{enumerate}
두 번째 명제는 Algebra의 보조정리
\ref{algebra-lemma-hom-from-finitely-presented}이고, 첫 번째
명제의 증명은 생략한다.

\medskip\noindent
이제 일반적인 $X$에 대하여 (1)을 증명하자.
$X$가 준콤팩트라고 가정하고 유한 아핀 열린 덮개
$X = U_1 \cup \ldots \cup U_m$을 선택하자. 여기서 $U_j$는
아핀이고 $\mathcal{L}|_{U_j} \cong \mathcal{O}_{U_j}$이다.
이 동형사상을 통해 상 $s|_{U_j}$는 어떤
$f_j \in \Gamma(U_j, \mathcal{O}_{U_j})$에 대응한다. 그러면
$X_s \cap U_j = D(f_j)$이다.
(\ref{properties-equation-hom-invert-s})의 핵의 원소 $\alpha/s^n$을 잡자.
그러면 $\alpha|_{X_s} = 0$이다.
따라서 $(\alpha|_{U_j})|_{D(f_j)} = 0$이다. 위에서 다룬 아핀
경우에 의해 $f_j^{e_j} \alpha|_{U_j} = 0$이 되는 어떤
$e_j \geq 0$이 존재한다. $e = \max(e_j)$라 하자. 그러면
$\alpha \otimes s^e$는 $U_j$에 제한하면 모든 $j$에 대하여
영이므로 그 자체가 영이다. $\alpha/s^n$은
$\alpha \otimes s^e/s^{n + e}$와 $M_{(s)}$ 안에서 같으므로
원하는 대로 $\alpha/s^n = 0$이다.

\medskip\noindent
(2)의 증명. $\mathcal{F}$가 유한 표시이므로 층
$\SheafHom_{\mathcal{O}_X}(\mathcal{F}, \mathcal{G})$는 준연접이다.
Schemes의 절 \ref{schemes-section-quasi-coherent}를 보라.
또한
$$
\SheafHom_{\mathcal{O}_X}(\mathcal{F},
\mathcal{G} \otimes_{\mathcal{O}_X} \mathcal{L}^{\otimes n}) =
\SheafHom_{\mathcal{O}_X}(\mathcal{F}, \mathcal{G})
\otimes_{\mathcal{O}_X} \mathcal{L}^{\otimes n}
$$
임은 모든 $n$에 대하여 명백하다. 따라서 이 경우에는 보조정리
\ref{properties-lemma-invert-s-sections}를
$\SheafHom_{\mathcal{O}_X}(\mathcal{F}, \mathcal{G})$에 적용하여
명제가 따른다.
\end{proof}

\section{준아핀 스킴}
\label{properties-section-quasi-affine}

\begin{definition}
\label{properties-definition-quasi-affine}
스킴 $X$가 준콤팩트이고 어떤 아핀 스킴의 열린 부분스킴과
동형이면 이를 {\it 준아핀}이라고 한다.
\end{definition}

\begin{lemma}
\label{properties-lemma-quasi-coherent-quasi-affine}
$A$를 환이라 하고 $U \subset \Spec(A)$를 준콤팩트 열린
부분스킴이라 하자. 준연접층 $\mathcal{F}$가 $U$ 위에 주어지면 표준 사상
$$
\widetilde{H^0(U, \mathcal{F})}|_U \to \mathcal{F}
$$
은 동형사상이다.
\end{lemma}

\begin{proof}
$j : U \to \Spec(A)$를 포함사상이라 쓰자. 그러면
$H^0(U, \mathcal{F}) = H^0(\Spec(A), j_*\mathcal{F})$이고,
$j_*\mathcal{F}$는 Schemes의 보조정리
\ref{schemes-lemma-push-forward-quasi-coherent}에 의해 준연접이다.
따라서 Schemes의 보조정리
\ref{schemes-lemma-equivalence-quasi-coherent}에 의해
$j_*\mathcal{F} = \widetilde{H^0(U, \mathcal{F})}$이다.
다시 $U$에 제한하면 보조정리를 얻는다.
\end{proof}

\begin{lemma}
\label{properties-lemma-invert-f-affine}
$X$를 스킴이라 하자. $f \in \Gamma(X, \mathcal{O}_X)$라 하자.
$X$가 준콤팩트이고 준분리이며 $X_f$가 아핀이라고 가정하자.
그러면 Schemes의 보조정리
\ref{schemes-lemma-morphism-into-affine}에서 오는 표준 사상
$$
j : X \longrightarrow \Spec(\Gamma(X, \mathcal{O}_X))
$$
은 $X_f = j^{-1}(D(f))$에서 표준 아핀 열린집합
$D(f) \subset \Spec(\Gamma(X, \mathcal{O}_X))$ 위로의 동형사상을
유도한다.
\end{lemma}

\begin{proof}
$j$가 보조정리 \ref{properties-lemma-invert-f-sections}에 의해 환의 동형사상
$\Gamma(X, \mathcal{O}_X)_f \to \mathcal{O}_X(X_f)$을 유도하므로
이는 명백하다.
\end{proof}

\begin{lemma}
\label{properties-lemma-quasi-affine}
$X$를 스킴이라 하자. 그러면 $X$가 준아핀일 필요충분조건은
Schemes의 보조정리 \ref{schemes-lemma-morphism-into-affine}에서 오는
표준 사상
$$
X \longrightarrow \Spec(\Gamma(X, \mathcal{O}_X))
$$
이 준콤팩트 열린 몰입인 것이다.
\end{lemma}

\begin{proof}
표시된 사상이 준콤팩트 열린 몰입이면 $X$는
$\Spec(\Gamma(X, \mathcal{O}_X))$의 준콤팩트 열린 부분스킴과
동형이고, 따라서 $X$는 분명 준아핀이다.

\medskip\noindent
$X$가 준아핀이라고 가정하자. 즉, $X \subset \Spec(R)$이
준콤팩트 열린집합이라고 하자. 특히 이로부터 $X$가 분리임이 따른다.
Schemes의 보조정리
\ref{schemes-lemma-subscheme-of-separated-scheme}를 보라.
$A = \Gamma(X, \mathcal{O}_X)$라 하자.
환 사상 $R \to A$를 생각하자. 이는
$R = \Gamma(\Spec(R), \mathcal{O}_{\Spec(R)})$과 층
$\mathcal{O}_{\Spec(R)}$의 제한 사상에서 온다.
Schemes의 보조정리 \ref{schemes-lemma-morphism-into-affine}에 의해
포함사상의 분해
$$
X \longrightarrow
\Spec(A) \longrightarrow
\Spec(R)
$$
를 얻는다. $x \in X$라 하자. $r \in R$을 선택하여
$x \in D(r)$이고 $D(r) \subset X$가 되게 하자.
$f \in A$를 $r$의 $A$에서의 상이라 쓰자. 위의 보조정리
\ref{properties-lemma-invert-f-sections}의 열린집합 $X_f$는
$D(r) \subset X$와 같으므로, 그 보조정리의 결론에 의해
$A_f \cong R_r$이다.
따라서 $D(r) \to \Spec(A)$는 표준 아핀 열린집합 $D(f)$,
즉 $\Spec(A)$ 안의 그 열린집합 위로의 동형사상이다. $X$는 이러한 아핀 열린집합
$D(f)$들로 덮이므로 결론이 따른다.
\end{proof}

\begin{lemma}
\label{properties-lemma-cartesian-diagram-quasi-affine}
$U \to V$를 준아핀 스킴들의 열린 몰입이라 하자. 그러면
$$
\xymatrix{
U \ar[d] \ar[rr]_-j & & \Spec(\Gamma(U, \mathcal{O}_U)) \ar[d] \\
U \ar[r] & V \ar[r]^-{j'} & \Spec(\Gamma(V, \mathcal{O}_V))
}
$$
은 데카르트 도식이다.
\end{lemma}

\begin{proof}
Schemes의 보조정리 \ref{schemes-lemma-morphism-into-affine}에 의해
도식은 가환한다.
$A = \Gamma(U, \mathcal{O}_U)$ 및 $B = \Gamma(V, \mathcal{O}_V)$라
쓰자. $g \in B$를 잡아 $V_g$가 아핀이고 $U$에 포함되게 하자.
이는 $f$가 $g$의 $A$에서의 상일 때 $U_f = V_g$라는
뜻이다. 보조정리 \ref{properties-lemma-invert-f-affine}에 의해 $j'$은
$V_g$에서 표준 열린집합 $D(g)$, 즉 $\Spec(B)$ 안의 그 열린집합
위로의 동형사상을
유도한다. 따라서
$V_g \times_{\Spec(B)} \Spec(A) \to \Spec(A)$는
$D(f) \subset \Spec(A)$ 위로의 동형사상이다. 보조정리
\ref{properties-lemma-invert-f-affine}를 다시 쓰면 $j$는 $U_f$를
$D(f)$ 위로 동형으로 보낸다. 따라서
$U_f = U_f \times_{\Spec(B)} \Spec(A)$이다. 보조정리
\ref{properties-lemma-quasi-affine}에 의해 $U$는 위와 같은 $V_g = U_f$들로
덮이므로 $U \to U \times_{\Spec(B)} \Spec(A)$는 동형사상이다.
\end{proof}

\begin{lemma}
\label{properties-lemma-quasi-affine-presentation}
$X$를 준아핀 스킴이라 하자. 정수 $n \geq 0$, 아핀 스킴 $T$,
그리고 사상 $T \to X$가 존재하여, 모든 사상 $X' \to X$ 중
$X'$이 아핀인 것에 대해 섬유곱 $X' \times_X T$는
$\mathbf{A}^n_{X'}$와 $X'$ 위에서 동형이다.
\end{lemma}

\begin{proof}
정의에 의해 어떤 환 $A$가 존재하여 $X$는 준콤팩트 열린 부분스킴
$U \subset \Spec(A)$와 동형이다. 표준 열린집합
$D(f) \subset \Spec(A)$들은 위상의 기저를 이룬다는 것을 상기하자.
Algebra의 절 \ref{algebra-section-spectrum-ring}을 보라.
$U$가 준콤팩트이므로
$f_1, \ldots, f_n \in A$를
$U = D(f_1) \cup \ldots \cup D(f_n)$이 되도록 선택할 수 있다.
따라서 $X = \Spec(A) \setminus V(I)$이고
$I = (f_1, \ldots, f_n)$이라고 가정해도 된다. 다음과 같이 놓자.
$$
T = \Spec(A[t, x_1, \ldots, x_n]/(f_1 x_1 + \ldots + f_n x_n - 1))
$$
구조사상 $T \to \Spec(A)$는 열린집합 $X$를 통해 분해되어
사상 $T \to X$를 준다. $X' = \Spec(A')$이고 사상 $X' \to X$가
환 사상 $A \to A'$에 대응한다면, 상들
$f'_1, \ldots, f'_n \in A'$은 $f_1, \ldots, f_n$의 상이며
$A'$에서 단위 아이디얼을 생성한다.
$1 = f'_1 a'_1 + \ldots + f'_n a'_n$이라고 하자.
밑변환 $X' \times_X T$는
$A'[t, x_1, \ldots, x_n]/(f'_1 x_1 + \ldots + f'_n x_n - 1)$의
스펙트럼이다. $A'$-대수 준동형
$$
\varphi :
A'[y_1, \ldots, y_n]
\longrightarrow
A'[t, x_1, \ldots, x_n, x_{n + 1}]/(f'_1 x_1 + \ldots + f'_n x_n - 1)
$$
이 $y_i$를 $a'_i t + x_i$로 보낼 때 동형사상이라고 주장한다.
이 주장은 보조정리의 증명을 끝낸다. $\varphi$의 역은
$A'$-대수 준동형
$$
\psi :
A'[t, x_1, \ldots, x_n, x_{n + 1}]/(f'_1 x_1 + \ldots + f'_n x_n - 1)
\longrightarrow
A'[y_1, \ldots, y_n]
$$
으로 주어지는데, 이는 $t$를
$-1 + f'_1 y_1 + \ldots + f'_n y_n$으로 보내고 $x_i$를
$y_i + a'_i - a'_i(f'_1 y_1 + \ldots + f'_n y_n)$으로 보낸다.
여기서 $i = 1, \ldots, n$이다. 이는 $\sum f'_ix_i$가
$$
\begin{matrix}
\sum f'_i(y_i + a'_i - a'_i(\sum f'_j y_j)) =
(\sum f'_iy_i) + 1 - (\sum f'_j y_j) = 1
\end{matrix}
$$
로 사상되므로 잘 정의된다. 두 사상이 서로 역임을 보려면
다음과 같이 계산하면 된다.
$$
\begin{matrix}
\varphi(\psi(t) = \varphi(-1 +  \sum f'_i y_i) =
-1 + \sum f'_i (a'_i t + x_i) = t \\
\varphi(\psi(x_i)) = \varphi(y_i + a'_i - a'_i(\sum f'_j y_j)) =
a'_i t + x_i + a'_i - a'_i(\sum f'_ja'_jt + f'_jx_j) = x_i \\
\psi(\varphi(y_i)) = \psi(a'_i t + x_i) =
a'_i(-1 + \sum f'_j y_j) + y_i + a'_i - a'_i(\sum f'_j y_j) = y_i
\end{matrix}
$$
이로써 증명이 끝난다.
\end{proof}

\section{평탄 가군}
\label{properties-section-flat}

\noindent
어떤 환 달린 공간 $(X, \mathcal{O}_X)$에서든
$\mathcal{O}_X$-가군이 (한 점에서) 평탄하다는 말의 뜻을 알고 있다.
Modules의 정의 \ref{modules-definition-flat}
(정의 \ref{modules-definition-flat-at-point})를 보라.
아핀 스킴 위의 준연접층에 대해서는 이것이 대수 장에서 정의한
개념과 일치한다.

\begin{lemma}
\label{properties-lemma-flat-module}
\begin{slogan}
평탄성은 가군과 층에서 같은 개념이다.
\end{slogan}
$X = \Spec(R)$을 아핀 스킴이라 하자.
$\mathcal{F} = \widetilde{M}$이라 하자. 여기서 $R$-가군은 $M$이다.
준연접층 $\mathcal{F}$가
평탄 $\mathcal{O}_X$-가군일 필요충분조건은 $M$이 평탄
$R$-가군인 것이다.
\end{lemma}

\begin{proof}
$\mathcal{F}$의 평탄성은 줄기에서 검사할 수 있다. Modules의
보조정리 \ref{modules-lemma-flat-stalks-flat}를 보라.
환 위의 가군에 대해서도 마찬가지이다. Algebra의 보조정리
\ref{algebra-lemma-flat-localization}를 보라. 그리고
$\mathcal{F}_x = M_{\mathfrak p}$이며 $x$가 $\mathfrak p$에
대응하므로 보조정리가 성립한다.
\end{proof}





\section{국소 자유 가군}
\label{properties-section-finite-locally-free}

\noindent
어떤 환 달린 공간에서든 $\mathcal{O}_X$-가군이 (유한) 국소
자유라는 말의 뜻을 알고 있다. 아핀 스킴에서는 이것이 대수 장에서
정의한 개념과 일치한다.

\begin{lemma}
\label{properties-lemma-locally-free-module}
$X = \Spec(R)$을 아핀 스킴이라 하자.
$\mathcal{F} = \widetilde{M}$이라 하자. 여기서 $R$-가군은 $M$이다.
준연접층 $\mathcal{F}$가
(유한) 국소 자유 $\mathcal{O}_X$-가군일 필요충분조건은 $M$이
(유한) 국소 자유 $R$-가군인 것이다.
\end{lemma}

\begin{proof}
정의에서 바로 따른다. Modules의 정의
\ref{modules-definition-locally-free}와 Algebra의 정의
\ref{algebra-definition-locally-free}를 보라.
\end{proof}

\noindent
유한 국소 자유 가군은 여러 가지 서로 다른 방식으로 특징지을 수 있다.

\begin{lemma}
\label{properties-lemma-finite-locally-free}
$X$를 스킴이라 하자. $\mathcal{F}$를 준연접
$\mathcal{O}_X$-가군이라 하자. 다음 조건들은 서로 동치이다.
\begin{enumerate}
\item $\mathcal{F}$는 유한 표시인 평탄 $\mathcal{O}_X$-가군이다.
\item $\mathcal{F}$는 유한 표시 $\mathcal{O}_X$-가군이고,
모든 $x \in X$에 대하여 줄기 $\mathcal{F}_x$는 자유
$\mathcal{O}_{X, x}$-가군이다.
\item $\mathcal{F}$는 국소 자유 유한형 $\mathcal{O}_X$-가군이다.
\item $\mathcal{F}$는 유한 국소 자유 $\mathcal{O}_X$-가군이다.
\item $\mathcal{F}$는 유한형 $\mathcal{O}_X$-가군이고,
모든 $x \in X$에 대하여 줄기 $\mathcal{F}_x$는 자유
$\mathcal{O}_{X, x}$-가군이며, 함수
$$
\rho_\mathcal{F} : X \to \mathbf{Z}, \quad
x \longmapsto
\dim_{\kappa(x)} \mathcal{F}_x \otimes_{\mathcal{O}_{X, x}} \kappa(x)
$$
는 $X$의 자리스키 위상에 관하여 국소 상수이다.
\end{enumerate}
\end{lemma}

\begin{proof}
이 보조정리는 곧바로 아핀인 경우로 환원된다. 그 경우에는
Algebra의 보조정리 \ref{algebra-lemma-finite-projective}를 다시
서술한 것이다. 이 번역에는 보조정리
\ref{properties-lemma-finite-type-module},
\ref{properties-lemma-finite-presentation-module},
\ref{properties-lemma-flat-module}, 그리고
\ref{properties-lemma-locally-free-module}를 사용한다.
\end{proof}

\begin{lemma}
\label{properties-lemma-finite-locally-free-reduced}
$X$를 축약 스킴이라 하자. $\mathcal{F}$를 준연접
$\mathcal{O}_X$-가군이라 하자. 그러면 보조정리
\ref{properties-lemma-finite-locally-free}의 서로 동치인 조건들은 다음 조건과도
동치이다.
\begin{enumerate}
\item[(6)] $\mathcal{F}$는 유한형 $\mathcal{O}_X$-가군이고 함수
$$
\rho_\mathcal{F} : X \to \mathbf{Z}, \quad
x \longmapsto
\dim_{\kappa(x)} \mathcal{F}_x \otimes_{\mathcal{O}_{X, x}} \kappa(x)
$$
는 $X$의 자리스키 위상에 관하여 국소 상수이다.
\end{enumerate}
\end{lemma}

\begin{proof}
이 보조정리는 곧바로 아핀인 경우로 환원된다. 그 경우에는
Algebra의 보조정리
\ref{algebra-lemma-finite-projective-reduced}를 다시 서술한 것이다.
\end{proof}




\section{국소 사영 가군}
\label{properties-section-locally-projective}

\noindent
대수 장에서 수행한 작업의 한 귀결로서, 다음과 같이 국소 사영
가군을 정의하는 것이 타당하다.

\begin{definition}
\label{properties-definition-locally-projective}
$X$를 스킴이라 하자. $\mathcal{F}$를 준연접
$\mathcal{O}_X$-가군이라 하자. $\mathcal{F}$가 {\it 국소 사영}이라는 것은
모든 아핀 열린집합
$U \subset X$에 대하여 $\mathcal{O}_X(U)$-가군
$\mathcal{F}(U)$가 사영이라는 뜻이다.
\end{definition}

\begin{lemma}
\label{properties-lemma-locally-projective}
$X$를 스킴이라 하자. $\mathcal{F}$를 준연접
$\mathcal{O}_X$-가군이라 하자. 다음 조건들은 서로 동치이다.
\begin{enumerate}
\item $\mathcal{F}$는 국소 사영이다.
\item 아핀 열린 덮개 $X = \bigcup U_i$가 존재하여
$\mathcal{O}_X(U_i)$-가군 $\mathcal{F}(U_i)$는 모든 $i$에
대하여 사영이다.
\end{enumerate}
특히 $X = \Spec(A)$이고 $\mathcal{F} = \widetilde{M}$이면
$\mathcal{F}$가 국소 사영일 필요충분조건은 $M$이 사영
$A$-가군인 것이다.
\end{lemma}

\begin{proof}
먼저 $M$이 사영 $A$-가군이고 $A \to B$가 환 사상이면
$M \otimes_A B$는 사영 $B$-가군이다. Algebra의 보조정리
\ref{algebra-lemma-ascend-properties-modules}를 보라.
따라서 $U$가 아핀 열린집합이고 $\mathcal{F}(U)$가 사영
$\mathcal{O}_X(U)$-가군이면, 표준 열린집합 $D(f)$도 아핀이고
$\mathcal{F}(D(f))$는 사영 $\mathcal{O}_X(D(f))$-가군이다.
이는 모든 $f \in \mathcal{O}_X(U)$에 대하여 성립한다.
(2)가 성립한다고 가정하자. $U \subset X$를 임의의 아핀
열린집합이라 하자. 열린 덮개
$U = \bigcup_{j = 1, \ldots, m} D(f_j)$를 찾을 수 있고, 이는
유한 개의 표준 열린집합 $D(f_j)$로 이루어지며,
각 $j$에 대하여 열린집합 $D(f_j)$는 어떤 $U_i$의 표준
열린집합이다. Schemes의 보조정리
\ref{schemes-lemma-standard-open-two-affines}를 보라.
$A = \mathcal{O}_X(U)$라 놓고, $M$을 $A$-가군이라 하되
$\mathcal{F}|_U$가 $M$에 대응한다고 하자. 그러면 $M_{f_j}$는 사영
$A_{f_j}$-가군이다. 따라서 $A \to B = \prod A_{f_j}$는
충실히 평탄인 환 사상이고 $M \otimes_A B$는 사영 $B$-가군이다.
그러므로 Algebra의 정리
\ref{algebra-theorem-ffdescent-projectivity}에 의해 $M$은 사영이다.
\end{proof}

\begin{lemma}
\label{properties-lemma-locally-projective-pullback}
$f : X \to Y$를 스킴의 사상이라 하자. $\mathcal{G}$를 준연접
$\mathcal{O}_Y$-가군이라 하자. $\mathcal{G}$가 $Y$ 위에서 국소
사영이면 $f^*\mathcal{G}$는 $X$ 위에서 국소 사영이다.
\end{lemma}

\begin{proof}
Algebra의 보조정리 \ref{algebra-lemma-ascend-properties-modules}와
보조정리 \ref{properties-lemma-locally-projective}에서 따른다.
\end{proof}

\section{준연접층의 연장}
\label{properties-section-extending-quasi-coherent-sheaves}

\noindent
열린 부분스킴 위에 주어진 준연접층이 전체 스킴으로 연장됨을
보일 수 있으면 때때로 유용하다.

\begin{lemma}
\label{properties-lemma-extend-trivial}
$j : U \to X$를 스킴들의 준콤팩트 열린 몰입이라 하자.
\begin{enumerate}
\item $U$ 위의 임의의 준연접층은 $X$ 위의 준연접층으로 연장된다.
\item $\mathcal{F}$를 $X$ 위의 준연접층이라 하자.
$\mathcal{G} \subset \mathcal{F}|_U$를 준연접 부분층이라 하자.
$\mathcal{H}$가 준연접 부분층으로서 존재하며, 이는
$\mathcal{F}$의 부분층이고
$\mathcal{H}|_U = \mathcal{G}$이고, 이는 $\mathcal{F}|_U$의
부분층으로서의 등식이다.
\item $\mathcal{F}$를 $X$ 위의 준연접층이라 하자.
$\mathcal{G}$를 $U$ 위의 준연접층이라 하자.
$\varphi : \mathcal{G} \to \mathcal{F}|_U$를
$\mathcal{O}_U$-가군의 사상이라 하자. 준연접층 $\mathcal{H}$가
$\mathcal{O}_X$-가군의 층으로서 존재하고, 사상
$\psi : \mathcal{H} \to \mathcal{F}$가 존재하여
$\mathcal{H}|_U = \mathcal{G}$이고 $\psi|_U = \varphi$이다.
\end{enumerate}
\end{lemma}

\begin{proof}
몰입은 분리이다(Schemes의 보조정리
\ref{schemes-lemma-immersions-monomorphisms}를 보라). 또한
가정에 의해 $j$는 준콤팩트이다. 따라서 임의의 준연접층
$\mathcal{G}$가 $U$ 위에 주어지면 층 $j_*\mathcal{G}$는 $X$로의 연장이다.
Schemes의 보조정리 \ref{schemes-lemma-push-forward-quasi-coherent}와
Sheaves의 절 \ref{sheaves-section-open-immersions}을 보라.

\medskip\noindent
$\mathcal{F}$, $\mathcal{G}$가 (2)와 같다고 가정하자. 그러면
$j_*\mathcal{G}$는 $X$ 위의 준연접층이다(위를 보라).
이는 $j_*j^*\mathcal{F}$의 부분층이다. 따라서 핵
$$
\mathcal{H} =
\Ker(\mathcal{F} \oplus j_* \mathcal{G}
\longrightarrow j_*j^*\mathcal{F})
$$
도 준연접이다. Schemes의 절
\ref{schemes-section-quasi-coherent}를 보라.
$\mathcal{H} \subset \mathcal{F}$이고
$\mathcal{H}|_U = \mathcal{G}$임은 형식적으로 확인할 수 있다
(Sheaves의 절 \ref{sheaves-section-open-immersions}의 내용을
다시 사용한다).

\medskip\noindent
(3)도 (2)와 같은 방법으로 증명한다. $\mathcal{H} = \Ker(\mathcal{F} \oplus j_* \mathcal{G}
\to j_*j^*\mathcal{F})$를 취하고, $\mathcal{F}$로 가는 명백한
사상 및 $\mathcal{G}$와의 $U$ 위에서의 명백한 식별을 사용하면 된다.
\end{proof}

\begin{lemma}
\label{properties-lemma-extend}
$X$를 준콤팩트 준분리 스킴이라 하자.
$U \subset X$를 준콤팩트 열린집합이라 하자.
$\mathcal{F}$를 준연접 $\mathcal{O}_X$-가군이라 하자.
$\mathcal{G} \subset \mathcal{F}|_U$를 유한형 준연접
$\mathcal{O}_U$-부분가군이라 하자. 그러면 유한형 준연접 부분가군
$\mathcal{G}' \subset \mathcal{F}$가 존재하여
$\mathcal{G}'|_U = \mathcal{G}$이다.
\end{lemma}

\begin{proof}
$n$을 아핀 열린집합 $U_i \subset X$,
$i = 1, \ldots , n$의 최소 개수라 하자. 이들은
$X = U \cup \bigcup U_i$가 되게 한다.
(여기서 $X$의 준콤팩트성을 사용한다.) $n = 1$인 경우에
보조정리를 증명할 수 있다고 하자. 그러면 $\mathcal{G}$를 차례로
$\mathcal{G}_1$로 $U \cup U_1$ 위에,
$\mathcal{G}_2$로 $U \cup U_1 \cup U_2$ 위에,
$\mathcal{G}_3$로 $U \cup U_1 \cup U_2 \cup U_3$ 위에
연장하는 식으로 계속할 수 있다. 따라서 $n = 1$인 경우로 환원된다.

\medskip\noindent
그러므로 $X = U \cup V$이고 $V$가 아핀이라고 가정해도 된다.
$X$가 준분리이고 $U$, $V$가 준콤팩트 열린집합이므로
$U \cap V$는 준콤팩트 열린집합이다. 계
$(V, U \cap V, \mathcal{F}|_V, \mathcal{G}|_{U \cap V})$에 대해
보조정리를 증명하면 충분하다. 실제로 그 결과 얻는 층
$\mathcal{G}'$를 $V$ 위에서 주어진 층 $\mathcal{G}$와
$U$ 위에서 $U \cap V$ 위의 공통값을 따라 붙일 수 있다.
따라서 $X$가 아핀인 경우로 환원된다.

\medskip\noindent
$X = \Spec(R)$이라고 가정하자. $\mathcal{F} = \widetilde M$이라
쓰자. 여기서 $R$-가군은 $M$이다. 위의 보조정리
\ref{properties-lemma-extend-trivial}에 의해 준연접 부분층
$\mathcal{H} \subset \mathcal{F}$를 찾을 수 있고, 이는
$\mathcal{G}$로 $U$ 위에서 제한된다.
$\mathcal{H} = \widetilde N$이라 쓰자. 여기서 $R$-가군은 $N$이다.
모든 $u \in U$에 대하여
어떤 $f \in R$이 존재하여 $u \in D(f) \subset U$이고
$N_f$는 유한 생성이다. 보조정리
\ref{properties-lemma-finite-type-module}를 보라.
$U$가 준콤팩트이므로 유한 개의 $D(f_i)$로 덮을 수 있고,
각 $N_{f_i}$는 유한 개의 원소, 이를테면
$x_{i, 1}/f_i^N, \ldots, x_{i, r_i}/f_i^N$로 생성된다.
$N' \subset N$을 원소 $x_{i, j}$들로 생성되는 부분가군이라 하자.
그러면 부분층
$\mathcal{G}' = \widetilde{N'} \subset \mathcal{H} \subset \mathcal{F}$
가 원하는 것이다.
\end{proof}

\begin{lemma}
\label{properties-lemma-quasi-coherent-colimit-finite-type}
$X$를 준콤팩트 준분리 스킴이라 하자. 임의의
$\mathcal{O}_X$-가군의 준연접층은 그 유한형 준연접
$\mathcal{O}_X$-부분가군들의 유향 여극한이다.
\end{lemma}

\begin{proof}
$\mathcal{G}_1$, $\mathcal{G}_2$가 유한형 준연접 부분층이면
$\mathcal{G}_1 \oplus \mathcal{G}_2 \to \mathcal{F}$의 상도
유한형 준연접 부분가군이므로 이 여극한은 유향이다.
$U \subset X$를 임의의 아핀 열린집합이라 하고,
$s \in \Gamma(U, \mathcal{F})$를 임의의 절단이라 하자.
$\mathcal{G} \subset \mathcal{F}|_U$를 $s$로 생성되는 부분층이라
하자. 그러면 분명히 $\mathcal{G}$는 준연접이고
$\mathcal{O}_U$-가군으로서 유한형이다. 보조정리
\ref{properties-lemma-extend}에 의해 $\mathcal{G}$는 유한형 준연접 부분층
$\mathcal{G}' \subset \mathcal{F}$의 제한이다.
$X$의 위상에는 아핀 열린집합들로 이루어진 기저가 있으므로,
$\mathcal{F}$의 모든 국소 절단은 국소적으로 어떤 유한형 준연접
부분가군에 포함된다. 이로써 결론이 따른다.
\end{proof}

\begin{lemma}
\label{properties-lemma-extend-finite-presentation}
$X$를 준콤팩트 준분리 스킴이라 하자. $\mathcal{F}$를 준연접
$\mathcal{O}_X$-가군이라 하자. $U \subset X$를 준콤팩트
열린집합이라 하자. $\mathcal{G}$를 유한 표시
$\mathcal{O}_U$-가군이라 하자.
$\varphi : \mathcal{G} \to \mathcal{F}|_U$를
$\mathcal{O}_U$-가군의 사상이라 하자. 그러면 유한 표시
$\mathcal{O}_X$-가군 $\mathcal{G}'$과 $\mathcal{O}_X$-가군의 사상
$\varphi' : \mathcal{G}' \to \mathcal{F}$가 존재하여
$\mathcal{G}'|_U = \mathcal{G}$이고 $\varphi'|_U = \varphi$이다.
\end{lemma}

\begin{proof}
증명의 시작은 보조정리 \ref{properties-lemma-extend}의 증명의 시작을 반복한다.
그래도 이를 자세히 적겠다.

\medskip\noindent
$n$을 아핀 열린집합 $U_i \subset X$,
$i = 1, \ldots , n$의 최소 개수라 하자. 이들은
$X = U \cup \bigcup U_i$가 되게 한다.
(여기서 $X$의 준콤팩트성을 사용한다.) $n = 1$인 경우에
보조정리를 증명할 수 있다고 하자. 그러면 쌍
$(\mathcal{G}, \varphi)$를 차례로
$(\mathcal{G}_1, \varphi_1)$로 $U \cup U_1$ 위에,
$(\mathcal{G}_2, \varphi_2)$로 $U \cup U_1 \cup U_2$ 위에,
$(\mathcal{G}_3, \varphi_3)$로 $U \cup U_1 \cup U_2 \cup U_3$ 위에
연장하는 식으로 계속할 수 있다. 따라서 $n = 1$인 경우로 환원된다.

\medskip\noindent
그러므로 $X = U \cup V$이고 $V$가 아핀이라고 가정해도 된다.
$X$가 준분리이고 $U$가 준콤팩트이므로
$U \cap V \subset V$는 준콤팩트이다. 계
$(V, U \cap V, \mathcal{F}|_V, \mathcal{G}|_{U \cap V}, \varphi|_{U \cap V})$
에 대해 보조정리를 증명하여 $(\mathcal{G}', \varphi')$를 얻고,
이 쌍이 $V$ 위에 있다고 하자. 그러면
$\mathcal{G}'$를 $V$ 위에서 주어진 층 $\mathcal{G}$와 $U$ 위에서
$U \cap V$ 위의 공통값을 따라 붙일 수 있고, 마찬가지로 사상
$\varphi'$을 사상 $\varphi$와 $U \cap V$ 위의 공통값을 따라
붙일 수 있다. 따라서 $X$가 아핀인 경우로 환원된다.

\medskip\noindent
$X = \Spec(R)$이라고 가정하자.
위의 보조정리 \ref{properties-lemma-extend-trivial}에 의해 준연접층
$\mathcal{H}$와 사상 $\psi : \mathcal{H} \to \mathcal{F}$를
찾을 수 있다. 이들은 $X$ 위에 있고, $\mathcal{G}$와
$\varphi$로 $U$ 위에서 제한된다.
보조정리 \ref{properties-lemma-extend}에 의해 유한형 준연접
$\mathcal{O}_X$-부분가군 $\mathcal{H}' \subset \mathcal{H}$를
$\mathcal{H}'|_U = \mathcal{G}$가 되도록 찾을 수 있다.
따라서 $\mathcal{H}$를 $\mathcal{H}'$로 바꾸고
$\psi$를 $\psi$의 $\mathcal{H}'$에 대한 제한으로 바꾸면
$\mathcal{H}$가 유한형이라고 가정할 수 있다.
보조정리 \ref{properties-lemma-finite-presentation-module}에 의해
$\mathcal{H} = \widetilde{N}$이고 $N$은 유한 생성
$R$-가군이다. 따라서 다음과 같은 준연접
$\mathcal{O}_X$-가군의 짧은 완전열을 이루는 전사사상이 존재한다.
$$
0 \to \mathcal{K} \to \mathcal{O}_X^{\oplus n} \to \mathcal{H} \to 0
$$
여기서 $\mathcal{K}$는 핵으로 정의된다.
$\mathcal{G}$가 유한 표시이고
$\mathcal{H}|_U = \mathcal{G}$이므로, Modules의 보조정리
\ref{modules-lemma-kernel-surjection-finite-free-onto-finite-presentation}에
의해 제한 $\mathcal{K}|_U$는 유한형
$\mathcal{O}_U$-가군이다. 따라서 보조정리 \ref{properties-lemma-extend}를
다시 사용하면 유한형 준연접
$\mathcal{O}_X$-부분가군 $\mathcal{K}' \subset \mathcal{K}$가
존재하여 $\mathcal{K}'|_U = \mathcal{K}|_U$이다. 보조정리에서
제기한 문제의 해답은 다음과 같이 놓는 것이다.
$$
\mathcal{G}' = \mathcal{O}_X^{\oplus n}/\mathcal{K}'
$$
이는 분명 유한 표시이고 $\mathcal{G}$를 $U$ 위에서 주도록 제한되며,
$\varphi'$은 합성
$$
\mathcal{G}' = \mathcal{O}_X^{\oplus n}/\mathcal{K}'
\to \mathcal{O}_X^{\oplus n}/\mathcal{K} = \mathcal{H} \xrightarrow{\psi}
\mathcal{F}.
$$
과 같다. 이로써 보조정리의 증명이 끝난다.
\end{proof}

\begin{lemma}
\label{properties-lemma-lift-finite-presentation}
$X$를 준콤팩트 준분리 스킴이라 하자. $U \subset X$를 준콤팩트
열린집합이라 하자. $\mathcal{G}$를 $\mathcal{O}_U$-가군이라 하자.
\begin{enumerate}
\item $\mathcal{G}$가 준연접이고 유한형이면 유한형 준연접
$\mathcal{O}_X$-가군 $\mathcal{G}'$가 존재하여
$\mathcal{G}'|_U = \mathcal{G}$이다.
\item $\mathcal{G}$가 유한 표시이면 유한 표시
$\mathcal{O}_X$-가군 $\mathcal{G}'$가 존재하여
$\mathcal{G}'|_U = \mathcal{G}$이다.
\end{enumerate}
\end{lemma}

\begin{proof}
(2)는 $\mathcal{F} = 0$인 보조정리
\ref{properties-lemma-extend-finite-presentation}의 특수한 경우이다.
(1)에 대해서는 먼저 보조정리 \ref{properties-lemma-extend-trivial}에 의해
$\mathcal{G} = \mathcal{F}|_U$라고 쓰되, 그 우변은 어떤 준연접
$\mathcal{O}_X$-가군에서 오도록 한다. 이어서
$\mathcal{G} = \mathcal{F}|_U$로 놓고 보조정리
\ref{properties-lemma-extend}를 적용한다.
\end{proof}

\noindent
다음 보조정리는 준콤팩트 준분리 스킴 위의 모든 준연접층이 유한 표시
$\mathcal{O}$-가군들의 여과 여극한임을 말한다. 실제로 다음
보조정리에서는 이를 유향 집합 위의 유향 여극한이라는 (아마 더
익숙한) 말로 다시 서술한다.

\begin{lemma}
\label{properties-lemma-directed-colimit-diagram-finite-presentation}
\begin{slogan}
준콤팩트 준분리 스킴 위의 준연접 가군은 유한 표시 가군들의
여과 여극한이다.
\end{slogan}
$X$를 스킴이라 하자. $X$가 준콤팩트이고 준분리라고 가정하자.
$\mathcal{F}$를 준연접 $\mathcal{O}_X$-가군이라 하자. 다음이 존재한다.
\begin{enumerate}
\item 여과 지표 범주 $\mathcal{I}$(Categories의 정의
\ref{categories-definition-directed}를 보라),
\item 도표 $\mathcal{I} \to \textit{Mod}(\mathcal{O}_X)$
(Categories의 절 \ref{categories-section-limits}를 보라),
$i \mapsto \mathcal{F}_i$,
\item $\mathcal{O}_X$-가군의 사상
$\varphi_i : \mathcal{F}_i \to \mathcal{F}$.
\end{enumerate}
여기서 각 $\mathcal{F}_i$는 유한 표시이고, 사상들 $\varphi_i$는
동형사상
$$
\colim_i \mathcal{F}_i
=
\mathcal{F}.
$$
을 유도한다.
\end{lemma}

\begin{proof}
집합 $I$를 선택하고 각 $i \in I$에 대하여 유한 표시
$\mathcal{O}_X$-가군과 $\mathcal{O}_X$-가군의 준동형
$\varphi_i : \mathcal{F}_i \to \mathcal{F}$를 다음 성질을
만족하도록 선택하자. 임의의 $\psi : \mathcal{G} \to \mathcal{F}$에
대하여, 여기서 $\mathcal{G}$는 유한 표시인데, 어떤 $i \in I$가
존재하여 동형사상 $\alpha : \mathcal{F}_i \to \mathcal{G}$가
$\varphi_i = \psi \circ \alpha$를 만족한다. 이러한 집합이 존재함은
Modules의 보조정리
\ref{modules-lemma-set-isomorphism-classes-finite-type-modules}에서
명백하다(그 증명도 보라).
$\mathcal{I}$를 $\Ob(\mathcal{I}) = I$인 범주라 쓰고,
$i, i' \in I$가 주어지면 다음과 같이 놓는다.
$$
\Mor_\mathcal{I}(i, i') =
\{\alpha : \mathcal{F}_i \to \mathcal{F}_{i'} \mid
\alpha \circ \varphi_{i'} = \varphi_i
\}.
$$
$\mathcal{I}$가 여과 범주이고
$\mathcal{F} = \colim_i \mathcal{F}_i$라고 주장한다.

\medskip\noindent
$i, i' \in I$라 하자. 그러면 사상
$$
\mathcal{F}_i \oplus \mathcal{F}_{i'} \longrightarrow \mathcal{F}
$$
을 생각할 수 있으며, 이는 $\varphi_i$와 $\varphi_{i'}$의
직합이다. 유한 표시 $\mathcal{O}_X$-가군들의 직합은 유한
표시이므로 어떤 $i'' \in I$가 존재하여
$\varphi_{i''} : \mathcal{F}_{i''} \to \mathcal{F}$가 위에 표시한
$\mathcal{F}$로 가는 화살표와 동형임을 알 수 있다. 가환 도식들
$$
\xymatrix{
\mathcal{F}_i \ar[r] \ar[d] & \mathcal{F} \ar@{=}[d] \\
\mathcal{F}_i \oplus \mathcal{F}_{i'} \ar[r] & \mathcal{F}
}
\quad
\text{and}
\quad
\xymatrix{
\mathcal{F}_{i'} \ar[r] \ar[d] & \mathcal{F} \ar@{=}[d] \\
\mathcal{F}_i \oplus \mathcal{F}_{i'} \ar[r] & \mathcal{F}
}
$$
이 있으므로 사상 $i \to i''$와 $i' \to i''$가
$\mathcal{I}$ 안에 있다. 다음으로 $i, i' \in I$와 사상
$\alpha, \beta : i \to i'$가 주어졌다고 하자(이는
$\mathcal{O}_X$-가군 사상
$\alpha, \beta : \mathcal{F}_i \to \mathcal{F}_{i'}$에 대응한다).
이 경우 쌍대등화자
$$
\mathcal{G} =
\Coker(
\mathcal{F}_i \xrightarrow{\alpha - \beta} \mathcal{F}_{i'}
)
$$
를 생각하자. $\mathcal{G}$는 유한 표시
$\mathcal{O}_X$-가군임에 유의하라. 범주 $\mathcal{I}$의 사상의
정의에 의해 $\varphi_{i'} \circ \alpha = \varphi_{i'} \circ \beta$이므로
유도된 사상 $\psi : \mathcal{G} \to \mathcal{F}$를 얻는다.
따라서 다시 쌍 $(\mathcal{G}, \psi)$는
쌍 $(\mathcal{F}_{i''}, \varphi_{i''})$와 어떤 $i''$에 대하여
동형이다. 따라서 사상 $i' \to i''$가 $\mathcal{I}$ 안에 존재하고
$\alpha$와 $\beta$를 같게 만든다. 이로써 범주 $\mathcal{I}$가
여과임을 보였다.

\medskip\noindent
이제 도표의 여극한이 $\mathcal{F}$임을 보여야 한다.
여극한의 정의와 범주 $\mathcal{I}$의 정의에 의해 표준 사상
$$
\varphi :
\colim_i \mathcal{F}_i
\longrightarrow
\mathcal{F}.
$$
이 있다. $x \in X$를 잡자. $\varphi_x$가 동형사상임을 보이자.
다음을 상기하라.
$$
(\colim_i \mathcal{F}_i)_x
=
\colim_i \mathcal{F}_{i, x},
$$
Sheaves의 절 \ref{sheaves-section-limits-sheaves}을 보라.
먼저 사상 $\varphi_x$가 단사임을 보이자.
$s \in \mathcal{F}_{i, x}$를 잡고, $s$가 $\mathcal{F}_x$에서
영으로 사상된다고 하자. 그러면 준콤팩트 열린집합 $U$가 존재하여
$s$는 $s \in \mathcal{F}_i(U)$에서 오고
$\varphi_i(s) = 0$이 $\mathcal{F}(U)$에서 성립한다.
보조정리 \ref{properties-lemma-extend}에 의해 유한형 준연접 부분층
$\mathcal{K} \subset \Ker(\varphi_i)$를 찾을 수 있고, 이는
$\mathcal{O}_U$-부분가군으로서 준연접이고, $\mathcal{F}_i$ 안에서
$s$로 생성되는 부분가군으로 제한된다.
$\mathcal{K}|_U = \mathcal{O}_U\cdot s \subset \mathcal{F}_i|_U$.
분명 $\mathcal{F}_i/\mathcal{K}$는 유한 표시이고 사상
$\varphi_i$는 몫사상
$\mathcal{F}_i \to \mathcal{F}_i/\mathcal{K}$를 통해 분해된다.
따라서 $i' \in I$와 사상
$\alpha : \mathcal{F}_i \to \mathcal{F}_{i'}$를
$\mathcal{I}$ 안에서 찾을 수 있고, 이 사상은 몫사상
$\mathcal{F}_i \to \mathcal{F}_i/\mathcal{K}$와 식별된다.
그러면 절단 $s$는 $\mathcal{F}_{i'}(U)$에서 영으로 사상되고,
특히
$(\colim_i \mathcal{F}_i)_x =
\colim_i \mathcal{F}_{i, x}$에서도 그러하다. 단사성이 따른다.
마지막으로 사상 $\varphi_x$가 전사임을 보이자.
$s \in \mathcal{F}_x$를 잡자. 준콤팩트 열린 근방
$U \subset X$를 $x$에 대해 선택하여 $s$가 절단
$s \in \mathcal{F}(U)$에 대응하게 하자. 사상
$s : \mathcal{O}_U \to \mathcal{F}$를 생각하자($s$에 의한 곱셈).
보조정리 \ref{properties-lemma-extend-finite-presentation}에 의해 유한 표시
$\mathcal{O}_X$-가군 $\mathcal{G}$와 $\mathcal{O}_X$-가군 사상
$\mathcal{G} \to \mathcal{F}$가 존재하여
$\mathcal{G}|_U \to \mathcal{F}|_U$는
$s : \mathcal{O}_U \to \mathcal{F}$와 식별된다.
다시 $\mathcal{I}$의 정의에 의해 어떤 $i \in I$가 존재하여
$\mathcal{G} \to \mathcal{F}$는
$\varphi_i : \mathcal{F}_i \to \mathcal{F}$와 동형이다.
분명 절단 $s' \in \mathcal{F}_i(U)$가 존재하여
$s \in \mathcal{F}(U)$로 사상된다. 이로써 전사성이 증명되고
보조정리의 증명이 끝난다.
\end{proof}

\begin{lemma}
\label{properties-lemma-directed-colimit-finite-presentation}
$X$를 스킴이라 하자. $X$가 준콤팩트이고 준분리라고 가정하자.
$\mathcal{F}$를 준연접 $\mathcal{O}_X$-가군이라 하자. 다음이 존재한다.
\begin{enumerate}
\item 유향 집합 $I$(Categories의 정의
\ref{categories-definition-directed-set}를 보라),
\item 계 $(\mathcal{F}_i, \varphi_{ii'})$가 $I$ 위에 있고
$\textit{Mod}(\mathcal{O}_X)$ 안에 놓인다
(Categories의 정의 \ref{categories-definition-system-over-poset}를 보라),
\item $\mathcal{O}_X$-가군의 사상
$\varphi_i : \mathcal{F}_i \to \mathcal{F}$.
\end{enumerate}
여기서 각 $\mathcal{F}_i$는 유한 표시이고, 사상들 $\varphi_i$는
동형사상
$$
\colim_i \mathcal{F}_i
=
\mathcal{F}.
$$
을 유도한다.
\end{lemma}

\begin{proof}
보조정리 \ref{properties-lemma-directed-colimit-diagram-finite-presentation}와
Categories의 보조정리 \ref{categories-lemma-directed-category-system}의
직접적인 귀결이다(층의 범주에는 $\mathcal{O}_X$-가군의 여극한이
존재한다는 사실도 함께 사용한다. Sheaves의 절
\ref{sheaves-section-limits-sheaves}을 보라).
\end{proof}

\begin{lemma}
\label{properties-lemma-finite-directed-colimit-surjective-maps}
$X$를 스킴이라 하자. $X$가 준콤팩트이고 준분리라고 가정하자.
$\mathcal{F}$를 유한형 준연접 $\mathcal{O}_X$-가군이라 하자.
그러면 $\mathcal{F} = \colim \mathcal{F}_i$라고 쓸 수 있고,
$\mathcal{F}_i$는 유한 표시이며 모든 전이사상
$\mathcal{F}_i \to \mathcal{F}_{i'}$는 전사이다.
\end{lemma}

\begin{proof}
$\mathcal{F} = \colim \mathcal{G}_i$를 유한 표시
$\mathcal{O}_X$-가군들의 여과 여극한으로 쓰자(보조정리
\ref{properties-lemma-directed-colimit-finite-presentation}).
$\mathcal{G}_i \to \mathcal{F}$가 어떤 $i$에 대하여
전사라고 주장한다. 실제로 유한 아핀 열린 덮개
$X = U_1 \cup \ldots \cup U_m$을 선택하자.
절단 $s_{jl} \in \mathcal{F}(U_j)$들을 선택하여
$\mathcal{F}|_{U_j}$를 생성하게 하자. 보조정리
\ref{properties-lemma-finite-type-module}를 보라. Sheaves의 보조정리
\ref{sheaves-lemma-directed-colimits-sections}에 의해 $s_{jl}$는
$\mathcal{G}_i \to \mathcal{F}$의 상에 충분히 큰 $i$에 대하여
속한다. 따라서 $\mathcal{G}_i \to \mathcal{F}$는 충분히 큰
$i$에 대하여 전사이다. 그러한 $i$를 선택하고
$\mathcal{K} \subset \mathcal{G}_i$를 사상
$\mathcal{G}_i \to \mathcal{F}$의 핵이라 하자.
$\mathcal{K} = \colim \mathcal{K}_a$를 그 유한형 준연접
부분가군들의 여과 여극한으로 쓰자(보조정리
\ref{properties-lemma-quasi-coherent-colimit-finite-type}). 그러면
$\mathcal{F} = \colim \mathcal{G}_i/\mathcal{K}_a$가 보조정리에서
제기한 문제의 해답이다.
\end{proof}

\begin{lemma}
\label{properties-lemma-application-directed-colimit}
$X$를 준콤팩트 준분리 스킴이라 하자.
$\mathcal{F}$를 유한형 준연접 $\mathcal{O}_X$-가군이라 하자.
$U \subset X$를 준콤팩트 열린집합이라 하고
$\mathcal{F}|_U$가 유한 표시라고 하자. 그러면
$\mathcal{O}_X$-가군의 사상
$\varphi : \mathcal{G} \to \mathcal{F}$가 존재하여
(a) $\mathcal{G}$는 유한 표시이고,
(b) $\varphi$는 전사이며,
(c) $\varphi|_U$는 동형사상이다.
\end{lemma}

\begin{proof}
$\mathcal{F} = \colim \mathcal{F}_i$를 유향 여극한으로 쓰되, 각
$\mathcal{F}_i$가 유한 표시가 되게 하자. 보조정리
\ref{properties-lemma-directed-colimit-finite-presentation}를 보라.
유한 아핀 열린 덮개 $X = \bigcup V_j$를 선택하고, 유한 개의 절단
$s_{jl} \in \mathcal{F}(V_j)$을 선택하여
$\mathcal{F}|_{V_j}$를 생성하게 하자. 보조정리
\ref{properties-lemma-finite-type-module}를 보라. Sheaves의 보조정리
\ref{sheaves-lemma-directed-colimits-sections}에 의해 $s_{jl}$는
$\mathcal{F}_i \to \mathcal{F}$의 상에 충분히 큰 $i$에 대하여
속한다. 따라서 $\mathcal{F}_i \to \mathcal{F}$는 충분히 큰
$i$에 대하여 전사이다. 그러한 $i$를 선택하고
$\mathcal{K} \subset \mathcal{F}_i$를 사상
$\mathcal{F}_i \to \mathcal{F}$의 핵이라 하자.
$\mathcal{F}_U$가 유한 표시이므로 $\mathcal{K}|_U$는 유한형이다.
Modules의 보조정리
\ref{modules-lemma-kernel-surjection-finite-free-onto-finite-presentation}를
보라. 따라서 유한형 준연접 부분가군
$\mathcal{K}' \subset \mathcal{K}$를
$\mathcal{K}'|_U = \mathcal{K}|_U$가 되도록 찾을 수 있다.
보조정리 \ref{properties-lemma-extend}를 보라. 그러면
$\mathcal{G} = \mathcal{F}_i/\mathcal{K}'$와 주어진 사상
$\mathcal{G} \to \mathcal{F}$가 해답이다.
\end{proof}

\noindent
$X$를 스킴이라 하자. 다음 보조정리에서는
{\it 유한 표시인 준연접 $\mathcal{O}_X$-대수 $\mathcal{A}$}라는
개념을 사용한다. 이는 모든 아핀 열린집합
$\Spec(R) \subset X$에 대하여 $\mathcal{A} = \widetilde{A}$이고,
여기서 $A$는 $R$-대수로서 유한 표시인 (가환) $R$-대수라는 뜻이다.

\begin{lemma}
\label{properties-lemma-algebra-directed-colimit-finite-presentation}
$X$를 스킴이라 하자. $X$가 준콤팩트이고 준분리라고 가정하자.
$\mathcal{A}$를 준연접 $\mathcal{O}_X$-대수라 하자. 다음이 존재한다.
\begin{enumerate}
\item 유향 집합 $I$(Categories의 정의
\ref{categories-definition-directed-set}를 보라),
\item 계 $(\mathcal{A}_i, \varphi_{ii'})$가 $I$ 위에서
$\mathcal{O}_X$-대수의 범주 안에 놓인다.
\item $\mathcal{O}_X$-대수의 사상
$\varphi_i : \mathcal{A}_i \to \mathcal{A}$.
\end{enumerate}
여기서 각 $\mathcal{A}_i$는 유한 표시 준연접
$\mathcal{O}_X$-대수이고, 사상들 $\varphi_i$는 동형사상
$$
\colim_i \mathcal{A}_i
=
\mathcal{A}.
$$
을 유도한다.
\end{lemma}

\begin{proof}
먼저 $\mathcal{A} = \colim_i \mathcal{F}_i$를 보조정리
\ref{properties-lemma-directed-colimit-finite-presentation}에서와 같이 유한 표시
준연접층들의 유향 여극한으로 쓰자. 각 $i$에 대하여
$\mathcal{B}_i = \text{Sym}(\mathcal{F}_i)$를
$\mathcal{F}_i$로 생성되는 $\mathcal{O}_X$ 위의 대칭대수라 하자.
$\mathcal{I}_i = \Ker(\mathcal{B}_i \to \mathcal{A})$라 쓰자.
$\mathcal{I}_i = \colim_j \mathcal{F}_{i, j}$라 쓰되,
$\mathcal{F}_{i, j}$는 $\mathcal{I}_i$의 유한형 준연접
부분가군이라 하자. 보조정리
\ref{properties-lemma-quasi-coherent-colimit-finite-type}를 보라.
$\mathcal{I}_{i, j} \subset \mathcal{I}_i$를
$\mathcal{B}_i$-아이디얼 가운데 $\mathcal{F}_{i, j}$로 생성되는
것과 같게 놓자.
$\mathcal{A}_{i, j} = \mathcal{B}_i/\mathcal{I}_{i, j}$라 놓자.
그러면 $\mathcal{A}_{i, j}$는 유한 표시 준연접
$\mathcal{O}_X$-대수이다. $(i, j) \leq (i', j')$를 다음과 같이
정의하자. $i \leq i'$이고 사상
$\mathcal{B}_i \to \mathcal{B}_{i'}$가 아이디얼
$\mathcal{I}_{i, j}$를 아이디얼 $\mathcal{I}_{i', j'}$ 안으로
보내야 한다. 그러면
$\mathcal{A} = \colim_{i, j} \mathcal{A}_{i, j}$임이 명백하다.
\end{proof}

\noindent
$X$를 스킴이라 하자. 다음 보조정리에서는
{\it 유한형 준연접 $\mathcal{O}_X$-대수 $\mathcal{A}$}라는 개념을
사용한다. 이는 모든 아핀 열린집합 $\Spec(R) \subset X$에 대하여
$\mathcal{A} = \widetilde{A}$이고, 여기서 $A$는 $R$-대수로서
유한형인 (가환) $R$-대수라는 뜻이다.

\begin{lemma}
\label{properties-lemma-algebra-directed-colimit-finite-type}
$X$를 스킴이라 하자. $X$가 준콤팩트이고 준분리라고 가정하자.
$\mathcal{A}$를 준연접 $\mathcal{O}_X$-대수라 하자. 그러면
$\mathcal{A}$는 그 유한형 준연접 $\mathcal{O}_X$-부분대수들의
유향 여극한이다.
\end{lemma}

\begin{proof}
$\mathcal{A}_1, \mathcal{A}_2 \subset \mathcal{A}$가 유한형 준연접
$\mathcal{O}_X$-부분대수이면
$\mathcal{A}_1 \otimes_{\mathcal{O}_X} \mathcal{A}_2 \to \mathcal{A}$의
상도 유한형 준연접 $\mathcal{O}_X$-부분대수이고(몇몇 세부사항은
생략한다), $\mathcal{A}_1$과 $\mathcal{A}_2$를 모두 포함한다.
이로써 이 계가 유향임을 알 수 있다.
$\mathcal{A}$가 이 계의 여극한임을 보이기 위해,
$\mathcal{A} = \colim_i \mathcal{A}_i$를 보조정리
\ref{properties-lemma-algebra-directed-colimit-finite-presentation}에서와 같이
유한 표시 준연접 $\mathcal{O}_X$-대수들의 유향 여극한으로 쓰자.
그러면 상들
$\mathcal{A}'_i = \Im(\mathcal{A}_i \to \mathcal{A})$은
$\mathcal{A}$의 유한형 준연접 부분대수이다.
$\mathcal{A}$가 이들의 여극한이므로 결과가 따른다.
\end{proof}

\noindent
$X$를 스킴이라 하자. 다음 보조정리에서는
{\it 유한 (각각 정수적) 준연접 $\mathcal{O}_X$-대수
$\mathcal{A}$}라는 개념을 사용한다. 이는 모든 아핀 열린집합
$\Spec(R) \subset X$에 대하여 $\mathcal{A} = \widetilde{A}$이고,
여기서 $A$는 $R$-대수로서 유한 (각각 정수적)인 (가환)
$R$-대수라는 뜻이다.

\begin{lemma}
\label{properties-lemma-finite-algebra-directed-colimit-finite-finitely-presented}
$X$를 스킴이라 하자. $X$가 준콤팩트이고 준분리라고 가정하자.
$\mathcal{A}$를 유한 준연접 $\mathcal{O}_X$-대수라 하자. 그러면
$\mathcal{A} = \colim \mathcal{A}_i$는 유한이고 유한 표시인 준연접
$\mathcal{O}_X$-대수들의 유향 여극한이며, 모든 전이사상
$\mathcal{A}_{i'} \to \mathcal{A}_i$는 전사이다.
\end{lemma}

\begin{proof}
보조정리 \ref{properties-lemma-finite-directed-colimit-surjective-maps}에 의해
유한 표시 $\mathcal{O}_X$-가군 $\mathcal{F}$와 전사사상
$\mathcal{F} \to \mathcal{A}$가 존재한다. 대수 구조를 사용하여
전사사상
$$
\text{Sym}^*_{\mathcal{O}_X}(\mathcal{F}) \longrightarrow \mathcal{A}
$$
을 얻는다. 그 핵을 $\mathcal{J}$라 쓰자.
$\mathcal{J} = \colim \mathcal{E}_i$를 유한형
$\mathcal{O}_X$-부분가군 $\mathcal{E}_i$들의 여과 여극한으로
쓰자(보조정리 \ref{properties-lemma-quasi-coherent-colimit-finite-type}).
다음과 같이 놓자.
$$
\mathcal{A}_i = \text{Sym}^*_{\mathcal{O}_X}(\mathcal{F})/(\mathcal{E}_i)
$$
여기서 $(\mathcal{E}_i)$는
$\mathcal{E}_i \to \text{Sym}^*_{\mathcal{O}_X}(\mathcal{F})$의
상으로 생성되는 아이디얼 층을 뜻한다.
그러면 각 $\mathcal{A}_i$는 유한 표시 $\mathcal{O}_X$-대수이고,
전이사상들은 전사이며,
$\mathcal{A} = \colim \mathcal{A}_i$이다. 증명을 끝내려면
$\mathcal{A}_i$가 유한 $\mathcal{O}_X$-대수임을 충분히 큰
$i$에 대하여
보이면 된다. 이를 위해 아핀 열린 덮개
$X = U_1 \cup \ldots \cup U_m$을 선택하자. 생성원들
$f_{j, 1}, \ldots, f_{j, N_j} \in \Gamma(U_i, \mathcal{F})$을 잡자.
$\mathcal{A}(U_j)$는 유한 $\mathcal{O}_X(U_j)$-대수이므로, 각
$k$에 대하여 일계수 다항식
$P_{j, k} \in \mathcal{O}(U_j)[T]$가 존재하여
$P_{j, k}(f_{j, k})$가 $\mathcal{A}(U_j)$에서 영이다.
구성에 의해 $\mathcal{A} = \colim \mathcal{A}_i$이므로,
$P_{j, k}(f_{j, k}) = 0$이 $\mathcal{A}_i(U_j)$에서 충분히 큰
모든 $i$에 대하여 성립한다. 그러한 $i$에 대하여 대수들
$\mathcal{A}_i$는 유한이다.
\end{proof}

\begin{lemma}
\label{properties-lemma-integral-algebra-directed-colimit-finite}
$X$를 스킴이라 하자. $X$가 준콤팩트이고 준분리라고 가정하자.
$\mathcal{A}$를 정수적 준연접 $\mathcal{O}_X$-대수라 하자. 그러면
\begin{enumerate}
\item $\mathcal{A}$는 그 유한 준연접 $\mathcal{O}_X$-부분대수들의
유향 여극한이고,
\item $\mathcal{A}$는 유한이고 유한 표시인 준연접
$\mathcal{O}_X$-대수들의 직접 여극한이다.
\end{enumerate}
\end{lemma}

\begin{proof}
보조정리 \ref{properties-lemma-algebra-directed-colimit-finite-type}에 의해
$\mathcal{A} = \colim \mathcal{A}_i$이고,
$\mathcal{A}_i \subset \mathcal{A}$는 유한형 준연접
$\mathcal{O}_X$-대수 전체를 달린다. 임의의 유한형 준연접
$\mathcal{O}_X$-부분대수 가운데 $\mathcal{A}$의 부분대수인 것은
유한이다(Algebra의
보조정리 \ref{algebra-lemma-characterize-finite-in-terms-of-integral}를
$\mathcal{A}_i(U) \subset \mathcal{A}(U)$에 적용하라. 여기서
$U$는 $X$ 안의 아핀 열린집합이다).
이로써 (1)이 증명된다.

\medskip\noindent
(2)를 증명하기 위해, 보조정리
\ref{properties-lemma-directed-colimit-finite-presentation}를 사용하여
$\mathcal{A} = \colim \mathcal{F}_i$를 유한 표시
$\mathcal{O}_X$-가군들의 여극한으로 쓰자. 각 $i$에 대하여
사상
$$
\text{Sym}^*_{\mathcal{O}_X}(\mathcal{F}_i) \longrightarrow \mathcal{A}
$$
의 핵을 $\mathcal{J}_i$라 하자.
$i' \geq i$에 대하여 유도된 사상
$\mathcal{J}_i \to \mathcal{J}_{i'}$가 있고,
$\mathcal{A} =
\colim \text{Sym}^*_{\mathcal{O}_X}(\mathcal{F}_i)/\mathcal{J}_i$이다.
또한 준연접 $\mathcal{O}_X$-대수
$\text{Sym}^*_{\mathcal{O}_X}(\mathcal{F}_i)/\mathcal{J}_i$는
유한이다(위를 보라). $\mathcal{J}_i = \colim \mathcal{E}_{ik}$를
유한 표시 $\mathcal{O}_X$-가군들의 여극한으로 쓰자.
$i' \geq i$와 $k$가 주어지면 어떤 $k'$이 존재하여
사상 $\mathcal{E}_{ik} \to \mathcal{E}_{i'k'}$가 있고, 도식
$$
\xymatrix{
\mathcal{J}_i \ar[r] & \mathcal{J}_{i'} \\
\mathcal{E}_{ik} \ar[u] \ar[r] & \mathcal{E}_{i'k'} \ar[u]
}
$$
을 가환하게 한다. 이는 Modules의 보조정리
\ref{modules-lemma-finite-presentation-quasi-compact-colimit}에서
따른다. 이로부터 사상
$$
\mathcal{A}_{ik} =
\text{Sym}^*_{\mathcal{O}_X}(\mathcal{F}_i)/(\mathcal{E}_{ik})
\longrightarrow
\text{Sym}^*_{\mathcal{O}_X}(\mathcal{F}_{i'})/(\mathcal{E}_{i'k'}) =
\mathcal{A}_{i'k'}
$$
이 유도된다. 여기서 $(\mathcal{E}_{ik})$는
$\mathcal{E}_{ik}$로 생성되는 아이디얼을 뜻한다. 준연접
$\mathcal{O}_X$-대수 $\mathcal{A}_{ki}$는 충분히 큰 $k$에 대하여
유한 표시이고 유한이다(보조정리
\ref{properties-lemma-finite-algebra-directed-colimit-finite-finitely-presented}의
증명을 보라). 마지막으로
$$
\colim \mathcal{A}_{ik} = \colim \mathcal{A}_i = \mathcal{A}
$$
이다. 첫 번째 등식은 보조정리
\ref{properties-lemma-finite-algebra-directed-colimit-finite-finitely-presented}의
증명에서 보였고, 두 번째 등식은 $\mathcal{A}$가 가군들
$\mathcal{F}_i$의 여극한이기 때문에 성립한다.
\end{proof}

\section{Gabber의 결과}
\label{properties-section-gabber}

\noindent
이 절에서는 모든 스킴 위에 기수 $\kappa$가 존재하여 각 준연접
가군 $\mathcal{F}$가 그 준연접 $\kappa$-생성 부분층들의 합집합이
되게 한다는 Gabber의 결과를 증명한다. 이로부터 스킴 위의 준연접층의
범주는 극한과 충분한 단사대상을 갖는 그로텐디크 아벨
범주임이 따른다\footnote{Akhil Mathew의
\href{https://amathew.wordpress.com/2011/07/30/quasi-coherent-sheaves-presentable-categories-and-a-result-of-gabber/}{블로그 글}에
훌륭하게 설명되어 있다.}.

\begin{definition}
\label{properties-definition-kappa-generated}
$(X, \mathcal{O}_X)$를 환 달린 공간이라 하자. $\kappa$를 무한
기수라 하자. $\mathcal{O}_X$-가군의 층 $\mathcal{F}$를
{\it $\kappa$-생성}이라고 하는 것은 다음 조건이 성립한다는 뜻이다.
열린 덮개 $X = \bigcup U_i$가 존재하여 $\mathcal{F}|_{U_i}$가
부분집합 $R_i \subset \mathcal{F}(U_i)$로 생성되고 그 기수가
$\kappa$ 이하이다.
\end{definition}

\noindent
많아야 $\kappa$개인 $\kappa$-생성 가군들의 직합은 다시
$\kappa$-생성임에 유의하라. 이는
$\kappa \otimes \kappa = \kappa$이기 때문이다. Sets의 절
\ref{sets-section-cardinals}을 보라. 특히 두 $\kappa$-생성 가군의
직합에 대해 이것이 성립한다. 또한 $\kappa$-생성 층의 몫은
$\kappa$-생성이다. (하지만 부분가군에 대해서도 같을 필요는 없다.)

\begin{lemma}
\label{properties-lemma-set-of-iso-classes}
$(X, \mathcal{O}_X)$를 환 달린 공간이라 하자. $\kappa$를 기수라
하자. 집합 $T$와 족 $(\mathcal{F}_t)_{t \in T}$가 존재하며,
이 족은 $\kappa$-생성 $\mathcal{O}_X$-가군들로 이루어지고 모든 $\kappa$-생성
$\mathcal{O}_X$-가군은 $\mathcal{F}_t$ 가운데 하나와 동형이다.
\end{lemma}

\begin{proof}
반복을 허용하지 않으면 $X$의 덮개들은 하나의 집합을 이룬다.
$X = \bigcup U_i$가 덮개이고 $\mathcal{F}_i$가
$\mathcal{O}_{U_i}$-가군이라고 하자. 그러면 붙임 사상들이 하나의
집합을 이루므로 $\mathcal{O}_X$-가군 $\mathcal{F}$ 중에서
$\mathcal{F}|_{U_i} \cong \mathcal{F}_i$라는 성질을 가진 것들의
동형류도 하나의 집합을 이룬다. 따라서 몫들
$\oplus_{k \in \kappa} \mathcal{O}_X \to \mathcal{F}$의
(동형류의) 모임이 임의의 환 달린 공간 $X$에 대해 집합임을 증명하는
것으로 환원된다. 이는 명백하다.
\end{proof}

\noindent
다음은 이 절의 제목이 가리키는 결과이다.

\begin{lemma}
\label{properties-lemma-colimit-kappa}
$X$를 스킴이라 하자. 기수 $\kappa$가 존재하여 모든 준연접
가군 $\mathcal{F}$는 그 준연접 $\kappa$-생성 부분가군들의 유향
여극한이다.
\end{lemma}

\begin{proof}
아핀 열린 덮개 $X = \bigcup_{i \in I} U_i$를 선택하자. 각 쌍
$i, j$에 대하여 아핀 열린 덮개
$U_i \cap U_j = \bigcup_{k \in I_{ij}} U_{ijk}$를 선택하자.
$U_i = \Spec(A_i)$ 및 $U_{ijk} = \Spec(A_{ijk})$라 쓰자.
$\kappa$를 $\geq$ 관계로 집합들 $I$, $I_{ij}$ 각각의 기수를
지배하는 임의의 무한 기수라 하자.

\medskip\noindent
$\mathcal{F}$를 준연접층이라 하자.
$M_i = \mathcal{F}(U_i)$ 및
$M_{ijk} = \mathcal{F}(U_{ijk})$라 놓자. 다음에 유의하라.
$$
M_i \otimes_{A_i} A_{ijk} = M_{ijk} = M_j \otimes_{A_j} A_{ijk}.
$$
Schemes의 보조정리 \ref{schemes-lemma-widetilde-pullback}를 보라.
선택공리를 사용하여 사상
$$
(i, j, k, m) \mapsto S(i, j, k, m)
$$
을 선택하자. 이는 모든 $i, j \in I$, $k \in I_{ij}$ 및
$m \in M_i$에 유한 부분집합
$S(i, j, k, m) \subset M_j$를 대응시키며,
$$
m \otimes 1 = \sum\nolimits_{m' \in S(i, j, k, m)} m' \otimes a_{m'}
$$
가 $M_{ijk}$ 안에서 어떤 $a_{m'} \in A_{ijk}$에 대하여 성립하게 한다.
또한 $S(i, i, k, m) = \{m\}$라고 약속하자. 이 약속은 위와 같은 모든
$i, j = i, k, m$에 적용한다. 이러한 사상을 하나 고정한다.

\medskip\noindent
부분집합들의 족 $\mathcal{S} = (S_i)_{i \in I}$가 주어졌고, 각
$S_i \subset M_i$의 기수가 $\kappa$ 이하이면
$\mathcal{S}' = (S'_i)$라 놓는데, 여기서
$$
S'_j = \bigcup\nolimits_{(i, k, m)\text{ such that }m \in S_i}
S(i, j, k, m)
$$
이다. $S_i \subset S'_i$임에 유의하라. 또한 $S'_i$는 기수가
많아야 $\kappa$인 집합에 걸친 유한집합들의 합집합이므로 기수가
많아야 $\kappa$이다.
$\mathcal{S}^{(0)} = \mathcal{S}$,
$\mathcal{S}^{(1)} = \mathcal{S}'$라 놓고 귀납적으로
$\mathcal{S}^{(n + 1)} = (\mathcal{S}^{(n)})'$라 놓자. 그리고
$\mathcal{S}^{(\infty)} = \bigcup_{n \geq 0} \mathcal{S}^{(n)}$라 놓자.
$\mathcal{S}^{(\infty)} = (S^{(\infty)}_i)$라 쓰면, 임의의 원소
$m \in S^{(\infty)}_i$에 대하여 $m$의 $M_{ijk}$ 안에서의 상은
유한합 $\sum m' \otimes a_{m'}$으로 쓸 수 있고, 여기서
$m' \in S_j^{(\infty)}$이다. 따라서
$$
N_i = A_i\text{-submodule of }M_i\text{ generated by }S^{(\infty)}_i
$$
라고 놓으면
$$
N_i \otimes_{A_i} A_{ijk} = N_j \otimes_{A_j} A_{ijk}.
$$
이다. 이는 $M_{ijk}$의 부분가군으로서의 등식이다. 그러므로
준연접 부분층 $\mathcal{G} \subset \mathcal{F}$가 존재하여
$\mathcal{G}(U_i) = N_i$이다. 또한 구성에 의해 층 $\mathcal{G}$는
$\kappa$-생성이다.

\medskip\noindent
$\{\mathcal{G}_t\}_{t \in T}$를 $\kappa$-생성 준연접 부분층들의
집합이라 하자. $t, t' \in T$이면
$\mathcal{G}_t + \mathcal{G}_{t'}$도 $\kappa$-생성 준연접
부분층이다. 실제로 이는 사상
$\mathcal{G}_t \oplus \mathcal{G}_{t'} \to \mathcal{F}$의 상이다.
따라서 포함으로 순서지은 이 계는 유향이다. 위의 논의는
$\mathcal{F}$의 $U_i$ 위의 모든 절단이 어떤
$\mathcal{G}_t$에 속함을 보인다($\mathcal{S}$를 주어진 절단이
$S_i$의 원소가 되도록 택하여 시작할 수 있기 때문이다). 따라서 원하는
대로 $\colim_t \mathcal{G}_t \to \mathcal{F}$는 단사이고 전사이다.
\end{proof}

\begin{proposition}
\label{properties-proposition-coherator}
$X$를 스킴이라 하자.
\begin{enumerate}
\item 범주 $\QCoh(\mathcal{O}_X)$는 그로텐디크 아벨 범주이다.
따라서 $\QCoh(\mathcal{O}_X)$는 충분한 단사대상과 모든 극한을 갖는다.
\item 포함 함자
$\QCoh(\mathcal{O}_X) \to \textit{Mod}(\mathcal{O}_X)$는
오른쪽 수반을 갖는다\footnote{이 함자를 때때로
{\it 코히레이터}라고 부른다.}
$$
Q : \textit{Mod}(\mathcal{O}_X) \longrightarrow \QCoh(\mathcal{O}_X)
$$
모든 준연접층 $\mathcal{F}$에 대하여 수반 사상
$Q(\mathcal{F}) \to \mathcal{F}$는 동형사상이다.
\end{enumerate}
\end{proposition}

\begin{proof}
(1)은 $\QCoh(\mathcal{O}_X)$가 (a) 모든 여극한을 갖고, (b) 여과
여극한이 완전하며, (c) 생성자를 갖는다는 뜻이다. Injectives의 절
\ref{injectives-section-grothendieck-conditions}를 보라.
Schemes의 절 \ref{schemes-section-quasi-coherent}에 의해
$\QCoh(\mathcal{O}_X)$ 안의 여극한은 존재하고
$\textit{Mod}(\mathcal{O}_X)$ 안의 여극한과 일치한다. Modules의
보조정리 \ref{modules-lemma-limits-colimits}에 의해 여과 여극한은
완전하다. 따라서 (a)와 (b)가 성립한다.
생성자 $U$를 구성하기 위해 보조정리 \ref{properties-lemma-colimit-kappa}와
같은 기수 $\kappa$를 선택하자. 보조정리
\ref{properties-lemma-set-of-iso-classes}에서와 같은 준연접층의 모임
$(\mathcal{F}_t)_{t \in T}$를 선택하자. 이들은 $\kappa$-생성이다. 다음과 같이 놓자.
$U = \bigoplus_{t \in T} \mathcal{F}_t$. 모든
$\QCoh(\mathcal{O}_X)$의 대상은 $\kappa$-생성 준연접 가군의
여과 여극한, 즉 $\mathcal{F}_t$와 동형인 대상들의 여과
여극한이므로 $U$가 생성자임은 명백하다.
극한과 단사대상에 관한 주장들은 모든 그로텐디크 아벨 범주에서
성립한다. Injectives의 정리
\ref{injectives-theorem-injective-embedding-grothendieck}와
보조정리 \ref{injectives-lemma-grothendieck-products}를 보라.

\medskip\noindent
(2)의 증명. $Q$를 구성하기 위해 다음의 일반적인 절차를 사용한다.
대상 $\mathcal{F}$가 $\textit{Mod}(\mathcal{O}_X)$에 주어지면
함자
$$
\QCoh(\mathcal{O}_X)^{opp} \longrightarrow \textit{Sets},\quad
\mathcal{G} \longmapsto \Hom_X(\mathcal{G}, \mathcal{F})
$$
를 생각한다. 이 함자는 여극한을 극한으로 보내므로 표현가능하다.
Injectives의 보조정리
\ref{injectives-lemma-grothendieck-brown}을 보라. 따라서 준연접층
$Q(\mathcal{F})$와 함자적 동형사상
$\Hom_X(\mathcal{G}, \mathcal{F}) = \Hom_X(\mathcal{G}, Q(\mathcal{F}))$
이 존재하며, 이는 $\mathcal{G}$가 $\QCoh(\mathcal{O}_X)$ 안에
대하여 성립한다. 요네다 보조정리(Categories의 보조정리
\ref{categories-lemma-yoneda})에 의해 구성
$\mathcal{F} \leadsto Q(\mathcal{F})$는 $\mathcal{F}$에 관하여
함자적이다. 구성에 의해 $Q$는 포함 함자의 오른쪽 수반이다.
사상 $Q(\mathcal{F}) \to \mathcal{F}$가 $\mathcal{F}$가 준연접일 때
동형사상이라는 사실은 포함 함자
$\QCoh(\mathcal{O}_X) \to \textit{Mod}(\mathcal{O}_X)$가
충실충만이라는 사실의 형식적인 귀결이다.
\end{proof}

\section{닫힌 부분집합에 지지가 있는 절}
\label{properties-section-sections-with-support-in-closed}

\noindent
위상 공간 $X$, 닫힌 부분집합 $Z \subset X$, 아벨 층 $\mathcal{F}$가
주어지면 지지가 $Z$에 포함되는 절단들의 부분층을 취할 수 있다.
$X$가 스킴이고 $Z$가 닫힌 부분스킴이며 $\mathcal{F}$가 준연접
가군이면, 스킴 이론적 의미에서 $Z$에 지지를 갖는 절단을 취하는
변형도 있다. 그러나 스킴의 경우에는 조심해야 한다. 얻어지는
$\mathcal{O}_X$-가군이 준연접이 아닐 수 있기 때문이다.

\begin{lemma}
\label{properties-lemma-quasi-coherent-finite-type-ideals}
$X$를 준콤팩트 준분리 스킴이라 하자.
$U \subset X$를 열린 부분스킴이라 하자. 다음 조건들은 서로 동치이다.
\begin{enumerate}
\item $U$는 $X$에서 역준콤팩트이다.
\item $U$는 준콤팩트이다.
\item $U$는 유한 개의 아핀 열린집합의 합집합이다.
\item 유한형 준연접 아이디얼 층 $\mathcal{I} \subset \mathcal{O}_X$가
존재하여 $X \setminus U = V(\mathcal{I})$이다 (집합으로서).
\end{enumerate}
\end{lemma}

\begin{proof}
(1), (2), (3)의 동치는 보조정리
\ref{properties-lemma-quasi-separated-quasi-compact-open-retrocompact}에서 따른다.
(1), (2), (3)을 가정하자. $T = X \setminus U$라 놓자.
Schemes의 보조정리 \ref{schemes-lemma-reduced-closed-subscheme}에 의해
유일한 준연접 아이디얼 층 $\mathcal{J}$가 존재하여 $T$ 위에 유도되는
약화된 닫힌 부분스킴 구조를 잘라낸다. $\mathcal{J}|_U = \mathcal{O}_U$임에
유의하자. 이는 유한형 $\mathcal{O}_U$-가군이다.
보조정리 \ref{properties-lemma-extend}에 의해 유한형 준연접 부분층
$\mathcal{I} \subset \mathcal{J}$가 존재하며 $\mathcal{I}|_U = \mathcal{J}|_U$
이다. 따라서 $X \setminus U = V(\mathcal{I})$이고 (4)를 얻는다.
반대로 (4)의 $\mathcal{I}$를 택하고 $W = \Spec(R) \subset X$가 아핀 열린집합이면,
$\mathcal{I}|_W = \widetilde{I}$이며 어떤 유한 생성 아이디얼
$I \subset R$가 존재한다. 보조정리 \ref{properties-lemma-finite-type-module}을 보라.
따라서 $U \cap W = \Spec(R) \setminus V(I)$는 준콤팩트이다.
대수의 보조정리 \ref{algebra-lemma-qc-open}을 보라. 그러므로 보조정리
\ref{properties-lemma-retrocompact}에 의해 $U \subset X$는 역준콤팩트이다.
\end{proof}
\begin{lemma}
\label{properties-lemma-sections-annihilated-by-ideal}
$X$를 스킴이라 하자.
$\mathcal{I} \subset \mathcal{O}_X$를 준연접 아이디얼 층이라 하자.
$\mathcal{F}$를 준연접 $\mathcal{O}_X$-가군이라 하자.
$\mathcal{O}_X$-가군 층 $\mathcal{F}'$를 생각하자. 이는 각 열린집합
$U \subset X$에 다음을 대응시킨다.
$$
\mathcal{F}'(U)
=
\{s \in \mathcal{F}(U) \mid
\mathcal{I}s = 0\}
$$
$\mathcal{I}$가 유한형이라고 가정하자. 그러면
\begin{enumerate}
\item $\mathcal{F}'$는 준연접 $\mathcal{O}_X$-가군 층이다.
\item 임의의 아핀 열린집합 $U \subset X$에서
$\mathcal{F}'(U) = \{s \in \mathcal{F}(U) \mid \mathcal{I}(U)s = 0\}$이고,
\item $\mathcal{F}'_x = \{s \in \mathcal{F}_x \mid \mathcal{I}_x s = 0\}$이다.
\end{enumerate}
\end{lemma}

\begin{proof}
$\mathcal{F}'$를 정의하는 규칙이 $\mathcal{F}$의 부분층을 준다는 것은
명백하다(층 조건은 쉽게 확인할 수 있다). 따라서 다른 명제들을 확인하기
위해 $X$ 위에서 국소적으로 작업하면 된다. 즉
$X = \Spec(A)$, $\mathcal{F} = \widetilde{M}$,
$\mathcal{I} = \widetilde{I}$라고 가정해도 된다. 이 경우
$\mathcal{F}'(U) = \{x \in M \mid Ix = 0\} =: M'$임은 명백하다.
이는 $\widetilde{I}$가 전체 절단 $I$로 생성되기 때문이며, (2)를 증명한다.
$\mathcal{F}'$가 준연접임을 보이려면 모든 $f \in A$에 대해
$\{x \in M_f \mid I_f x = 0\} = (M')_f$임을 보이면 충분하다.
$I = (g_1, \ldots, g_t)$라고 쓰자. 이는 $\mathcal{I}$가 유한형이므로
가능하다. 보조정리 \ref{properties-lemma-finite-type-module}을 보라.
$x = y/f^n$이고 $I_fx = 0$이면, 이는 모든 $i$에 대해 어떤
$m \geq 0$가 존재하여 $f^mg_ix = 0$임을 뜻한다.
하나의 $m$을 택할 수 있으며, 이것은 모든 $i$에 대해 성립한다
(여기서 $I$가 유한 생성이라는 사실을 사용한다). 그러면
$f^mx \in M'$이고 $x/f^n = f^mx/f^{n + m}$은
$(M')_f$에서 성립하므로 원하는 결론을 얻는다.
(3)의 증명은 유사하므로 생략한다.
\end{proof}

\begin{definition}
\label{properties-definition-subsheaf-sections-annihilated-by-ideal}
$X$를 스킴이라 하자.
$\mathcal{I} \subset \mathcal{O}_X$를 유한형 준연접 아이디얼 층이라 하자.
$\mathcal{F}$를 준연접 $\mathcal{O}_X$-가군이라 하자.
위의 보조정리 \ref{properties-lemma-sections-annihilated-by-ideal}에서 정의한
부분층 $\mathcal{F}' \subset \mathcal{F}$를
{\it $\mathcal{I}$에 의해 소멸되는 절단의 부분층}이라고 부른다.
\end{definition}

\begin{lemma}
\label{properties-lemma-push-sections-annihilated-by-ideal}
$f : X \to Y$를 스킴의 준콤팩트 준분리 사상이라 하자.
$\mathcal{I} \subset \mathcal{O}_Y$를 유한형 준연접 아이디얼 층이라 하자.
$\mathcal{F}$를 준연접 $\mathcal{O}_X$-가군이라 하자.
$\mathcal{F}' \subset \mathcal{F}$를
$f^{-1}\mathcal{I}\mathcal{O}_X$에 의해 소멸되는 절단의 부분층이라 하자.
그러면 $f_*\mathcal{F}' \subset f_*\mathcal{F}$는
$\mathcal{I}$에 의해 소멸되는 절단의 부분층이다.
\end{lemma}

\begin{proof}
생략한다. (힌트: $f$가 준콤팩트 준분리라는 가정은
$f_*\mathcal{F}$가 준연접임을 뜻하므로, 보조정리
\ref{properties-lemma-sections-annihilated-by-ideal}를 $\mathcal{I}$와
$f_*\mathcal{F}$에 적용할 수 있다.)
\end{proof}

\noindent
위상 공간 위의 아벨 층에 대해서는 닫힌 부분집합에 지지가 있는 절단의
부분층을 이미 Modules의 주석 \ref{modules-remark-sections-support-in-closed}에서
논의했다. 준연접 가군의 경우 이 부분가군은 항상 준연접인 것은 아니지만,
닫힌 부분집합의 여집합이 역준콤팩트이면 준연접이다.
\begin{lemma}
\label{properties-lemma-sections-supported-on-closed-subset}
$X$를 스킴이라 하자. $Z \subset X$를 닫힌 부분집합이라 하자.
$\mathcal{F}$를 준연접 $\mathcal{O}_X$-가군이라 하자.
$\mathcal{O}_X$-가군 층 $\mathcal{F}'$를 생각하자. 이는 각 열린집합
$U \subset X$에 다음을 대응시킨다.
$$
\mathcal{F}'(U)
=
\{s \in \mathcal{F}(U) \mid
\text{the support of }s\text{ is contained in }Z \cap U\}
$$
$X \setminus Z$가 $X$의 역준콤팩트 열린집합이면 다음이 성립한다.
\begin{enumerate}
\item 아핀 열린집합 $U \subset X$에 대하여 유한 생성 아이디얼
$I \subset \mathcal{O}_X(U)$가 존재하여 $Z \cap U = V(I)$이다.
\item (1)의 $U$와 $I$에 대하여
$\mathcal{F}'(U) = \{x \in \mathcal{F}(U) \mid
I^nx = 0 \text{ for some } n\}$이다.
\item $\mathcal{F}'$는 준연접 $\mathcal{O}_X$-가군 층이다.
\end{enumerate}
\end{lemma}

\begin{proof}
(1)은 Algebra의 보조정리 \ref{algebra-lemma-qc-open}이다.
$U = \Spec(A)$라 놓고 $I$는 (1)에서와 같이 택하자.
그러면 $\mathcal{F}|_U$는 어떤 $A$-가군 $M$에 대응하는
준연접 층이다. 다음이 성립한다.
$$
\mathcal{F}'(U) = \{x \in M \mid x = 0\text{ in }M_\mathfrak p
\text{ for all }\mathfrak p \not \in Z\}.
$$
이는 Modules의 정의 \ref{modules-definition-support}에 의한 것이다.
따라서 $x \in \mathcal{F}'(U)$일 필요충분조건은
$V(\text{Ann}(x)) \subset V(I)$이다. Algebra의 보조정리
\ref{algebra-lemma-support-element}을 보라. $I$가 유한 생성이므로
이는 $I^n x = 0$가 어떤 $n$에 대해 성립하는 것과 동치이다.
이로써 (2)가 증명된다.

\medskip\noindent
(3)의 증명. 열린집합 $U \subset X$가 주어지면 정확한 열
$$
0 \to \mathcal{F}'(U) \to \mathcal{F}(U) \to \mathcal{F}(U \setminus Z)
$$
이 있음에 유의하자. 포함 사상
$j : X \setminus Z \to X$를 나타내면,
$\mathcal{F}(U \setminus Z)$는
$j_*(\mathcal{F}|_{X \setminus Z})$의 $U$ 위 절단임을 관찰할 수 있다.
따라서 정확한 열
$$
0 \to \mathcal{F}' \to \mathcal{F} \to j_*(\mathcal{F}|_{X \setminus Z})
$$
이 있다. 제한 $\mathcal{F}|_{X \setminus Z}$는 준연접이다.
그러므로 $j_*(\mathcal{F}|_{X \setminus Z})$는 준연접이다.
이는 Schemes의 보조정리 \ref{schemes-lemma-push-forward-quasi-coherent}와
$j$가 준콤팩트라는 가정(모든 열린 몰입은 분리되어 있다)에 따른다.
따라서 $\mathcal{F}'$는 준연접 가군 사이의 사상의 핵으로서 준연접이다.
Schemes의 절 \ref{schemes-section-quasi-coherent}을 보라.
\end{proof}

\begin{definition}
\label{properties-definition-subsheaf-sections-supported-on-closed}
$X$를 스킴이라 하자.
$T \subset X$를 닫힌 부분집합이라 하고 그 여집합이
$X$에서 역준콤팩트라고 하자.
$\mathcal{F}$를 준연접 $\mathcal{O}_X$-가군이라 하자.
보조정리 \ref{properties-lemma-sections-supported-on-closed-subset}에서 정의한
준연접 부분층 $\mathcal{F}' \subset \mathcal{F}$를
{\it $T$에 지지가 있는 절단의 부분층}이라고 부른다.
\end{definition}

\begin{lemma}
\label{properties-lemma-push-sections-supported-on-closed-subset}
$f : X \to Y$를 스킴의 준콤팩트 준분리 사상이라 하자.
$Z \subset Y$를 닫힌 부분집합이라 하고 $Y \setminus Z$가
$Y$에서 역준콤팩트라고 하자. $\mathcal{F}$를 준연접
$\mathcal{O}_X$-가군이라 하자. $\mathcal{F}' \subset \mathcal{F}$를
$f^{-1}Z$에 지지가 있는 절단의 부분층이라 하자.
그러면 $f_*\mathcal{F}' \subset f_*\mathcal{F}$는
$Z$에 지지가 있는 절단의 부분층이다.
\end{lemma}

\begin{proof}
생략한다. (힌트: 먼저 $X \setminus f^{-1}Z$가 $X$에서 역준콤팩트임을
보인다. 이는 $Y \setminus Z$가 $Y$에서 역준콤팩트이기 때문이다. 따라서
보조정리 \ref{properties-lemma-sections-supported-on-closed-subset}를
$f^{-1}Z$와 $\mathcal{F}$에 적용할 수 있다. $f$가 준콤팩트 준분리이므로
$f_*\mathcal{F}$는 준연접이다. 그러므로 보조정리
\ref{properties-lemma-sections-supported-on-closed-subset}를 $Z$와
$f_*\mathcal{F}$에 적용할 수 있다. 마지막으로 층들을 직접 대응시킨다.)
\end{proof}
\section{준연접 층의 절단}
\label{properties-section-sections-quasi-coherent}

\noindent
다음은 아핀 스펙트럼의 준콤팩트 열린집합 위에서 준연접 층의
절단을 계산한 것이다.

\begin{lemma}
\label{properties-lemma-sections-over-quasi-compact-open-in-affine}
$A$를 환이라 하자.
$I \subset A$를 유한 생성 아이디얼이라 하자.
$M$을 $A$-가군이라 하자.
그러면 다음의 표준 사상이 존재한다.
$$
\colim_n \Hom_A(I^n, M)
\longrightarrow
\Gamma(\Spec(A) \setminus V(I), \widetilde{M}).
$$
이 사상은 항상 단사이다.
모든 $x \in M$에 대하여 $Ix = 0 \Rightarrow x = 0$이면
이 사상은 동형사상이다. 일반적으로
$M_n = \{x \in M \mid I^nx = 0\}$이라 놓으면 다음 동형사상이 존재한다.
$$
\colim_n \Hom_A(I^n, M/M_n)
\longrightarrow
\Gamma(\Spec(A) \setminus V(I), \widetilde{M}).
$$
\end{lemma}

\begin{proof}
$I^{n + 1} \subset I^n$이고 $M_n \subset M_{n + 1}$이므로,
이 사상들을 통한 합성을 사용하여 $A$-가군의 표준 사상
$$
\Hom_A(I^n, M)
\longrightarrow
\Hom_A(I^{n + 1}, M)
$$
및
$$
\Hom_A(I^n, M/M_n)
\longrightarrow
\Hom_A(I^{n + 1}, M/M_{n + 1})
$$
을 얻고, 이를 계에서 전이 사상으로 사용한다.
$A$-가군 사상 $\varphi : I^n \to M$이 주어지면
층 사상 $\widetilde{\varphi} : \widetilde{I^n} \to \widetilde{M}$을 얻으며,
이를 열린집합 $\Spec(A) \setminus V(I)$에 제한할 수 있다.
$\widetilde{I^n}$을 이 열린집합에 제한하면 구조층이 되므로,
$\Gamma(\Spec(A) \setminus V(I), \widetilde{M})$의 원소를 얻는다.
이것이 계 $\Hom_A(I^n, M)$의 전이 사상과 양립한다는 확인은 생략한다.
이로써 첫 번째 화살표를 얻는다.
두 번째 화살표를 얻기 위해
$\widetilde{M}$과 $\widetilde{M/M_n}$은 열린집합
$\Spec(A) \setminus V(I)$ 위에서 일치한다는 점을 관찰한다.
이는 층 $\widetilde{M_n}$의 지지가 명백히 $V(I)$에 포함되기 때문이다.
따라서 앞과 같은 장치를 사용할 수 있다.

\medskip\noindent
이 화살표를 대수적으로 정의하는 방법을 다음에 계산한다.
$I = (f_1, \ldots, f_t)$라 하자. 그러면
$\Spec(A) \setminus V(I) = \bigcup_{i = 1, \ldots, t} D(f_i)$이다.
따라서
$$
0 \to
\Gamma(\Spec(A) \setminus V(I), \widetilde{M}) \to
\bigoplus\nolimits_i M_{f_i} \to
\bigoplus\nolimits_{i, j} M_{f_if_j}
$$
은 정확하다. $\varphi : I^n \to M$이 $A$-가군 사상이라고 하자.
$\varphi(f_i^n)/f_i^n \in M_{f_i}$인 원소들의 벡터를 생각하자.
이 벡터가 위 정확한 열의 두 번째 직합에서 영사상으로 간다는 것은
쉽게 알 수 있다. 따라서
$\Gamma(\Spec(A) \setminus V(I), \widetilde{M})$의 원소를 얻는다.
이 설명이 앞에서 주어진 설명과 일치한다는 확인은 생략한다.

\medskip\noindent
이 설명을 사용하여 첫 번째 화살표가 단사임을 보이자.
즉 $\varphi$가 영사상으로 간다면, 각 $i$에 대해
$\varphi(f_i^n)/f_i^n$는 $M_{f_i}$에서 영이다.
다시 말해 각 $i$에 대해 $f_i^m\varphi(f_i^n) = 0$인 어떤
$m \geq 0$가 존재한다.
하나의 $m$을 택할 수 있으며, 이는 모든 $i$에 대해 성립한다.
그러면 $\varphi(f_i^{n + m}) = 0$가 모든 $i$에 대해 성립함을 알 수 있다.
이는 $\varphi|_{I^{t(n + m - 1) + 1}} = 0$가 됨을 쉽게 알 수 있으며,
다시 말해 $\varphi$가 여극한의 $t(n + m - 1) + 1$번째 항에서
영사상으로 간다는 뜻이다. 따라서 단사성이 따른다.
\medskip\noindent
각 $M_n = 0$임에 유의하자. 이는
$Ix = 0 \Rightarrow x = 0$가 $x \in M$에 대해 성립하는 경우이다.
따라서 이 보조정리의 증명을
끝내려면 두 번째 화살표가 동형사상임을 보이면 충분하다.

\medskip\noindent
이 보조정리의 두 번째 사상의 역을 구성해 보자.
$s \in \Gamma(\Spec(A) \setminus V(I), \widetilde{M})$라 하자.
이는 위의 정확한 열의 첫 번째 직합에서 $x_i/f_i^n$라는 벡터에
대응하며, 여기서 $x_i \in M$이다.
따라서 각 $i, j$에 대하여 어떤 $m \geq 0$가 존재하여
$f_i^m f_j^m (f_j^n x_i - f_i^n x_j) = 0$가 $M$에서 성립한다.
하나의 $m$을 택할 수 있으며, 이는 모든 쌍 $i, j$에 대해 성립한다.
$x_i$를 $f_i^mx_i$로, $n$을 $n + m$으로 바꾸면
$f_j^nx_i = f_i^nx_j$가 $M$에서 모든 $i, j$에 대해 성립함을 알 수 있다.
다음을 정의하자.
$$
K_n = \{x \in M \mid f_1^nx = \ldots = f_t^nx = 0\}
$$
다음과 같은 $A$-가군 사상이 존재한다고 주장한다.
$$
\varphi :
I^{t(n - 1) + 1}
\longrightarrow
M/K_n
$$
이는 합
$f_1^{e_1} \ldots f_t^{e_t}$를 $\sum e_i = t(n - 1) + 1$일 때
$K_n$으로 나눈 다음 표현식의 동치류로 보낸다.
$f_1^{e_1} \ldots f_i^{e_i - n} \ldots f_t^{e_t}x_i$
여기서 $i$는 $e_i \geq n$이 되도록 택한다
(그러한 $i$가 적어도 하나 존재한다).
이것이 실제로 잘 정의됨을 보이기 위해 다음을 가정하자.
$$
\sum\nolimits_{E = (e_1, \ldots, e_t), |E| = t(n - 1) + 1}
a_E f_1^{e_1} \ldots f_t^{e_t} = 0
$$
이는 계수 $a_E$가 $A$에 있는 단항식들 사이의 관계이다.
그렇다면 이를 다음 원소로 보낼 것이다.
$$
z =
\sum\nolimits_{E = (e_1, \ldots, e_t), |E| = t(n - 1) + 1}
a_E f_1^{e_1} \ldots f_{i(E)}^{e_{i(E)} - n} \ldots f_t^{e_t}x_{i(E)}
$$
각 다중지수 $E$에 대해 특정한 $i(E)$를 택했으며,
$e_{i(E)} \geq n$이다.
이를 $f_j^n$으로 곱하면 임의의 $j$에 대해 영이 됨에 유의하자.
위에서 얻은 관계 $f_j^nx_i = f_i^nx_j$에 의해 다음과 같이 된다.
\begin{align*}
f_j^nz & = \sum\nolimits_{E = (e_1, \ldots, e_t), |E| = t(n - 1) + 1}
a_E f_1^{e_1} \ldots f_j^{e_j + n}
\ldots f_{i(E)}^{e_{i(E)} - n} \ldots f_t^{e_t}x_{i(E)} \\
& =
\sum\nolimits_{E = (e_1, \ldots, e_t), |E| = t(n - 1) + 1}
a_E f_1^{e_1} \ldots f_t^{e_t}x_j
= 0.
\end{align*}
따라서 $z \in K_n$이고 모든 관계가 $M/K_n$에서 영으로 보내짐을
알 수 있다. 이것이 주장을 증명한다.

\medskip\noindent
$K_n \subset M_{t(n - 1) + 1}$임에 유의하자. 따라서 사상 $\varphi$는
특히 $A$-가군 사상
$I^{t(n - 1) + 1} \to M/M_{t(n - 1) + 1}$을 준다.
이로써 보조정리의 두 번째 화살표가 전사임이 증명된다.
단사성의 증명은 생략한다.
\end{proof}
\begin{example}
\label{properties-example-does-not-work-in-general}
첫 번째로 표시된 사상이 동형사상이 아님을 보이는 두 예를
보이겠다. 보조정리 \ref{properties-lemma-sections-over-quasi-compact-open-in-affine}를 보라.

\medskip\noindent
$k$를 체라 하자. 다음 환을 생각하자.
$$
A = k[x, y, z_1, z_2, \ldots]/(x^nz_n).
$$
$I = (x)$라 놓고 $M = A$라 하자. 그러면 원소 $y/x$는
$\Spec(A)$의 구조층에서 $D(x) = \Spec(A) \setminus V(I)$ 위의
절단을 정의한다.
$y/x$가 표준 사상
$\colim \Hom_A(I^n, A) \to A_x = \mathcal{O}(D(x))$의 상에
속하지 않는다고 주장한다. 만일 그렇다면 준동형
$\varphi : I^n \to A$로부터 어떤 $n$에 대해 왔을 것이다.
$a = \varphi(x^n)$이라 놓자. 그러면
$x^m(xa - x^ny) = 0$인 어떤 $m > 0$가 존재할 것이다. 이는
$x^{m + 1}a = x^{m + n}y$임을 뜻한다. 다시 이는
$\varphi(x^{n + m + 1}) = x^{m + n}y$임을 뜻한다.
따라서 다음이 성립해야 하므로 모순에 이른다.
$$
0 = \varphi(0) = \varphi(z_{n + m + 1} x^{n + m + 1}) =
x^{m + n}y z_{n + m + 1}
$$
이는 환 $A$에서 참이 아니다.

\medskip\noindent
$k$를 체라 하자. 다음 환을 생각하자.
$$
A = k[f, g, x, y, \{a_n, b_n\}_{n \geq 1}]/
(fy - gx, \{a_nf^n + b_ng^n\}_{n \geq 1}).
$$
$I = (f, g)$라 놓고 $M = A$라 하자. 그러면
$x/f \in A_f$와 $y/g \in A_g$는 $A_{fg}$에서 같은 원소로 간다.
따라서 이들은 $s$라는 절단을 정의하며, $\Spec(A)$의 구조층에서
$D(f) \cup D(g) = \Spec(A) \setminus V(I)$ 위의 절단이다. 그러나
첫 번째 표시된 화살표의 원천에서 오는 어떤 $n \geq 0$에 대해서도
$s$가 $A$-가군 사상 $\varphi : I^n \to A$로부터 유도되지는 않는다.
보조정리
\ref{properties-lemma-sections-over-quasi-compact-open-in-affine}를 보라.
구체적으로 그러한 가군 사상이 주어졌다고 하고
$x_n = \varphi(f^n)$ 및 $y_n = \varphi(g^n)$이라 놓자.
그러면 $f^mx_n = f^{n + m - 1}x$ 및
$g^my_n = g^{n + m - 1}y$가 어떤 $m \geq 0$에 대해 성립한다
(보조정리의 증명 참조). 그런데 그러면
다음이 성립한다.
$0 = \varphi(0) =
\varphi(a_{n + m}f^{n + m} + b_{n + m}g^{n + m}) =
a_{n + m}f^{n + m - 1}x + b_{n + m}g^{n + m - 1}y$는 성립하지 않는다.
이는 환 $A$에서 성립하지 않는다.
\end{example}

\noindent
다음 보조정리는 뇌터 환의 경우에 개선할 것이다. Cohomology of Schemes의
보조정리 \ref{coherent-lemma-homs-over-open}을 보라.

\begin{lemma}
\label{properties-lemma-sections-over-quasi-compact-open}
$X$를 준콤팩트 스킴이라 하자.
$\mathcal{I} \subset \mathcal{O}_X$를 유한형 준연접 아이디얼 층이라 하자.
$Z \subset X$를 $\mathcal{I}$로 정의되는 닫힌 부분스킴이라 하고
$U = X \setminus Z$라 놓자.
$\mathcal{F}$를 준연접 $\mathcal{O}_X$-가군이라 하자.
표준 사상
$$
\colim_n \Hom_{\mathcal{O}_X}(\mathcal{I}^n,
\mathcal{F})
\longrightarrow
\Gamma(U, \mathcal{F})
$$
은 단사이다. 더 나아가 $X$가 준분리라고 가정하자.
$\mathcal{F}_n \subset \mathcal{F}$를 $\mathcal{I}^n$에 의해 소멸되는
절단들의 부분층이라 하자. 표준 사상
$$
\colim_n \Hom_{\mathcal{O}_X}(\mathcal{I}^n,
\mathcal{F}/\mathcal{F}_n)
\longrightarrow
\Gamma(U, \mathcal{F})
$$
은 동형사상이다.
\end{lemma}
\begin{proof}
$\Spec(A) = W \subset X$를 아핀 열린집합이라 하자.
$\mathcal{F}|_W = \widetilde{M}$이라 쓰자. 여기서 어떤
$A$-가군 $M$이다. 또한 $\mathcal{I}|_W = \widetilde{I}$라 쓰자.
여기서 $I \subset A$는 어떤 유한형 아이디얼이다. 보조정리의 첫 번째 표시 사상을
$W$에 제한하면 보조정리
\ref{properties-lemma-sections-over-quasi-compact-open-in-affine}의 첫 번째
표시 사상을 얻는다. $X$를 유한 개의 아핀 열린집합으로 덮을 수 있으므로
이로부터 보조정리의 첫 번째 표시 사상이 단사임이 따른다.

\medskip\noindent
$\mathcal{F}_n|_W = \widetilde{M_n}$임을 알 수 있으며,
$M_n \subset M$은 보조정리
\ref{properties-lemma-sections-over-quasi-compact-open-in-affine}에서 정의한
것이다(세부사항은 생략한다). 이 보조정리에 의해 다음 전단사가 존재한다.
$$
\colim_n  \Hom_{\mathcal{O}_W}(
\mathcal{I}^n|_W, (\mathcal{F}/\mathcal{F}_n)|_W)
\longrightarrow
\Gamma(U \cap W, \mathcal{F})
$$
이는 그러한 아핀 열린집합 $W$마다 성립한다.

\medskip\noindent
보조정리의 두 번째 표시 화살표가 전단사임을 보이기 위해
유한 아핀 열린 덮개
$X = \bigcup_{j = 1, \ldots, m} W_j$를 택하자.
단사성은 위의 결과와 덮개의 유한성에서 즉시 따른다.
$X$가 준분리이면 각 쌍 $j, j'$에 대해 유한 아핀 열린 덮개
$$
W_j \cap W_{j'} = \bigcup\nolimits_{k = 1, \ldots, m_{jj'}} W_{jj'k}.
$$
를 택하자.
$s \in \Gamma(U, \mathcal{F})$라 하자. 위에서 본 것처럼 각 $j$에 대해
어떤 $n_j$와 사상
$\varphi_j : \mathcal{I}^{n_j}|_{W_j} \to
(\mathcal{F}/\mathcal{F}_{n_j})|_{W_j}$가 존재하며,
이는 $s|_{U \cap W_j}$에 대응한다.
같은 방식으로 각 세쌍 $(j, j', k)$에 대해 어떤 정수 $n_{jj'k}$가
존재하여 $\varphi_j$와 $\varphi_{j'}$를
$\mathcal{I}^{n_{jj'k}} \to \mathcal{F}/\mathcal{F}_{n_{jj'k}}$라는
사상으로 제한한 것들이 $W_{jj'k}$ 위에서 일치한다.
$n = \max\{n_j, n_{jj'k}\}$라 놓으면 $\varphi_j$들은
$\mathcal{I}^n \to \mathcal{F}/\mathcal{F}_n$이라는 사상으로
$X$ 위에서 붙는다. 이로써 사상의 전사성이 증명된다.
\end{proof}
\section{충분한 가역층}
\label{properties-section-ample}

\noindent
Modules의 보조정리 \ref{modules-lemma-s-open}를 상기하자.
가역층 $\mathcal{L}$이 국소 환 달린 공간 $X$ 위에 주어지고,
그 전역 절단 $s$가 $\mathcal{L}$에 주어지면
집합 $X_s = \{x \in X \mid s \not \in \mathfrak m_x\mathcal{L}_x\}$는
열려 있다. 일반적으로
$X_s \cap X_{s'} = X_{ss'}$임을 관찰할 수 있다. 여기서 $ss'$는
절단 $s \otimes s' \in \Gamma(X, \mathcal{L} \otimes \mathcal{L}')$을
뜻한다.

\begin{definition}
\label{properties-definition-ample}
\begin{reference}
\cite[II Definition 4.5.3]{EGA}
\end{reference}
$X$를 스킴이라 하자.
$\mathcal{L}$을 가역 $\mathcal{O}_X$-가군이라 하자.
$\mathcal{L}$이 {\it 충분하다}는 것은 다음을 뜻한다.
\begin{enumerate}
\item $X$가 준콤팩트이고,
\item 모든 $x \in X$에 대해 어떤 $n \geq 1$과
$s \in \Gamma(X, \mathcal{L}^{\otimes n})$가 존재하여
$x \in X_s$이고 $X_s$가 아핀이다.
\end{enumerate}
\end{definition}

\begin{lemma}
\label{properties-lemma-ample-power-ample}
\begin{reference}
\cite[II Proposition 4.5.6(i)]{EGA}
\end{reference}
$X$를 스킴이라 하자. $\mathcal{L}$을 가역 $\mathcal{O}_X$-가군이라 하자.
$n \geq 1$이라 하자. 그러면 $\mathcal{L}$이 충분할 필요충분조건은
$\mathcal{L}^{\otimes n}$이 충분한 것이다.
\end{lemma}

\begin{proof}
이는 $X_{s^n} = X_s$라는 사실에서 따른다.
\end{proof}

\begin{lemma}
\label{properties-lemma-ample-on-closed}
$X$를 스킴이라 하자.
$\mathcal{L}$을 충분한 가역 $\mathcal{O}_X$-가군이라 하자.
임의의 닫힌 부분스킴 $Z \subset X$에 대해
$\mathcal{L}$의 $Z$로의 제한은 충분하다.
\end{lemma}

\begin{proof}
이는 준콤팩트 공간의 닫힌 부분집합이 준콤팩트이고,
아핀 스킴의 닫힌 부분스킴이 아핀이라는 사실에서 명백하다
(스킴, 보조정리 \ref{schemes-lemma-closed-immersion-affine-case} 참조).
\end{proof}

\begin{lemma}
\label{properties-lemma-affine-cap-s-open}
$X$를 스킴이라 하자. $\mathcal{L}$을 가역 $\mathcal{O}_X$-가군이라 하자.
$s \in \Gamma(X, \mathcal{L})$라 하자. 임의의 아핀 $U \subset X$에 대해
교집합 $U \cap X_s$는 아핀이다.
\end{lemma}

\begin{proof}
이는 다음 대수 문제로 옮겨진다.
$R$을 환이라 하자. $N$을 가역 $R$-가군
(즉, 계수 1의 국소 자유 가군)이라 하자.
$s \in N$을 원소라 하자.
그러면 $V = \{\mathfrak p \mid s \not \in \mathfrak p N\}$는
$\Spec(R)$의 아핀 열린 부분집합이다.

\medskip\noindent
$A = \bigoplus_{n \geq 0} A_n$을 $N$의 대칭 대수
(가환 대수이다)라 하고 $s$를 $A_1$의 원소로 보자.
$B = A/(s - 1)A$로 두자. 이는 $R$-대수이며 그 구성이 임의의
기저변환 $R \to R'$와 가환한다. 따라서
$B' = B \otimes_R R'$는 $s$가
$N' = N \otimes_R R'$에서 영으로 사상되면 영환이다.
그러므로 $x \in \Spec(R) \setminus V$이면
$B \otimes_R \kappa(x) = 0$이다.
$\Spec(B) \to \Spec(R)$는
$V$를 통해 인수분해된다. 실제로 $x \not \in V$ 위의
섬유는 비어 있기 때문이다. 한편
$\Spec(R') \subset V$가 아핀 열린집합이면 $s$는 $N'$의
기저 원소로 사상되며, 따라서
$B' = R'[s]/(s - 1) \cong R'$임을 알 수 있다.
그러므로 $\Spec(B) \to V$는 동형사상이고 $V$는
실제로 아핀이다.
\end{proof}

\begin{lemma}
\label{properties-lemma-ample-tensor-globally-generated}
\begin{reference}
\cite[II Proposition 4.5.6(ii)]{EGA}
\end{reference}
$X$를 스킴이라 하자. $\mathcal{L}$과 $\mathcal{M}$을
가역 $\mathcal{O}_X$-가군이라 하자. 다음 조건을 만족한다고 하자.
\begin{enumerate}
\item $\mathcal{L}$은 충분하고,
\item 열린집합 $X_t$들이 (여기서
$t \in \Gamma(X, \mathcal{M}^{\otimes m})$이고 $m > 0$)
$X$를 덮는다.
\end{enumerate}
그러면 $\mathcal{L} \otimes \mathcal{M}$은 충분하다.
\end{lemma}

\begin{proof}
정의 \ref{properties-definition-ample}의 조건들을 확인하자.
$\mathcal{L}$이 충분하므로 $X$는 준콤팩트이다.
$x \in X$를 잡자. 어떤 $n \geq 1$, $m \geq 1$,
$s \in \Gamma(X, \mathcal{L}^{\otimes n})$, 그리고
$t \in \Gamma(X, \mathcal{M}^{\otimes m})$가 존재하여
$x \in X_s$, $x \in X_t$이고 $X_s$가 아핀이 되도록 택하자.
그러면
$s^mt^n \in \Gamma(X, (\mathcal{L} \otimes \mathcal{M})^{\otimes nm})$이고,
$x \in X_{s^mt^n}$이다. 또한 보조정리
\ref{properties-lemma-affine-cap-s-open}에 의해 $X_{s^mt^n}$은 아핀이다.
\end{proof}
\begin{lemma}
\label{properties-lemma-affine-s-opens}
$X$를 스킴이라 하자. $\mathcal{L}$을 가역 $\mathcal{O}_X$-가군이라 하자.
열린집합 $X_s$들을 생각하자. 여기서
$s \in \Gamma(X, \mathcal{L}^{\otimes n})$이고
$n \geq 1$이다. 이 열린집합들이 $X$의 위상에 대한 기저를
이룬다고 가정하자. 이 열린집합들 중 아핀인 $X_s$들이
$X$의 위상에 대한 기저를 이룬다.
\end{lemma}

\begin{proof}
$x \in X$를 잡자. 아핀 열린 근방
$\Spec(R) = U \subset X$를 $x$의 근방으로 택하자.
가정에 의해 어떤
$n \geq 1$과 $s \in \Gamma(X, \mathcal{L}^{\otimes n})$이 존재하여
$X_s \subset U$이다. 위의 보조정리
\ref{properties-lemma-affine-cap-s-open}에 의해 교집합
$X_s = U \cap X_s$는 아핀이다.
$U$를 임의로 작게 택할 수 있으므로 결론을 얻는다.
\end{proof}

\begin{lemma}
\label{properties-lemma-affine-s-opens-cover-quasi-separated}
$X$를 스킴이라 하고 $\mathcal{L}$을 가역 $\mathcal{O}_X$-가군이라 하자.
모든 점 $x$에 대해, $X$에서 어떤 $n \geq 1$과
$s \in \Gamma(X, \mathcal{L}^{\otimes n})$가 존재하여
$x \in X_s$이고 $X_s$가 아핀이라고 가정하자. 그러면 $X$는 분리된다.
\end{lemma}

\begin{proof}
먼저 $X$가 준분리임을 보이자. 가정에 의해
$X$를 $X_s$ 꼴의 아핀 열린집합들로 덮을 수 있다. 위의 보조정리
\ref{properties-lemma-affine-cap-s-open}에 의해
이러한 두 집합의 교집합은 아핀이다. 따라서 스킴의 보조정리
\ref{schemes-lemma-characterize-quasi-separated}에 의해 $X$는
준분리이다.

\medskip\noindent
$X$가 분리임을 보이기 위해 스킴의 보조정리
\ref{schemes-lemma-valuative-criterion-separatedness}의
값매김 판정법을 사용할 수 있다.
따라서 값매김환 $A$와 그 분수체 $K$ 및
두 사상 $f, g : \Spec(A) \to X$를 잡되,
두 합성
$\Spec(K) \to \Spec(A) \to X$가 일치한다고 하자.
$A$가 국소이므로 어떤 $p, q \ge 1$, $s \in \Gamma(X,
\mathcal{L}^{\otimes p})$, 그리고 $t \in \Gamma(X,
\mathcal{L}^{\otimes q})$가 존재하여
$X_s$와 $X_t$가 아핀이고, $f(\Spec A) \subseteq X_s$이며,
$g(\Spec A)
\subseteq X_t$이다. 이제 $s$를 $s^q$로,
$t$를 $t^p$로, $\mathcal{L}$을 $\mathcal{L}^{\otimes pq}$로
바꾼다. 이는 $X_s = X_{s^q}$ 및 $X_t = X_{t^p}$이므로
문제가 없으며, 이제 $s$와 $t$는 같은 층 $\mathcal{L}$의
절단이다.

\medskip\noindent
준연접 가군 $f^*\mathcal{L}$은 $A$-가군 $M$에 대응하고,
$g^*\mathcal{L}$은 아핀 스킴 위의 준연접 가군 분류에 의해
$A$-가군 $N$에 대응한다
(스킴, 보조정리 \ref{schemes-lemma-quasi-coherent-affine}).
$A$-가군 $M$과 $N$은 계수 $1$의 국소 자유 가군이다
(보조정리 \ref{properties-lemma-locally-free-module}). 또한 $A$가
국소이므로 이들은 자유 가군이다(대수, 보조정리
\ref{algebra-lemma-K0-local}). 따라서
$M$과 $N$을 각각 $A$-부분가군인
$M \otimes_A K$와 $N
\otimes_A K$로 식별할 수 있다.
$f|_{\Spec(K)} = g|_{\Spec(K)}$라는 등식은 동형사상
$\phi \colon M \otimes_A K \to N
\otimes_A K$를 결정한다.

\medskip\noindent
당김에 대응하는 원소를 각각 $x \in M$, $y \in N$이라 하자.
이는 $s$의 당김이며 $f$와 $g$에 따른 것이다. 이들은
$\phi(x \otimes 1) = y \otimes 1$을 만족한다. $f$의 상은
$X_s$에 포함되므로 $x \not\in \mathfrak{m}_A M$이다.
즉 $x$는 $M$을 생성한다. 따라서 $\phi$는
$M$을 $N$의 부분가군으로 보내는 동형사상을 결정하며,
$y$가 그 부분가군을 생성한다.
$t$를 사용하여 대칭적으로 논하면 $\phi^{-1}$은
$N$을 $M$의 부분가군으로 보내는 동형사상을 결정한다.
결과적으로 $\phi$는 $M$과 $N$의 동형사상으로 제한된다.
$x$가 $M$을 생성하고 그 상 $y$가 $N$을 생성하므로
$y \not\in \mathfrak{m}_A N$이다. 그러므로
$g(\Spec(A)) \subseteq X_s$이다. $X_s$는 아핀이므로
스킴의 보조정리 \ref{schemes-lemma-affine-separated}에 의해
분리되고, 이에 따라 $f = g$이다.
\end{proof}

\begin{lemma}
\label{properties-lemma-ample-separated}
$X$를 스킴이라 하자. $X$ 위에 충분한 가역층이 존재하면
$X$는 분리된다.
\end{lemma}

\begin{proof}
이는 보조정리 \ref{properties-lemma-affine-s-opens-cover-quasi-separated}와
정의 \ref{properties-definition-ample}에서 즉시 따른다.
\end{proof}

\begin{lemma}
\label{properties-lemma-map-into-proj}
$X$를 스킴이라 하자.
$\mathcal{L}$을 가역 $\mathcal{O}_X$-가군이라 하자.
$S = \Gamma_*(X, \mathcal{L})$을 등급환으로 두자.
$X$의 모든 점이 열린 부분스킴 $X_s$ 중 하나에 포함된다고 하자.
여기서 $s \in S_{+}$는 동차 원소이다. 그러면 표준적인
스킴 사상
$$
f : X \longrightarrow Y = \text{Proj}(S),
$$
이 존재하며, 이는 $S$의 동차 스펙트럼으로 가는 사상이다
(구성, 절 \ref{constructions-section-proj} 참조).
이 사상은 다음 성질을 갖는다.
\begin{enumerate}
\item $f^{-1}(D_{+}(s)) = X_s$ (모든 동차 $s \in S_{+}$에 대해),
\item 곱셈 사상과 양립하는 $\mathcal{O}_X$-가군 사상
$f^*\mathcal{O}_Y(n) \to \mathcal{L}^{\otimes n}$이 존재한다.
구성, 식 (\ref{constructions-equation-multiply}) 참조,
\item
$S_n \to \Gamma(Y, \mathcal{O}_Y(n)) \to \Gamma(X, \mathcal{L}^{\otimes n})$
의 합성은 항등사상이고,
\item 모든 $x \in X$에 대해 어떤 정수 $d \geq 1$과
열린 근방 $U \subset X$가 $x$의 근방으로 존재하여
$f^*\mathcal{O}_Y(dn)|_U \to \mathcal{L}^{\otimes dn}|_U$가
모든 $n \in \mathbf{Z}$에 대해 동형사상이다.
\end{enumerate}
\end{lemma}

\begin{proof}
$\psi : S \to \Gamma_*(X, \mathcal{L})$를 항등사상이라 하자.
구성의 보조정리
\ref{constructions-lemma-invertible-map-into-proj}의
세 쌍 $(U(\psi), r_{\mathcal{L}, \psi}, \theta)$를 사용할 것이다.
가정에 의해 열린 부분스킴 $U(\psi)$는 $X$와 같다.
따라서 $r_{\mathcal{L}, \psi} : U(\psi) \to Y$는 $X$ 전체에서
정의된다. $f = r_{\mathcal{L}, \psi}$로 둔다.
(2)의 사상들은 $\theta$의 성분들이다.
(3)은 위에서 인용한 보조정리의 조건 (2)에서 따른다.
(1)은 (3)과 그 보조정리의 조건 (1)을 결합하면 따른다.
(4)는 구성의 보조정리
\ref{constructions-lemma-invertible-map-into-proj}의
마지막 서술에서 따른다. 그곳에서 언급한 사상 $\alpha$가
동형사상이기 때문이다.
\end{proof}

\begin{lemma}
\label{properties-lemma-map-into-proj-quasi-compact}
$X$를 스킴이라 하자. $\mathcal{L}$을 가역 $\mathcal{O}_X$-가군이라 하자.
$S = \Gamma_*(X, \mathcal{L})$로 두자.
(a) $X$의 모든 점이 열린 부분스킴 $X_s$ 중 하나에 포함되고,
여기서 $s \in S_{+}$는 동차 원소이며, (b) $X$가 준콤팩트라고 가정하자.
그러면 위의 보조정리 \ref{properties-lemma-map-into-proj}의 표준 사상
$f : X \longrightarrow \text{Proj}(S)$는 준콤팩트이고 조밀한 상을 갖는다.
\end{lemma}

\begin{proof}
$f$가 준콤팩트임을 보이려면
$f^{-1}(D_{+}(s))$가 준콤팩트임을 모든 동차 $s \in S_{+}$에 대해
보이면 충분하다.
$X = \bigcup_{i = 1, \ldots, n} X_i$를 아핀 열린집합들의 유한 합으로 쓰자.
보조정리 \ref{properties-lemma-affine-cap-s-open}에 의해 각 교집합
$X_s \cap X_i$는 아핀이다. 따라서
$X_s = \bigcup_{i = 1, \ldots, n} X_s \cap X_i$는 준콤팩트이다.
$f$의 상이 조밀하지 않다고 가정하여 모순을 얻자. 그러면
열린집합 $D_+(s)$들이, 여기서 $s \in S_+$는 동차 원소이며,
$\text{Proj}(S)$의 위상에 대한 기저를 이루므로,
$s$를 택하여 $D_+(s) \not = \emptyset$이고
$f(X) \cap D_+(s) = \emptyset$가 되게 할 수 있다.
보조정리 \ref{properties-lemma-map-into-proj}에 의해
이는 $X_s = \emptyset$임을 뜻한다. 보조정리
\ref{properties-lemma-invert-s-sections}에 의해 이는 어떤 $s^n$이
$\mathcal{L}^{\otimes n\deg(s)}$의 영절단임을 뜻한다.
그러면 $D_+(s) = \emptyset$이 되어 원하는 모순을 얻는다.
\end{proof}

\begin{lemma}
\label{properties-lemma-ample-immersion-into-proj}
$X$를 스킴이라 하자. $\mathcal{L}$을 가역 $\mathcal{O}_X$-가군이라 하자.
$S = \Gamma_*(X, \mathcal{L})$로 두자.
$\mathcal{L}$이 충분하다고 가정하자. 그러면 보조정리
\ref{properties-lemma-map-into-proj}의 표준적인 스킴 사상
$f : X \longrightarrow \text{Proj}(S)$은
조밀한 상을 갖는 열린 몰입이다.
\end{lemma}

\begin{proof}
보조정리 \ref{properties-lemma-affine-s-opens-cover-quasi-separated}에 의해
$X$가 준분리임을 알 수 있다. 다음과 같은 유한 개의
동차 원소 $s_1, \ldots, s_n \in S_{+}$를 택하여
$X_{s_i}$들이 아핀이고 $X = \bigcup X_{s_i}$가 되게 하자.
$s_i$의 차수를 $d_i$라 하자. $D_{+}(s_i)$의 역상은
$f$에 의해 $X_{s_i}$이다. 보조정리
\ref{properties-lemma-map-into-proj}를 참조하라.
보조정리 \ref{properties-lemma-invert-s-sections}에 의해 환 사상
$$
(S^{(d_i)})_{(s_i)} = \Gamma(D_{+}(s_i), \mathcal{O}_{\text{Proj}(S)})
\longrightarrow
\Gamma(X_{s_i}, \mathcal{O}_X)
$$
은 동형사상이다. 따라서 $f$는 동형사상
$X_{s_i} \to D_{+}(s_i)$를 유도한다. 그러므로 $f$는
$X$를 열린 부분스킴
$\bigcup_{i = 1, \ldots, n} D_{+}(s_i)$ 위로 보내는 동형사상이며,
이는 $\text{Proj}(S)$의 부분스킴이다.
상은 보조정리 \ref{properties-lemma-map-into-proj-quasi-compact}에 의해 조밀하다.
\end{proof}

\begin{lemma}
\label{properties-lemma-open-in-proj-ample}
$X$를 스킴이라 하자.
$S$를 등급환이라 하자. $X$가 준콤팩트이고,
다음과 같은 열린 몰입이 존재한다고 가정하자.
$$
j : X \longrightarrow Y = \text{Proj}(S).
$$
그러면 $j^*\mathcal{O}_Y(d)$는 어떤 $d > 0$에 대해
가역이고 충분한 층이다.
\end{lemma}

\begin{proof}
이는 구성, 보조정리 \ref{constructions-lemma-ample-on-proj}이다.
\end{proof}

\begin{proposition}
\label{properties-proposition-characterize-ample}
$X$를 준콤팩트 스킴이라 하자.
$\mathcal{L}$을 $X$ 위의 가역층이라 하자.
$S = \Gamma_*(X, \mathcal{L})$로 두자.
다음 조건들은 서로 동치이다.
\begin{enumerate}
\item
\label{properties-item-ample}
$\mathcal{L}$은 충분하다,
\item
\label{properties-item-immersion}
열린집합 $X_s$들이 (동차 원소 $s \in S_{+}$에 대해)
$X$를 덮고, 연관된 사상 $X \to \text{Proj}(S)$가
열린 몰입이다,
\item
\label{properties-item-s-basis}
열린집합 $X_s$들이 (동차 원소 $s \in S_{+}$에 대해)
$X$의 위상에 대한 기저를 이룬다,
\item
\label{properties-item-s-affine-basis}
열린집합 $X_s$들 중 아핀인 것들이
(동차 원소 $s \in S_{+}$에 대해) $X$의 위상에 대한 기저를 이룬다,
\item
\label{properties-item-qc-gg}
모든 준연접층 $\mathcal{F}$가 $X$ 위에 있을 때 다음 표준 사상들의
상들의 합
$$
\Gamma(X, \mathcal{F} \otimes_{\mathcal{O}_X} \mathcal{L}^{\otimes n})
\otimes_{\mathbf{Z}} \mathcal{L}^{\otimes -n}
\longrightarrow
\mathcal{F}
$$
($n \geq 1$인 경우)이 $\mathcal{F}$와 같다,
\item
\label{properties-item-qc-i-gg}
(\ref{properties-item-qc-gg})와 같은 성질이 $\mathcal{F}$가
모든 준연접 아이디얼층을 달리하는 경우에도 성립한다,
\item
\label{properties-item-c-gg}
$X$가 준분리이고 유한형 준연접층 $\mathcal{F}$가 $X$ 위에 있을 때
마다 어떤 정수 $n_0$가 존재하여
$\mathcal{F} \otimes_{\mathcal{O}_X} \mathcal{L}^{\otimes n}$이
모든 $n \geq n_0$에 대해 전역 생성된다,
\item
\label{properties-item-c-q}
$X$가 준분리이고 유한형 준연접층 $\mathcal{F}$가 $X$ 위에 있을 때
마다 어떤 정수 $n > 0$, $k \geq 0$가 존재하여
$\mathcal{F}$가 $k$개 복사본의 직합의 몫인
$\mathcal{L}^{\otimes - n}$의,
\item
\label{properties-item-c-i-q}
(\ref{properties-item-c-q})와 같은 성질이 $\mathcal{F}$가 $X$ 위의
유한형 아이디얼층 전체를 달리하는 경우에도 성립한다.
\end{enumerate}
\end{proposition}
\begin{proof}
보조정리 \ref{properties-lemma-ample-immersion-into-proj}는
(\ref{properties-item-ample}) $\Rightarrow$ (\ref{properties-item-immersion})이다.
보조정리 \ref{properties-lemma-ample-power-ample}와
\ref{properties-lemma-open-in-proj-ample}는
(\ref{properties-item-ample}) $\Leftarrow$ (\ref{properties-item-immersion})을 준다.
(\ref{properties-item-immersion}) $\Rightarrow$
(\ref{properties-item-s-affine-basis}) $\Rightarrow$
(\ref{properties-item-s-basis})라는 함의는 구성의 절
\ref{constructions-section-proj}에서 명백하다.
보조정리 \ref{properties-lemma-affine-s-opens}는
(\ref{properties-item-s-basis}) $\Rightarrow$ (\ref{properties-item-ample})이다.
따라서 처음 네 조건은 모두 동치이다.

\medskip\noindent
동치인 조건 (1)--(4)를 가정하자.
특히 $X$는 분리된 스킴 $\text{Proj}(S)$의 열린 부분스킴이므로 분리된다.
 준연접층 $\mathcal{F}$를 $X$ 위에서 잡자.
 동차 원소 $s \in S_{+}$를 택하여 $X_s$가 아핀이 되게 하자.
 임의의 절단 $m \in \Gamma(X_s, \mathcal{F})$가
위 (\ref{properties-item-qc-gg})에 표시한 사상 중 하나의 상에 속한다고 주장하자.
이는 이러한 아핀 열린집합 $X_s$들이 $X$를 덮으므로
(\ref{properties-item-qc-gg})를 뜻한다.
실제로 보조정리 \ref{properties-lemma-invert-s-sections}에 의해
$m$은 어떤 $m' \otimes s^{-n}$의 상으로 쓸 수 있는데,
어떤 $n \geq 1$과 어떤
$m' \in \Gamma(X, \mathcal{F} \otimes \mathcal{L}^{\otimes n})$에 대해 그렇다.
이로써 주장이 증명된다.

\medskip\noindent
분명히 (\ref{properties-item-qc-gg}) $\Rightarrow$ (\ref{properties-item-qc-i-gg})이다.
(\ref{properties-item-qc-i-gg})를 가정하고 $\mathcal{L}$이
충분함을 보이자. $x \in X$를 잡자. $U \subset X$인 아핀 열린집합을
택하자. 이는 $x$를 포함한다. $Z = X \setminus U$로 두자.
$Z$를 감소된 닫힌 부분스킴으로 생각할 수 있는데, 스킴
절 \ref{schemes-section-reduced}를 참조하라.
닫힌 부분스킴에 대응하는 준연접 아이디얼층
$\mathcal{I} \subset \mathcal{O}_X$를 두자. 이 층은 $Z$에 대응한다.
가정 (\ref{properties-item-qc-i-gg})에 의해 어떤 $n \geq 1$과 절단
$s \in \Gamma(X, \mathcal{I} \otimes \mathcal{L}^{\otimes n})$가 존재하여
$s$가 $x$에서 소멸하지 않는다(더 정확히는
$s \not \in \mathfrak m_x \mathcal{I}_x \otimes \mathcal{L}_x^{\otimes n}$).
$s$를 $\mathcal{L}^{\otimes n}$의 절단으로 생각할 수 있다.
$Z$를 따라 명백히 소멸하므로 $X_s \subset U$이다.
따라서 $X_s$는 아핀이다. 보조정리 \ref{properties-lemma-affine-cap-s-open}을 보라.
이로써 $\mathcal{L}$이 충분함이 증명된다.
이 시점에서 (1)--(6)이 서로 동치임을 증명했다.

\medskip\noindent
동치인 조건 (1)--(6)을 가정하자. 다음에서는
전역 생성된 두 가군층의 텐서곱도 전역 생성된다는 사실을
더 이상 언급하지 않고 사용한다(구성, 보조정리
\ref{modules-lemma-tensor-product-globally-generated} 참조).
(1)에 의해 원소 $s_i \in S_{d_i}$를 찾을 수 있으며 $d_i \geq 1$이고,
$X = \bigcup_{i = 1, \ldots, n} X_{s_i}$가 되게 할 수 있다.
$d = d_1\ldots d_n$으로 두자. 그러면 $\mathcal{L}^{\otimes d}$는
$$
s_1^{d/d_1}, \ldots, s_n^{d/d_n}.
$$
에 의해 전역 생성된다.
이는 $\mathcal{L}^{\otimes j}$가 전역 생성되면
$\mathcal{L}^{\otimes j + dn}$도 모든 $n \geq 0$에 대해
전역 생성된다는 뜻이다.
$j \in \{0, \ldots, d - 1\}$를 고정하자. 모든 점 $x \in X$에 대해
어떤 $n \geq 1$과 $s$라는 $\mathcal{L}^{j + dn}$의 전역 절단이 존재하며,
이 절단은 $x$에서 소멸하지 않는다. 이는 (\ref{properties-item-qc-gg})를
$\mathcal{F} = \mathcal{L}^{\otimes j}$와 충분한 가역층
$\mathcal{L}^{\otimes d}$에 적용하면 따른다. $X$가 준콤팩트이므로
어느 점에서도 소멸하지 않는 정수 $n_i$들과
$s_i$라는 $\mathcal{L}^{\otimes j + dn_i}$의 전역 절단들의
유한 목록을 택할 수 있는데, 이는 $X$의 어느 점에서도 성립한다.
$\mathcal{L}^{\otimes d}$가 전역 생성되므로
$\mathcal{L}^{\otimes j + dn}$도 전역 생성되는데, 여기서
$n = \max\{n_i\}$이다. 이를 모든 법 $d$의 합동류에 대해 증명했으므로
어떤 $n_0 = n_0(\mathcal{L})$가 존재하여
$\mathcal{L}^{\otimes n}$이 모든 $n \geq n_0$에 대해 전역 생성된다.
이제 $\mathcal{F}$가 전역 생성되면
 $\mathcal{F} \otimes \mathcal{L}^{\otimes n}$도 모든 $n \geq n_0$에 대해
전역 생성됨을 알 수 있다.

\medskip\noindent
동치인 조건 (1)--(6)을 계속 가정하자.
 유한형 준연접층 $\mathcal{F}$를 잡자. 이 층은
 $\mathcal{O}_X$-가군의 준연접층이다.
(\ref{properties-item-qc-gg})의 표준 사상의 상을
$\mathcal{F}_n \subset \mathcal{F}$라 하자.
구성에 의해 $\mathcal{F}_n \otimes \mathcal{L}^{\otimes n}$은
전역 생성된다. (\ref{properties-item-qc-gg})에 의해
$\mathcal{F}$는 $\mathcal{F}_n$ 부분층들의 합이며, $n \geq 1$이다.
Modules, 보조정리
\ref{modules-lemma-finite-type-quasi-compact-colimit}에 의해
$\mathcal{F} = \sum_{n = 1, \ldots, N} \mathcal{F}_n$임을 안다.
이는 어떤 $N \geq 1$에 대해 성립한다.
따라서 $\mathcal{F} \otimes \mathcal{L}^{\otimes n}$은
$n \geq N + n_0(\mathcal{L})$이면 전역 생성된다.
여기서 $n_0(\mathcal{L})$은 위에서와 같다. 그러므로 (1)--(6)은
(\ref{properties-item-c-gg})를 함의한다.

\medskip\noindent
(\ref{properties-item-c-gg})를 가정하자. 준연접층 $\mathcal{F}$를 잡자.
이 층은 $\mathcal{O}_X$-가군의 유한형 층이다.
(\ref{properties-item-c-gg})에 의해 어떤 정수 $n \geq 1$이 존재하여
다음의 표준 사상이
$$
\Gamma(X, \mathcal{F} \otimes_{\mathcal{O}_X} \mathcal{L}^{\otimes n})
\otimes_{\mathbf{Z}} \mathcal{L}^{\otimes -n}
\longrightarrow
\mathcal{F}
$$
전사적이다. $I$를
$\Gamma(X, \mathcal{F} \otimes_{\mathcal{O}_X} \mathcal{L}^{\otimes n})$의
유한 부분집합들의 집합으로 두고 포함관계로 부분순서화하자.
그러면 $I$는 유향 부분순서집합이다.
$i = \{s_1, \ldots, s_{r(i)}\}$에 대해
$\mathcal{F}_i \subset \mathcal{F}$를 다음 사상의 상으로 두자.
$$
\bigoplus\nolimits_{j = 1, \ldots, r(i)} \mathcal{L}^{\otimes -n}
\longrightarrow
\mathcal{F}
$$
이 사상은 $s_j$에 의한 곱셈을 $j$번째 인자에 적용한다.
위 전사성으로부터
$\mathcal{F} = \colim_{i \in I} \mathcal{F}_i$임을 얻는다.
따라서 Modules, 보조정리
\ref{modules-lemma-finite-type-quasi-compact-colimit}를 적용하면
$\mathcal{F} = \mathcal{F}_i$임을 어떤 $i$에 대해 결론낼 수 있다.
그러므로 (\ref{properties-item-c-q})를 증명했다. 즉,
(\ref{properties-item-c-gg}) $\Rightarrow$ (\ref{properties-item-c-q})이다.

\medskip\noindent
(\ref{properties-item-c-q}) $\Rightarrow$ (\ref{properties-item-c-i-q})라는 함의는 자명하다.

\medskip\noindent
마지막으로 (\ref{properties-item-c-i-q})를 가정하자.
$\mathcal{I} \subset \mathcal{O}_X$인 준연접 아이디얼층을 잡자.
보조정리 \ref{properties-lemma-quasi-coherent-colimit-finite-type}에 의해
($X$가 준분리라는 조건을 사용하는 부분이다)
$\mathcal{I} = \colim_\alpha I_\alpha$임을 안다.
각 $I_\alpha$는 유한형 준연접층이다.
가정에 의해 각 $I_\alpha$는
$\mathcal{L}$의 음의 텐서 거듭제곱들의 몫이므로
$\mathcal{I}$에 대해서도 같은 결론을 얻는다
(물론 거듭제곱의 유한성이나 유계성은 없다).
따라서 (\ref{properties-item-c-i-q})는
(\ref{properties-item-qc-i-gg})를 함의한다.
이로써 명제의 증명이 끝난다.
\end{proof}

\begin{lemma}
\label{properties-lemma-ample-on-locally-closed}
스킴 $X$를 잡자. $\mathcal{L}$을 충분한 가역
$\mathcal{O}_X$-가군이라 하자. $i : X' \to X$를 스킴의 사상이라 하자.
다음 조건 중 적어도 하나가 성립한다고 가정하자.
\begin{enumerate}
\item $i$가 준콤팩트 몰입이다,
\item $X'$가 준콤팩트이고 $i$가 몰입이다,
\item $i$가 준콤팩트이고 $X'$와 $i(X')$ 사이의 위상동형을 유도한다,
\item $X'$가 준콤팩트이고 $i$가 $X'$와 $i(X')$ 사이의 위상동형을 유도한다.
\end{enumerate}
그러면 $i^*\mathcal{L}$은 $X'$ 위에서 충분하다.
\end{lemma}

\begin{proof}
조건 (1)과 (3)에서는 $X'$가 준콤팩트임에 주목하자.
이는 정의 \ref{properties-definition-ample}에 의해 $X$가 준콤팩트이기 때문이다.
따라서 (2)와 (4)만 증명하면 충분하다. (2)는 (4)의 특수한 경우이므로
(4)만 증명하면 충분하다.

\medskip\noindent
조건 (4)가 성립한다고 하자. $s \in \Gamma(X, \mathcal{L}^{\otimes d})$에 대해
$s' = i^*s$를 $s$의 $X'$로의 당김으로 표기하자. 그러면
$s'$는 $(i^*\mathcal{L})^{\otimes d}$의 절단이다. 명제
\ref{properties-proposition-characterize-ample}에 의해
$X_s$들, 여기서 $s \in \Gamma(X, \mathcal{L}^{\otimes d})$는
$X$의 위상에 대한 기저를 이룬다. Modules, 비고
\ref{modules-remark-pullback-s-open}에 의해
$X'_{s'} = i^{-1}(X_s)$이고, $X' \to i(X')$가 위상동형이므로
$X'_{s'}$들이 $X'$의 위상에 대한 기저를 이룬다.
따라서 명제 \ref{properties-proposition-characterize-ample}에 의해
$i^*\mathcal{L}$은 충분하다.
\end{proof}

\begin{lemma}
\label{properties-lemma-ample-on-product}
$S$를 준분리 스킴이라 하자. $X$, $Y$를 $S$ 위의 스킴이라 하자.
$\mathcal{L}$을 충분한 가역 $\mathcal{O}_X$-가군이라 하고,
$\mathcal{N}$을 충분한 가역 $\mathcal{O}_Y$-가군이라 하자.
그러면
$\mathcal{M} = \text{pr}_1^*\mathcal{L}
\otimes_{\mathcal{O}_{X \times_S Y}} \text{pr}_2^*\mathcal{N}$
은 $X \times_S Y$ 위의 충분한 가역층이다.
\end{lemma}

\begin{proof}
사상 $i : X \times_S Y \to X \times Y$는 준콤팩트 몰입이다.
스킴, 보조정리 \ref{schemes-lemma-fibre-product-after-map}를 보라.
한편 $\mathcal{M}$은 $i$로 당긴 $X \times Y$ 위의 대응하는
가역 가군이다. 보조정리
\ref{properties-lemma-ample-on-locally-closed}에 의해
$X \times Y$에 대해 명제를 증명하면 충분하다.
$\mathcal{M}$에 대해 정의 \ref{properties-definition-ample}의
(1)과 (2)를 $X \times Y$ 위에서 확인하자.

\medskip\noindent
$X$와 $Y$가 준콤팩트이므로 $X \times Y$도 준콤팩트이다.
점 $z \in X \times Y$를 잡자. 사영을 $x \in X$와 $y \in Y$라 하자.
$n > 0$과 $s \in \Gamma(X, \mathcal{L}^{\otimes n})$을 택하여
$X_s$가 $x$의 아핀 열린 근방이 되게 하자.
$m > 0$과 $t \in \Gamma(Y, \mathcal{N}^{\otimes m})$을 택하여
$Y_t$가 $y$의 아핀 열린 근방이 되게 하자.
그러면 $r = \text{pr}_1^*s \otimes \text{pr}_2^*t$는
$\mathcal{M}$의 절단이고
$(X \times Y)_r = X_s \times Y_t$이다.
이는 $z$의 아핀 열린 근방이므로 증명이 끝난다.
\end{proof}

\section{아핀 스킴과 준아핀 스킴}
\label{properties-section-affine-quasi-affine}

\begin{lemma}
\label{properties-lemma-quasi-affine-O-ample}
스킴 $X$를 잡자.
$X$가 준아핀일 필요충분조건은 $\mathcal{O}_X$가 충분한 것이다.
\end{lemma}

\begin{proof}
$X$가 준아핀이라고 하자. $A = \Gamma(X, \mathcal{O}_X)$로 두자.
보조정리 \ref{properties-lemma-quasi-affine}에 따른 열린 몰입
$$
j : X \longrightarrow \Spec(A)
$$
을 생각하자. $\Spec(A) = \text{Proj}(A[T])$임에 주목하자.
구성, 예 \ref{constructions-example-trivial-proj}를 보라.
따라서 보조정리 \ref{properties-lemma-open-in-proj-ample}를 적용하여
$\mathcal{O}_X$가 충분함을 얻을 수 있다.

\medskip\noindent
$\mathcal{O}_X$가 충분하다고 하자.
$\Gamma_*(X, \mathcal{O}_X) \cong A[T]$가 등급환으로서 성립함에 주목하자.
따라서 보조정리 \ref{properties-lemma-ample-immersion-into-proj}와
\ref{properties-lemma-quasi-affine}로부터 결과를 얻는다.
위에서 본 $\Spec(A) = \text{Proj}(A[T])$가 임의의 환
$A$에 대해 성립함을 함께 사용했다.
\end{proof}

\begin{lemma}
\label{properties-lemma-quasi-affine-locally-closed}
$X$를 준아핀 스킴이라 하자. 임의의 준콤팩트 몰입
$i : X' \to X$에 대해 $X'$는 준아핀이다.
\end{lemma}

\begin{proof}
이는 충분한 가역층에 관한 내용을 사용하지 않고 쪽지에 직접 증명할 수 있다.
독자가 직접 해 보기를 권한다.
$X$가 준아핀이므로 보조정리 \ref{properties-lemma-quasi-affine-O-ample}에 의해
$\mathcal{O}_X$는 충분하다.
그러면 보조정리 \ref{properties-lemma-ample-on-locally-closed}에 의해
$\mathcal{O}_{X'}$는 충분하다. 따라서 보조정리
\ref{properties-lemma-quasi-affine-O-ample}에 의해 $X'$는 준아핀이다.
\end{proof}

\begin{lemma}
\label{properties-lemma-characterize-affine}
스킴 $X$를 잡자. 다음 조건을 만족하는 유한 개의 원소
$f_1, \ldots, f_n \in \Gamma(X, \mathcal{O}_X)$가 존재한다고 하자.
\begin{enumerate}
\item 각 $X_{f_i}$가 $X$의 아핀 열린집합이다,
\item $f_1, \ldots, f_n$이 $\Gamma(X, \mathcal{O}_X)$에서 생성하는
아이디얼이 단위원 아이디얼이다.
\end{enumerate}
그러면 $X$는 아핀이다.
\end{lemma}

\begin{proof}
명제의 $f_1, \ldots, f_n$을 택했다고 하자.
$1 = \sum g_i f_i$로 쓸 수 있는데, 여기서 어떤
$g_j \in \Gamma(X, \mathcal{O}_X)$를 택할 수 있으므로
$X = \bigcup X_{f_i}$임은 자명하다. (어떤 점에서도
$f_i$들이 모두 소멸할 수 없다.) 각 $X_{f_i}$는 아핀이므로
준콤팩트이다. 따라서 $X$도 준콤팩트이다.
그러므로 위의 보조정리 \ref{properties-lemma-quasi-affine-O-ample}에 의해
$X$는 준아핀이다.
열린 몰입
$$
j : X \to \Spec(\Gamma(X, \mathcal{O}_X)),
$$
을 보조정리 \ref{properties-lemma-quasi-affine}에 따라 생각하자.
우변의 표준 열린집합 $D(f_i)$의 역상은 좌변의 $X_{f_i}$와 같고,
사상 $j$는
$X_{f_i} \cong D(f_i)$라는 동형을 유도한다. 보조정리
\ref{properties-lemma-invert-f-affine}를 보라.
$f_i$들이 단위 아이디얼을 생성하므로
$\Spec(\Gamma(X, \mathcal{O}_X))
= \bigcup_{i = 1, \ldots, n} D(f_i)$이다. 따라서 $j$는 동형이다.
\end{proof}

\section{준연접층과 충분한 가역층}
\label{properties-section-ample-quasi-coherent}

\noindent
이 절의 주제: 충분한 가역층이 있으면 모든 준연접층은
등급 가군에서 유래한다.

\begin{situation}
\label{properties-situation-ample}
스킴 $X$를 잡자.
충분한 가역층 $\mathcal{L}$를 $X$ 위에 잡자.
$S = \Gamma_*(X, \mathcal{L})$을 등급환으로 두자.
$Y = \text{Proj}(S)$로 두자.
보조정리 \ref{properties-lemma-map-into-proj}의 표준 사상
$f : X \to Y$를 잡자.
이 사상에는 $\mathbf{Z}$-등급 $\mathcal{O}_X$-대수 사상
$\bigoplus f^*\mathcal{O}_Y(n) \to \bigoplus \mathcal{L}^{\otimes n}$
이 함께 주어진다.
\end{situation}

\noindent
다음 보조정리는 사실 다음 보조정리의 특수한 경우이지만,
먼저 그 유효성을 지적하는 것이 좋을 듯하다.

\begin{lemma}
\label{properties-lemma-ample-gcd-is-one}
상황 \ref{properties-situation-ample}에서.
표준 사상 $f : X \to Y$는 $X$를 열린 부분스킴
$W = W_1 \subset Y$로 보낸다.
여기서 $\mathcal{O}_Y(1)$은 가역이고 모든 곱셈 사상
$\mathcal{O}_Y(n) \otimes_{\mathcal{O}_Y} \mathcal{O}_Y(m) \to
\mathcal{O}_Y(n + m)$
은 동형이다(구성, 보조정리
\ref{constructions-lemma-where-invertible}를 보라).
더욱이 사상
$f^*\mathcal{O}_Y(n) \to \mathcal{L}^{\otimes n}$은 모두 동형이다.
\end{lemma}

\begin{proof}
명제 \ref{properties-proposition-characterize-ample}에 의해 어떤 정수
$n_0$가 존재하여 $\mathcal{L}^{\otimes n}$은 모든
$n \geq n_0$에 대해 전역 생성된다. 점 $x \in X$를 잡자.
위 결과에 의해 $a \in S_{n_0}$와 $b \in S_{n_0 + 1}$을 택하여
$a$와 $b$가 $x$에서 소멸하지 않게 할 수 있다. 따라서
$f(x) \in D_{+}(a) \cap D_{+}(b) = D_{+}(ab)$이다.
구성, 보조정리 \ref{constructions-lemma-where-invertible}에 의해
$f(x) \in W_1$임을 얻는다. 이제
$f$의 사상을 구성할 때 사용한 구성, 보조정리
\ref{constructions-lemma-invertible-map-into-proj}에 의해
사상
$f^*\mathcal{O}_Y(n_0) \to \mathcal{L}^{\otimes n_0}$와
$f^*\mathcal{O}_Y(n_0 + 1) \to \mathcal{L}^{\otimes n_0 + 1}$은
$x$의 근방에서 동형이다. 대수 구조와 $f$가 $W$로 사상된다는
사실과의 호환성으로부터 모든 사상
$f^*\mathcal{O}_Y(n) \to \mathcal{L}^{\otimes n}$이
$x$의 근방에서 동형임을 결론낸다. 따라서 증명되었다.
\end{proof}

\noindent
Modules, 정의 \ref{modules-definition-gamma-star}에서 상기하자.
국소 환 달린 공간 $X$, 가역층 $\mathcal{L}$,
$\mathcal{O}_X$-가군 $\mathcal{F}$가 주어지면
등급 $\Gamma_*(X, \mathcal{L})$-가군
$$
\Gamma_*(X, \mathcal{L}, \mathcal{F}) =
\bigoplus\nolimits_{n \in \mathbf{Z}}
\Gamma(X, \mathcal{F} \otimes_{\mathcal{O}_X} \mathcal{L}^{\otimes n}).
$$
을 얻는다. 다음 보조정리는 상황 \ref{properties-situation-ample}에서
이 등급 가군으로부터 준연접 $\mathcal{O}_X$-가군 $\mathcal{F}$를
복원할 수 있다고 말한다. 구성, 보조정리
\ref{constructions-lemma-quasi-coherent-projective-space}도 보라.
그곳에서는 특수한 경우 $X = \mathbf{P}^n_R$에 대해 이 보조정리를 증명한다.

\begin{lemma}
\label{properties-lemma-ample-quasi-coherent}
상황 \ref{properties-situation-ample}에서.
$\mathcal{F}$를 $X$ 위의 준연접층이라 하자.
$M = \Gamma_*(X, \mathcal{L}, \mathcal{F})$를 등급 $S$-가군으로 두자.
다음 동형들이 존재한다.
$$
f^*\widetilde{M} \longrightarrow \mathcal{F}
$$
이는 $\mathcal{F}$에 함자적이며 다음 사상을 만족한다.
$M_0 \to \Gamma(\text{Proj}(S), \widetilde{M}) \to \Gamma(X, \mathcal{F})$
는 항등 사상이다.
\end{lemma}

\begin{proof}
$s \in S_{+}$를 동차 원소로 택하여 $X_s$가 $X$의 아핀 열린집합이 되게 하자.
$\widetilde{M}|_{D_{+}(s)}$는
$S_{(s)}$-가군 $M_{(s)}$에 대응함을 상기하자.
구성, 보조정리 \ref{constructions-lemma-proj-sheaves}를 보라.
또한 $f^{-1}(D_{+}(s)) = X_s$임을 상기하자.
$X$가 충분한 가역층을 가지므로 준콤팩트이고
준분리임을 절 \ref{properties-section-ample}에서 보았다.
보조정리 \ref{properties-lemma-invert-s-sections}에 의해 표준 동형
$M_{(s)} = \Gamma_*(X, \mathcal{L}, \mathcal{F})_{(s)} \to
\Gamma(X_s, \mathcal{F})$이 존재한다.
$\mathcal{F}$가 준연접이므로 이는 표준 동형
$$
f^*\widetilde{M}|_{X_s} \to \mathcal{F}|_{X_s}
$$
으로 이어진다.
$\mathcal{L}$이 $X$ 위에서 충분하므로 $X$는
$X_s$ 꼴의 아핀 열린집합들로 덮인다.
따라서 표시된 사상들이 겹침에서 이어 붙는다는 것을 증명하면 충분하다.
이 증명은 생략한다.
\end{proof}

\begin{remark}
\label{properties-remark-neurotic}
보조정리 \ref{properties-lemma-ample-quasi-coherent}의 가정과 표기를 사용하자.
그 보조정리에서 표시한 사상을 $\theta_\mathcal{F}$로 쓰자.
보조정리 \ref{properties-lemma-ample-gcd-is-one}의 동형
$f^*\mathcal{O}_Y(n) \to \mathcal{L}^{\otimes n}$은 바로
$\theta_{\mathcal{L}^{\otimes n}}$임에 주목하자.
곱셈 사상
$$
\widetilde{M} \otimes_{\mathcal{O}_Y} \mathcal{O}_Y(n)
\longrightarrow
\widetilde{M(n)}
$$
을 생각하자. 구성, 방정식
(\ref{constructions-equation-multiply-more-generally})를 보라.
이를 $X$로 당겨 다음을 생각하자.
$$
\xymatrix{
f^*\widetilde{M} \otimes_{\mathcal{O}_X} f^*\mathcal{O}_Y(n)
\ar[r]
\ar[d]_{\theta_\mathcal{F} \otimes \theta_{\mathcal{L}^{\otimes n}}}
&
f^*\widetilde{M(n)}
\ar[d]^{\theta_{\mathcal{F} \otimes \mathcal{L}^{\otimes n}}}
\\
\mathcal{F} \otimes \mathcal{L}^{\otimes n} \ar[r]^{\text{id}} &
\mathcal{F} \otimes \mathcal{L}^{\otimes n}
}
$$
여기서 명백한 식별
$M(n) = \Gamma_*(X, \mathcal{L}, \mathcal{F} \otimes \mathcal{L}^{\otimes n})$
을 사용했다. 이 도표는 가환한다. 증명은 생략한다.
\end{remark}

\noindent
다음 보조정리는 보조정리
\ref{properties-lemma-ample-quasi-coherent}에서 (또는 그 역으로) 얻을 수 있을 것 같지만,
그 증명을 다시 쓰는 편이 더 간단하다.

\begin{lemma}
\label{properties-lemma-proj-quasi-coherent}
$S$를 등급환이라 하고 $X = \text{Proj}(S)$가 준콤팩트라고 하자.
$\mathcal{F}$를 준연접 $\mathcal{O}_X$-가군이라 하자.
$M = \bigoplus_{n \in \mathbf{Z}} \Gamma(X, \mathcal{F}(n))$을
등급 $S$-가군으로 두자. 구성, 절
\ref{constructions-section-invertible-on-proj}를 보라.
사상
$$
\widetilde{M} \longrightarrow \mathcal{F}
$$
은 구성, 보조정리
\ref{constructions-lemma-comparison-proj-quasi-coherent}의 사상이며 동형이다.
$X$가 표준 열린집합 $D_+(f)$들로 덮이고 각 $f$의 차수가 $1$이면,
유도된 사상
$M_n \to \Gamma(X, \mathcal{F}(n))$은 항등 사상이다.
\end{lemma}

\begin{proof}
$X$가 준콤팩트이므로 양의 차수를 갖는 동차 원소
$f_1, \ldots, f_n \in S$를 찾아
$X = D_+(f_1) \cup \ldots \cup D_+(f_n)$으로 쓸 수 있다.
$d$를 $f_1, \ldots, f_n$의 차수들의 최소공배수라 하자.
$f_i$를 거듭제곱으로 바꾸면 각 $f_i$의 차수가 $d$라고 가정할 수 있다.
그러면 $\mathcal{L} = \mathcal{O}_X(d)$는 가역이고,
곱셈 사상
$\mathcal{O}_X(ad) \otimes \mathcal{O}_X(bd) \to \mathcal{O}_X((a + b)d)$
은 동형이며, 각 $f_i$는 전역 절단 $s_i$를 정하여
$\mathcal{L}$의 절단이 되고 $X_{s_i} = D_+(f_i)$이다. 구성, 보조정리
\ref{constructions-lemma-where-invertible}와
\ref{constructions-lemma-principal-open}을 보라.
따라서
$\Gamma(X, \mathcal{F}(ad)) =
\Gamma(X, \mathcal{F} \otimes \mathcal{L}^{\otimes a})$이다.
$\widetilde{M}|_{D_{+}(f_i)}$는
$S_{(f_i)}$-가군 $M_{(f_i)}$에 대응함을 상기하자. 구성,
보조정리 \ref{constructions-lemma-proj-sheaves}를 보라.
$f_i$의 차수가 $d$이므로 $M_{(f_i)}$의 동형류는
$M$의 동차 성분들 중 $d$로 나누어지는 것에만 의존한다.
더 정확히는 $M_{(f_i)}$의 동형류는 등급
$\Gamma_*(X, \mathcal{L})$-가군
$\Gamma_*(X, \mathcal{L}, \mathcal{F})$와
$s_i$ 및 $f_i$의 상이 놓이는 $\Gamma_*(X, \mathcal{L})$에만 의존한다.
가정에 의해 $X$는 준콤팩트이고, 구성, 보조정리
\ref{constructions-lemma-proj-separated}에 의해 분리되어 있다.
보조정리 \ref{properties-lemma-invert-s-sections}에 의해 표준 동형
$$
M_{(f_i)} = \Gamma_*(X, \mathcal{L}, \mathcal{F})_{(s_i)} \to
\Gamma(X_{s_i}, \mathcal{F}).
$$
이 존재한다. 구성, 보조정리
\ref{constructions-lemma-comparison-proj-quasi-coherent}에서
사상을 구성하는 방법을 보면 이 사상이 $D_+(f_i)$ 위에서 동형임을 알 수 있다.
따라서 $X$가 이 열린집합들로 덮이므로 전체에서도 동형이다.
마지막 명제의 증명은 생략한다.
\end{proof}

\section{적절한 아핀 열린집합 찾기}
\label{properties-section-finding-affine-opens}

\noindent
이 절에서는 더 일반적이거나 덜 일반적인 상황에서 아핀 열린집합의
존재에 관한 결과를 모은다.

\begin{lemma}
\label{properties-lemma-maximal-points-affine}
$X$를 준분리 스킴이라 하자.
$Z_1, \ldots, Z_n$을 서로 다른 $X$의 기약 성분들이라 하자.
위상, 절 \ref{topology-section-irreducible-components}를 보라.
$\eta_i \in Z_i$를 그 일반점이라 하자. 스킴이 준콤팩트임에 관한
스킴, 보조정리 \ref{schemes-lemma-scheme-sober}를 보라.
$\eta_i \in U_i$인 아핀 열린 근방들을 택하여
$U_i \cap U_j = \emptyset$가 모든 $i \not = j$에 대해 성립하게 할 수 있다.
특히 $U = U_1 \cup \ldots \cup U_n$은
$\eta_1, \ldots, \eta_n$을 모두 포함하는 아핀 열린집합이다.
\end{lemma}

\begin{proof}
임의의 아핀 열린집합 $V_i$를 택하자. 이는 $\eta_i$를 포함하고
닫힌집합 $Z_1 \cup \ldots \hat Z_i \ldots \cup Z_n$과 서로소이다.
$X$가 준분리이므로 각 $i$에 대해
$W_i = \bigcup_{j, j \not = i} V_i \cap V_j$는
$V_i$의 준콤팩트 열린집합이며 $\eta_i$를 포함하지 않는다.
대수, 보조정리 \ref{algebra-lemma-standard-open-containing-maximal-point}에 의해
$U_i \subset V_i$인 열린 근방을 찾아 $\eta_i$를 포함하고
$W_i$와 서로소가 되게 할 수 있다.
마지막으로 $U$는 아핀이다. 왜냐하면 이는
$R_1 \times \ldots \times R_n$의 스펙트럼이고,
여기서 $R_i = \mathcal{O}_X(U_i)$이기 때문이다.
스킴, 보조정리 \ref{schemes-lemma-disjoint-union-affines}를 보라.
\end{proof}

\begin{remark}
\label{properties-remark-maximal-points-affine}
위 보조정리 \ref{properties-lemma-maximal-points-affine}는 $X$가
준분리되지 않으면 거짓이다. 예를 들어 보자.
$R = \mathbf{Q}[x, y_1, y_2, \ldots]/((x-i)y_i)$로 두자.
$\mathfrak p = (y_1, y_2, \ldots)$라는 극소 소 아이디얼을
$R$에서 생각하자.
$\Spec(R)$의 두 복사본을
(준콤팩트가 아닌) 열린집합
$\Spec(R) \setminus V(\mathfrak p)$를 따라 붙여
스킴 $X$를 얻자(붙임은 스킴, 예
\ref{schemes-example-affine-space-zero-doubled}와 같다).
그러면 $X$의 $\mathfrak p$에 대응하는 두 극대점은
공통의 아핀 열린집합에 포함되지 않는다.
그 이유는 $\Spec(R)$에서 $\mathfrak p$를 포함하는 임의의 열린집합이
$x = i$, $y_j = 0$ ($j \not = i$)인 “직선”들을
매개변수 $y_i$와 함께 무한히 많이 포함하기 때문이다.
세부 사항은 생략한다.
\end{remark}

\noindent
위 예에도 불구하고 유한한 기약 닫힌 부분집합의 “대부분”에 대해서는,
적어도 $X$가 준콤팩트라면 위 보조정리
\ref{properties-lemma-maximal-points-affine}를 적용할 수 있다.
이는 $X$가 분리된 조밀 열린집합을 포함하기 때문에 성립한다.

\begin{lemma}
\label{properties-lemma-quasi-compact-dense-open-separated}
$X$를 준콤팩트 스킴이라 하자.
분리된 조밀 열린집합 $V \subset X$가 존재한다.
\end{lemma}

\begin{proof}
$X = \bigcup_{i = 1, \ldots, n} U_i$가 $n$개의 아핀 열린 부분스킴의
합이라고 하자. $n$에 대한 귀납법으로 보조정리를 증명한다.
$n = 1$이면 자명하다.
귀납 가정에 의해 분리된 조밀 열린 부분스킴
$V' \subset \bigcup_{i = 1, \ldots, n - 1} U_i$가 존재한다고 하자.
다음을 생각하자.
$$
V = V' \amalg (U_n \setminus \overline{V'}).
$$
$V$가 분리된 조밀 열린 부분스킴인 $X$의 조밀 열린 부분스킴임은 자명하다.
\end{proof}

\noindent
$X$가 준분리이면서 준콤팩트이더라도 분리된 준콤팩트 조밀 열린집합이
존재하지 않을 수 있음에 주의하자. 예는 예, 보조정리
\ref{examples-lemma-no-dense-separated-quasi-compact-open-in-qcqs}를 보라.
다음은 위 보조정리 \ref{properties-lemma-maximal-points-affine}의 약간의 보강이다.

\begin{lemma}
\label{properties-lemma-point-and-maximal-points-affine}
$X$를 준분리 스킴이라 하자. $Z_1, \ldots, Z_n$을
$X$의 서로 다른 기약 성분들이라 하자. 그 일반점은
$\eta_i \in Z_i$라 하자.
$x \in X$를 임의의 점이라 하자.
아핀 열린집합 $U \subset X$가 존재하여
$x$와 모든 $\eta_i$를 포함한다.
\end{lemma}

\begin{proof}
$x \in Z_1 \cap \ldots \cap Z_r$이고
$x \not \in Z_{r + 1}, \ldots, Z_n$이라고 하자.
그러면 $W \subset X$인 아핀 열린집합을 택하여
$x \in W$이고 $W \cap Z_i = \emptyset$가
$i = r + 1, \ldots, n$에 대해 성립하게 할 수 있다.
분명히 $\eta_i \in W$는 $i = 1, \ldots, r$에 대해 성립한다.
보조정리 \ref{properties-lemma-maximal-points-affine}에 의해
서로소인 아핀 열린집합 $U_i \subset X$들을 택하여
$\eta_i \in U_i$가 $i = r + 1, \ldots, n$에 대해 성립하게 할 수 있다.
$X$가 준분리이므로 열린집합 $W \cap U_i$는 준콤팩트이고
$\eta_i$를 포함하지 않으며, 이는 $i = r + 1, \ldots, n$에 대해 성립한다.
따라서 대수, 보조정리
\ref{algebra-lemma-standard-open-containing-maximal-point}에 의해
$U_i$를 줄여 $W \cap U_i = \emptyset$가
$i = r + 1, \ldots, n$에 대해 성립하게 할 수 있다.
그러면
$U = W \cup \bigcup_{i = r + 1, \ldots, n} U_i$는 서로소이고,
따라서 (스킴, 보조정리
\ref{schemes-lemma-disjoint-union-affines}에 의해) 적절한 아핀 열린집합이다.
\end{proof}

\begin{lemma}
\label{properties-lemma-ample-finite-set-in-affine}
스킴 $X$를 잡자. 다음 조건 중 하나가 성립한다고 하자.
\begin{enumerate}
\item 스킴 $X$는 준아핀이다.
\item 스킴 $X$는 아핀 스킴의 국소 닫힌 부분스킴과 동형이다.
\item $X$ 위에 충분한 가역층이 존재한다.
\item 스킴 $X$는 $\text{Proj}(S)$의 국소 닫힌 부분스킴과
동형이며, 여기서 $S$는 등급환이다.
\end{enumerate}
임의의 유한 부분집합 $E \subset X$에 대해
아핀 열린집합 $U \subset X$가 존재하여 $E \subset U$이다.
\end{lemma}

\begin{proof}
Properties, 정의 \ref{properties-definition-quasi-affine}에 의해
준아핀 스킴은 아핀 스킴의 준콤팩트 열린 부분스킴이다.
임의의 아핀 스킴 $\Spec(R)$는 $\text{Proj}(R[X])$와 동형이며,
$R[X]$에 $\deg(X) = 1$로 등급을 준다.
명제 \ref{properties-proposition-characterize-ample}에 의해
$X$가 충분한 가역층을 가지면 $X$는
$\text{Proj}(S)$의 열린 부분스킴과 동형이며, 여기서 $S$는 등급환이다.
따라서 경우 (4)에 대해 보조정리를 증명하면 충분하다.
(독자가 경우 (2)를 직접 증명해 보기를 권한다.)

\medskip\noindent
따라서 $X \subset \text{Proj}(S)$가 국소 닫힌 부분스킴이라고 하자.
여기서 $S$는 등급환이다. $T = \overline{X} \setminus X$로 두자.
표준 열린집합 $D_{+}(f)$들이
$\text{Proj}(S)$의 위상에 대한 기저를 이룸을 상기하자.
$E$가 유한이므로 다음을 만족하는 유한 개의 동차 원소
$f_i \in S_{+}$를 택할 수 있다.
$$
E \subset
D_{+}(f_1) \cup \ldots \cup D_{+}(f_n) \subset
\text{Proj}(S) \setminus T
$$
$E = \{\mathfrak p_1, \ldots, \mathfrak p_m\}$를
$\text{Proj}(S)$의 부분집합으로 보자.
아이디얼
$I = (f_1, \ldots, f_n) \subset S$를 생각하자.
 $I \not \subset \mathfrak p_j$가 모든 $j = 1, \ldots, m$에 대해
성립하므로 대수, 보조정리
\ref{algebra-lemma-graded-silly}에 의해
동차 원소 $f \in I$를 택하여
$f \not \in \mathfrak p_j$가 모든 $j = 1, \ldots, m$에 대해
성립하게 할 수 있다.
그러면
$E \subset D_{+}(f) \subset
D_{+}(f_1) \cup \ldots \cup D_{+}(f_n)$이다.
$D_{+}(f)$가 $T$와 만나지 않으므로
$X \cap D_{+}(f)$는 아핀 스킴 $D_{+}(f)$의 닫힌 부분스킴이다.
따라서 이는 원하는 $X$의 아핀 열린집합이다.
\end{proof}

\begin{lemma}
\label{properties-lemma-ample-finite-set-in-principal-affine}
스킴 $X$를 잡자. $\mathcal{L}$을 $X$ 위의 충분한 가역층이라 하자.
다음을 생각하자.
$$
E \subset W \subset X
$$
여기서 $E$는 유한하고 $W$는 $X$에서 열린집합이다.
그러면 어떤 $n > 0$과 절단
$s \in \Gamma(X, \mathcal{L}^{\otimes n})$이 존재하여
$X_s$가 아핀이고 $E \subset X_s \subset W$이다.
\end{lemma}

\begin{proof}
독자는 보조정리 \ref{properties-lemma-ample-finite-set-in-affine}의 증명을
수정하여 이 보조정리를 증명할 수 있다. 여기서는 그 보조정리로부터
이 결과를 이끌어 내겠다.
보조정리 \ref{properties-lemma-ample-finite-set-in-affine}에 의해
$U \subset W$인 아핀 열린집합을 택하여 $E \subset U$가 되게 할 수 있다.
등급환
$S = \Gamma_*(X, \mathcal{L}) =
\bigoplus_{n \geq 0} \Gamma(X, \mathcal{L}^{\otimes n})$
을 생각하자.
각 $x \in E$에 대해 절단들의 등급 아이디얼
$\mathfrak p_x \subset S$를 두자. 이는 $x$에서 소멸하는 절단들로
이루어지므로 $\mathfrak p_x$가 소 아이디얼임은 자명하다.
또한 $\mathcal{L}$의 어떤 거듭제곱이 전역 생성되므로
$S_{+} \not \subset \mathfrak p_x$이다.
모든 점에서 소멸하는 절단들의 등급 아이디얼
$I \subset S$를 $X \setminus U$에 대해 두자.
$X_s$들이 위상의 기저를 이루므로
 $I \not \subset \mathfrak p_x$가 모든 $x \in E$에 대해 성립한다.
(등급) 소 아이디얼 피하기, 즉 대수, 보조정리
\ref{algebra-lemma-graded-silly}에 의해 동차 원소
$s \in I$를 택하여 $s \not \in \mathfrak p_x$가
모든 $x \in E$에 대해 성립하게 할 수 있다.
그러면 $E \subset X_s \subset U$이고
보조정리 \ref{properties-lemma-affine-cap-s-open}에 의해 $X_s$는 아핀이다.
\end{proof}

\begin{lemma}
\label{properties-lemma-quasi-affine-invertible-nonvanishing-section}
$X$를 준아핀 스킴이라 하자. $\mathcal{L}$을 가역
$\mathcal{O}_X$-가군이라 하자. $E \subset W \subset X$라 하고
$E$는 유한하며 $W$는 열린집합이라고 하자. 그러면
$s \in \Gamma(X, \mathcal{L})$인 절단이 존재하여
$X_s$는 아핀이고 $E \subset X_s \subset W$이다.
\end{lemma}

\begin{proof}
이 보조정리의 증명은 대수, 보조정리
\ref{algebra-lemma-silly}의 증명과 많은 공통점이 있다.
$E = \{x_1, \ldots, x_n\}$이라 하자.
$E = W = \emptyset$이면 $s = 0$으로 충분하다.
$W \not = \emptyset$이면 필요하다면 점을 하나 추가하여
$E \not = \emptyset$라고 가정할 수 있다.
따라서 $n \geq 1$이라 가정할 수 있다.
$n$에 대한 귀납법으로 보조정리를 증명한다.

\medskip\noindent
기초 단계: $n = 1$. $W$를 $x_1$의 $W$ 안의
아핀 열린 근방으로 바꾼 뒤 $W$가 아핀이라고 가정할 수 있다.
보조정리 \ref{properties-lemma-quasi-affine-O-ample}와
명제 \ref{properties-proposition-characterize-ample}를 결합하면
모든 준연접 $\mathcal{O}_X$-가군이 전역 생성됨을 알 수 있다.
따라서 $s$라는 $\mathcal{L}$의 전역 절단을 택하여
$x_1$에서 소멸하지 않게 할 수 있다.
한편 $Z \subset X$를 $X \setminus W$ 위에 유도된
감소된 닫힌 부분스킴이라 하자.
준연접 아이디얼층 $\mathcal{I}$인 $Z$에 전역 생성을 적용하여
전역 절단 $f$를 $\mathcal{I}$에서 택하면
$x_1$에서 소멸하지 않게 할 수 있다.
그러면 $s' = fs$는 $\mathcal{L}$의 전역 절단이고
$x_1$에서 소멸하지 않으며
$X_{s'} \subset W$이다. 따라서 $X_{s'}$는 아핀이다.
보조정리 \ref{properties-lemma-affine-cap-s-open}을 보라.

\medskip\noindent
$n > 1$일 때의 귀납 단계. $x_i \leadsto x_j$ 특수화가
$i \not = j$인 경우 존재한다면,
$\{x_1, \ldots, x_n\} \setminus \{x_i\}$에 대해
보조정리를 증명하는 것으로 충분하고 귀납에 의해 끝난다.
따라서 $x_i$들 사이에 특수화가 없다고 가정할 수 있다.
보조정리 \ref{properties-lemma-ample-finite-set-in-affine} 또는
보조정리 \ref{properties-lemma-ample-finite-set-in-principal-affine}에 의해
$W$가 아핀이라고 가정할 수 있다.
귀납 가정에 의해 전역 절단 $s$를 $\mathcal{L}$에서 찾아
$X_s \subset W$가 아핀이고 $x_1, \ldots, x_{n - 1}$을 포함하게 할 수 있다.
$x_n \in X_s$이면 끝난다.
$s$가 $x_n$에서 영이라고 가정하자. $n = 1$인 경우에 의해
전역 절단 $s'$를 $\mathcal{L}$에서 찾아
$\{x_n\} \subset X_{s'} \subset
W \setminus \overline{\{x_1, \ldots, x_{n - 1}\}}$
가 되게 할 수 있다.
여기서 $x_n$이 $x_1, \ldots, x_{n - 1}$의 특수화가 아니라는
사실을 사용했다.
그러면 $s + s'$는 $\mathcal{L}$의 전역 절단이고
$x_1, \ldots, x_n$에서 소멸하지 않으며
$X_{s + s'} \subset W$이다. 앞에서와 같이 결론을 얻는다.
\end{proof}

\begin{lemma}
\label{properties-lemma-ring-affine-open-injective-local-ring}
$X$를 스킴이라 하고 $x \in X$를 점이라 하자. $U \subset X$인
$x$의 아핀 열린 근방이 존재하여 표준 사상
$\mathcal{O}_X(U) \to \mathcal{O}_{X, x}$가 다음 각 경우에 단사이다.
\begin{enumerate}
\item $X$가 정수적이다.
\item $X$가 국소 뇌터이다.
\item $X$가 축약이고 기약 성분의 개수가 유한하다.
\end{enumerate}
\end{lemma}

\begin{proof}
대수로 옮기면 이는 대수의 보조정리
\ref{algebra-lemma-subring-of-local-ring}에서 따른다.
\end{proof}

\begin{lemma}
\label{properties-lemma-standard-open-two-affines}
$U$, $V$를 아핀 스킴이라 하고 $W \to U$와 $W \to V$를
열린 몰입이라 하자. 임의의 $w \in W$에 대하여 $W' \subset W$인
$w$의 아핀 열린 근방이 존재하여 $W'$가 $U$와 $V$ 모두의
표준 열린집합으로 사상된다.
\end{lemma}

\begin{proof}
$X$를 $U$와 $V$를 $W$를 따라 붙여 얻은 스킴이라 하고
Schemes의 보조정리 \ref{schemes-lemma-standard-open-two-affines}를
적용한다.
\end{proof}
